%Version 3.1 December 2024
% See section 11 of the User Manual for version history
%
%%%%%%%%%%%%%%%%%%%%%%%%%%%%%%%%%%%%%%%%%%%%%%%%%%%%%%%%%%%%%%%%%%%%%%
%%                                                                 %%
%% Please do not use \input{...} to include other tex files.       %%
%% Submit your LaTeX manuscript as one .tex document.              %%
%%                                                                 %%
%% All additional figures and files should be attached             %%
%% separately and not embedded in the \TeX\ document itself.       %%
%%                                                                 %%
%%%%%%%%%%%%%%%%%%%%%%%%%%%%%%%%%%%%%%%%%%%%%%%%%%%%%%%%%%%%%%%%%%%%%

%%\documentclass[referee,sn-basic]{sn-jnl}% referee option is meant for double line spacing

%%=======================================================%%
%% to print line numbers in the margin use lineno option %%
%%=======================================================%% 

%%\documentclass[lineno,pdflatex,sn-basic]{sn-jnl}% Basic Springer Nature Reference Style/Chemistry Reference Style

%%=========================================================================================%%
%% the documentclass is set to pdflatex as default. You can delete it if not appropriate.  %%
%%=========================================================================================%%

%%\documentclass[sn-basic]{sn-jnl}% Basic Springer Nature Reference Style/Chemistry Reference Style

%%Note: the following reference styles support Namedate and Numbered referencing. By default the style follows the most common style. To switch between the options you can add or remove “Numbered” in the optional parenthesis. 
%%The option is available for: sn-basic.bst, sn-chicago.bst%  
 
% \documentclass[pdflatex,sn-nature]{sn-jnl}% Style for submissions to Nature Portfolio journals
%%\documentclass[pdflatex,sn-basic]{sn-jnl}% Basic Springer Nature Reference Style/Chemistry Reference Style
\documentclass[pdflatex,sn-mathphys-num]{sn-jnl}% Math and Physical Sciences Numbered Reference Style
%%\documentclass[pdflatex,sn-mathphys-ay]{sn-jnl}% Math and Physical Sciences Author Year Reference Style
% \documentclass[pdflatex,sn-aps]{sn-jnl}% American Physical Society (APS) Reference Style
%%\documentclass[pdflatex,sn-vancouver-num]{sn-jnl}% Vancouver Numbered Reference Style
%%\documentclass[pdflatex,sn-vancouver-ay]{sn-jnl}% Vancouver Author Year Reference Style
%%\documentclass[pdflatex,sn-apa]{sn-jnl}% APA Reference Style
%%\documentclass[pdflatex,sn-chicago]{sn-jnl}% Chicago-based Humanities Reference Style

%%%% Standard Packages
%%<additional latex packages if required can be included here>

\usepackage{graphicx}%
\usepackage{multirow}%
\usepackage{amsmath,amssymb,amsfonts}%
\usepackage{amsthm}%
\usepackage{mathrsfs}%
\usepackage[title]{appendix}%
\usepackage{xcolor}%
\usepackage{textcomp}%
\usepackage{manyfoot}%
\usepackage{booktabs}%
\usepackage{algorithm}%
\usepackage{algorithmicx}%
\usepackage{algpseudocode}%
\usepackage{listings}%



%%%% Revision markup (round 2) -- same conventions as in the 063 project
% \added{new text}            additions (underlined in the marked build)
% \deleted{old text}          deletions (struck out)
% \replaced{new text}{old text}   NOTE: changes-package order: NEW first, OLD second
% \changed{revised text}      changed text in equations (blue), where \added/\deleted cannot be used
% \mycomment{who}{text}       red working notes, removed in the clean build
% The clean manuscript is generated with:  python3 make_clean_manuscript.py manuscript_revised_template.tex manuscript_revised_clean.tex
\newcommand{\mycomment}[2]{\textbf{\textcolor{red}{[#1: #2]}}}
\newcommand{\changed}[1]{\begingroup\color{blue}#1\endgroup}
\newcommand{\pending}[1]{\textcolor{red}{[#1]}}% no placeholders left after part 11 (all w_s = 2L numbers final); macro kept in case new data arrive
\newcommand{\prelim}[1]{\textcolor{red}{#1}}% provisional (red, no brackets): normalised 3D values, which await the homogeneous references with elastic side layers (strip runs submitted 2026-10-28)


\usepackage{pgfplots}


\DeclareRobustCommand{\mathdefault}[1]{#1}

\usepackage{nicefrac}
\usepackage{xspace}

\newcommand{\vMLabel}{[\textbf{vM}]\xspace}
\newcommand{\EqJcLabel}{[\textbf{Eq$\boldsymbol{J_c}$}]\xspace}
\newcommand{\EqSigmaStarLabel}{[\textbf{Eq$\boldsymbol{\sigma^*}$}]\xspace}
\newcommand{\EqPiStarLabel}{[\textbf{Eq$\boldsymbol{\Pi^*}$}]\xspace}

%%%%

%%%%%=============================================================================%%%%
%%%%  Remarks: This template is provided to aid authors with the preparation
%%%%  of original research articles intended for submission to journals published 
%%%%  by Springer Nature. The guidance has been prepared in partnership with 
%%%%  production teams to conform to Springer Nature technical requirements. 
%%%%  Editorial and presentation requirements differ among journal portfolios and 
%%%%  research disciplines. You may find sections in this template are irrelevant 
%%%%  to your work and are empowered to omit any such section if allowed by the 
%%%%  journal you intend to submit to. The submission guidelines and policies 
%%%%  of the journal take precedence. A detailed User Manual is available in the 
%%%%  template package for technical guidance.
%%%%%=============================================================================%%%%

%% as per the requirement new theorem styles can be included as shown below
\theoremstyle{thmstyleone}%
\newtheorem{theorem}{Theorem}%  meant for continuous numbers
%%\newtheorem{theorem}{Theorem}[section]% meant for sectionwise numbers
%% optional argument [theorem] produces theorem numbering sequence instead of independent numbers for Proposition
\newtheorem{proposition}[theorem]{Proposition}% 
%%\newtheorem{proposition}{Proposition}% to get separate numbers for theorem and proposition etc.

\theoremstyle{thmstyletwo}%
\newtheorem{example}{Example}%
\newtheorem{remark}{Remark}%

\theoremstyle{thmstylethree}%
\newtheorem{definition}{Definition}%

\raggedbottom
%%\unnumbered% uncomment this for unnumbered level heads

\begin{document}

\title[Influence of J2-plasticity on the effective crack resistance of a simplified porous material in a numerical experiment]{Influence of J2-plasticity on the effective crack resistance of a simplified porous material in a numerical experiment}

%%=============================================================%%
%% GivenName	-> \fnm{Joergen W.}
%% Particle	-> \spfx{van der} -> surname prefix
%% FamilyName	-> \sur{Ploeg}
%% Suffix	-> \sfx{IV}
%% \author*[1,2]{\fnm{Joergen W.} \spfx{van der} \sur{Ploeg} 
%%  \sfx{IV}}\email{iauthor@gmail.com}
%%=============================================================%%

\author*[1]{\fnm{Alexander} \sur{Schl\"{u}ter}}
\email{alexander.schlueter@tu-darmstadt.de}

\author[1]{\fnm{Ronjit} \sur{Medda}}

\author[1]{\fnm{Ralf} \sur{M\"{u}ller}}

\affil*[1]{%
  \orgdiv{Institut f\"ur Mechanik, FB Bau- und Umweltingenieurwissenschaften},
  \orgname{TU Darmstadt},
  \orgaddress{%
    \street{Franziska-Braun-Stra{\ss}e 7},
    \city{Darmstadt},
    \postcode{64287},
    \country{Germany}
  }
}

%%==================================%%
%% Sample for unstructured abstract %%
%%==================================%%

\abstract{  
The effective crack resistance of a simplified porous material with ductile bulk behavior is determined in a numerical experiment. Following Hossain et al.~\cite{hossain_effective_2014}, a resolved periodic arrangement of pores is subjected to surfing boundary conditions, the crack evolution is computed with a phase-field model without a prescribed crack path, and the maximum of the J-integral over one pore-to-pore cycle is taken as the effective crack resistance. A von~Mises plasticity model is compared with two linear elastic models with the same fracture energy or the same one-dimensional strength. In two dimensions, the effective crack resistance increases with the wall width and saturates. Relative to the homogeneous material, pores toughen the linear elastic materials by up to a factor of two, the plastic material by at most 1.34, and weaken it for thin walls. The energy balance and the ligament stress state before re-nucleation identify two regimes: thin ligaments yield across their width and re-nucleate at the one-dimensional strength, wide ligaments re-nucleate at a constrained notch at a higher stress. In three dimensions, with spherical pores, the crack front advances beside the pores without unstable ligament failure, and \prelim{the effective crack resistance does not exceed the homogeneous value for wall widths up to $2L$}.
}

%%================================%%
%% Sample for structured abstract %%
%%================================%%

% \abstract{\textbf{Purpose:} The abstract serves both as a general introduction to the topic and as a brief, non-technical summary of the main results and their implications. The abstract must not include subheadings (unless expressly permitted in the journal's Instructions to Authors), equations or citations. As a guide the abstract should not exceed 200 words. Most journals do not set a hard limit however authors are advised to check the author instructions for the journal they are submitting to.
% 
% \textbf{Methods:} The abstract serves both as a general introduction to the topic and as a brief, non-technical summary of the main results and their implications. The abstract must not include subheadings (unless expressly permitted in the journal's Instructions to Authors), equations or citations. As a guide the abstract should not exceed 200 words. Most journals do not set a hard limit however authors are advised to check the author instructions for the journal they are submitting to.
% 
% \textbf{Results:} The abstract serves both as a general introduction to the topic and as a brief, non-technical summary of the main results and their implications. The abstract must not include subheadings (unless expressly permitted in the journal's Instructions to Authors), equations or citations. As a guide the abstract should not exceed 200 words. Most journals do not set a hard limit however authors are advised to check the author instructions for the journal they are submitting to.
% 
% \textbf{Conclusion:} The abstract serves both as a general introduction to the topic and as a brief, non-technical summary of the main results and their implications. The abstract must not include subheadings (unless expressly permitted in the journal's Instructions to Authors), equations or citations. As a guide the abstract should not exceed 200 words. Most journals do not set a hard limit however authors are advised to check the author instructions for the journal they are submitting to.}

\keywords{Ductile Fracture, Porous Metals, Effective Crack Resistance, Phase Field Model of Ductile Fracture, Numerical Experiment}

%%\pacs[JEL Classification]{D8, H51}

%%\pacs[MSC Classification]{35A01, 65L10, 65L12, 65L20, 65L70}

\maketitle

\section{Introduction}\label{sec:intro}

Porous and other heterogeneous materials are increasingly used in load-bearing components. Additive manufacturing processes such as wire arc additive manufacturing (WAAM) allow materials to be combined and microstructures to be designed through the process parameters, as demonstrated by Hauser et al.~\cite{hauser_porosity_2021}, and pores are a typical defect of these processes. The strength assessment of such defect-rich materials is often based on fracture mechanics, which requires a fracture resistance parameter of the material. Relevant quantities are the critical stress intensity factor $K_{Ic}$ of Irwin's criterion $K_I = K_{Ic}$, the critical energy release rate $G_c$ of Griffith's criterion $G = G_c$ and, more generally, a critical value $J_c$ of the J-integral $J$ in the criterion
\begin{equation}
    J = J_c,
    \label{eq:fracture_criterion_J}
\end{equation}
where \( J_c \) denotes the critical value of the J-integral and serves as a material parameter.
Such criteria
remain applicable up to the onset of crack growth under monotonic loading even
in the presence of plasticity, as discussed by Gross and Seelig~\cite{gross_fracture_2011}.
For a heterogeneous material, the fracture resistance entering eq.~(\ref{eq:fracture_criterion_J}) is not a property of the bulk material alone but depends on the microstructure. How it can be determined for a porous material with a ductile bulk material is the subject of this work.



For elastic and plastic properties, the step from the microstructure to a macroscopic material model is provided by homogenization. Its central concept is the representative volume element, defined by Hill~\cite{hill_elastic_1963} as the smallest sample of a heterogeneous material that can be treated as a statistically uniform continuum. Fracture does not fit into this framework in a direct way. A crack is a localized and evolving discontinuity, its growth is a non-stationary process that may be unstable, and the separation of length scales that underlies the notion of a representative volume element is lost once the crack interacts with individual microstructural features. For this reason, the term \textsl{fracture toughness} is used in this work only for a macroscopic material property in the sense of eq.~(\ref{eq:fracture_criterion_J}), whereas the value obtained from a numerical experiment under specific conditions, i.e., the maximum energetic driving force required to force crack advance through the considered microstructure, is called the \textsl{effective crack resistance}. Three classes of approaches to the latter quantity can be distinguished in the literature.

Rigorous results exist for the homogenization of surface energies of Griffith type. Braides, Defranceschi and Vitali~\cite{braides_homogenization_1996} proved that free-discontinuity functionals with a periodically oscillating surface energy, as they appear in the variational formulation of brittle fracture by Francfort and Marigo~\cite{francfort_revisiting_1998} and Bourdin et al.~\cite{bourdin_variational_2008}, $\Gamma$-converge to a functional with a homogenized and, in general, anisotropic surface energy. Cagnetti, Dal Maso, Scardia and Zeppieri~\cite{cagnetti_stochastic_2019} extended this result to stationary random microstructures. Schneider~\cite{schneider_fftbased_2020} and Ernesti and Schneider~\cite{ernesti_fast_2021,ernesti_computing_2022} use the resulting cell formula as a variational definition of an effective crack energy and compute it with FFT-based solvers, also for anisotropic fracture resistance. Michel and Suquet~\cite{michel_merits_2022} assess the merits and limitations of this definition. They relate it to limit-load problems and show that it is a lower bound of a more physically motivated effective fracture energy, namely the energy dissipated along the actual crack path through a unit cell. The homogenized surface energy is the energy of the minimal crack surface weighted by the local toughness. It does not account for the elastic interaction of the crack with the microstructure, for crack trapping and the unstable crack jumps that follow re-nucleation, or for inelastic bulk dissipation, which are the mechanisms that govern toughening in many materials, see the discussion by Lebihain~\cite{lebihain_brittle_2021}.

A second class of approaches describes the interaction of a crack front with the microstructure directly. Gao and Rice~\cite{gao_firstorder_1989} and Bower and Ortiz~\cite{bower_threedimensional_1991} analyzed the trapping of a crack front by arrays of tough obstacles with perturbation methods and with numerical methods for three-dimensional cracks, and Lebihain, Leblond and Ponson~\cite{lebihain_effective_2020} derived the effective toughness of periodic heterogeneous materials including out-of-plane excursions of the crack. These approaches are restricted to linear elastic bulk behavior and to small contrasts of the toughness or small perturbations of a planar crack.

The third class is the numerical experiment of Hossain et al.~\cite{hossain_effective_2014}. A region of the material that resolves the microstructure is subjected to so-called surfing boundary conditions, which enforce complete separation at a prescribed macroscopic crack speed, and the maximum of the energetic driving force recorded at the boundary is taken as the effective crack resistance. Since the crack evolution inside the region is computed with a phase-field model, crack trapping, re-nucleation, branching and unstable crack jumps are included without restriction. The method was developed further by Hsueh et al.~\cite{hsueh_stress_2018} for layered materials and by Brach et al.~\cite{brach_anisotropy_2019} for anisotropic effective crack resistance.
Recent contributions underline that this remains an active field of research. Liu and Shen~\cite{liu_effective_2025} propose a phase-field formulation for anisotropic effective fracture toughness that accounts for a quasi-worst-case translation of a representative volume element. Aranda and Segurado~\cite{aranda_cracklength_2025} introduce a crack-length control technique for phase-field fracture in FFT homogenization and use it to estimate effective fracture properties of composites and porous materials. Man et al.~\cite{man_collaborative_2025} study how crack guiding and trapping in bioinspired interfaces control effective toughness, while Henry~\cite{henry_phasefield_2025} shows for dynamic crack propagation in disordered heterogeneous materials that the spatial distribution of soft inclusions can strongly affect the apparent fracture energy even when the effective elastic properties are similar.
A recent related development is the variable stiffness boundary condition of Hao et al.~\cite{hao_variable_2026}, which builds on the surfing boundary condition but translates a simple uniform remote displacement into a non-uniform load through an engineered stiffness gradient, thereby facilitating implementation in experiments and commercial finite-element software. The limitations of this third class are of a different nature than those of the first two. The recorded value depends on the loading of the experiment and on the size and the content of the resolved region, and the representativity of this region in the sense of Hill is not guaranteed but has to be assessed for every microstructure. The result is therefore an operational quantity of the experiment and not a homogenized material property. With few exceptions, the contributions cited above are limited to linear elastic material behavior and to two-dimensional problems.

Phase-field models of fracture, see the review by Ambati et al.~\cite{ambati_review_2015}, have been coupled with plasticity in several ways. Gradient damage models coupled with plasticity and their relation to cohesive cracks were analyzed by Alessi, Marigo and Vidoli~\cite{alessi_gradient_2014}. Variational formulations at finite strains were proposed by Miehe, Aldakheel and co-workers~\cite{miehe_phasefield_ductile_2016,miehe_phasefield_porous_2016} and by Borden et al.~\cite{borden_phasefield_2016}, and small-strain formulations in which the plastic energy is degraded together with the elastic energy were proposed by Ambati, Gerasimov and De Lorenzis~\cite{ambati_ductile_2015} and by Kuhn et al.~\cite{kuhn_phase_2016} and Noll et al.~\cite{noll_3d_2020}. In the latter formulation, a cubic degradation function allows the yield stress and the hardening modulus to be prescribed directly, and it is the formulation used in this work.









In previous studies by the authors of this contribution~\cite{schlueter_ermittlung_2025,schluter_determination_2025},
the aforementioned framework was employed to determine the effective crack resistance of a simplified porous material,
assuming linear-elastic bulk behavior. It was demonstrated that
the renucleation mechanism of cracks that are trapped in pores
determines the effective crack resistance
and can act as a toughening factor if the pores are
spaced far enough apart. A saturation of the effective crack resistance was obtained above
a specific minimum distance between pores. This minimum distance depends on the
length at which cracks renucleate after being trapped, which among other factors also
depends on the pore shape.
These observations should be understood as results for the specific porous material and numerical experiment investigated in those studies, not as a general statement for arbitrary porous materials.

As a first step toward accounting for ductile bulk behavior in this type of numerical experiment,
the effect of a non-linear elastic Ramberg--Osgood-type material law was also studied by Schlüter et al.~\cite{schluter_influence_2025}.

For the specific idealized porous material and the crack-growth experiment following Hossain et al.,
the results indicate that the energy dissipation capacity represented by the material law,
or the fracture resistance of the bulk material in that setting,
has a stronger influence on the computed effective crack resistance than the tensile strength.

In the present work, we aim to extend this analysis
toward additively manufactured metals, which typically exhibit pronounced ductile behavior.
Our previous work considered only reversible elastic material behavior. The Ramberg--Osgood law is non-linear and corresponds to a deformation-theory description of plasticity, which agrees with von~Mises plasticity under monotonic, proportional loading, but it does not capture unloading. Here, the effect of true energy dissipation due to plasticity is studied with a von Mises plasticity model.
Since local unloading inevitably occurs during crack propagation,
even under monotonically increasing external loads,
this more closely matches true ductile fracture behavior.

The open questions addressed in this work are the following. First, how does plastic dissipation change the effective crack resistance obtained in the numerical experiment, and by which mechanisms? Compared with a physical experiment, the simulation gives access to the complete energy balance of the specimen and to the stress state at the location of crack re-nucleation. Both are used here to partition the energy dissipated per pore-to-pore cycle into crack surface energy, stored plastic energy and viscous dissipation and to characterize the stress state at re-nucleation, and the sensitivity of the result with respect to the yield stress and the hardening modulus is examined. Second, does the two-dimensional idealization of cylindrical pores carry over to spherical pores in three dimensions? Third, by how much does the result depend on the initial crack and on the size of the resolved region? The last point does not replace a representativity study, but it quantifies the dependence for the microstructure considered.

This paper is structured as follows.
Section~\ref{sec:numericalExperiment} introduces the numerical experiment
proposed by Hossain et al.~\cite{hossain_effective_2014}.
Section~\ref{sec:simulation_phase_field_models} presents the phase-field
fracture formulation for both linear-elastic and von Mises-type plasticity material models.
Section~\ref{sec:2Dresults} focuses on the computation
of the effective crack resistance for a simplified porous material
under two-dimensional plane strain assumptions
 using both constitutive descriptions, followed by the energy partition of the experiment, the stress state before re-nucleation, the sensitivity with respect to the plasticity parameters, and the influence of the isotropic approximation of the embedding and of the size of the resolved region.
Lastly, we extend the model to three dimensions and compare our findings for the von Mises-type plasticity model
for both the two-dimensional and three-dimensional cases.




\section{\texorpdfstring{A Numerical Experiment for the Crack Resistance of a Resolved Porous Region}{A Numerical Experiment for the Crack Resistance of a Resolved Porous Region}}
\label{sec:numericalExperiment}

% In this work, we limit ourselves to the description of brittle fracture and linear elastic materials. 

\begin{figure}[htb]
    \centering
    \resizebox{0.6\textwidth}{!}{\input{pics/num_experiment.pdf_t}}
    \caption{%
        Sketch of the numerical experiment introduced by Hossain et al.~\cite{hossain_effective_2014}. 
        Rather than conducting a macroscopic experiment, a region that resolves the microstructural features is fractured in a numerical simulation while the energetic driving force required for fracture is measured.
    }
    \label{fig:numerical_experiment}
\end{figure}

The numerical experiment proposed by Hossain et al.~\cite{hossain_effective_2014} for the determination of an effective crack resistance is conceptually closer to physical experiments than to the classical homogenization concepts commonly presented in micromechanics textbooks such as those by Gross and Seelig~\cite{gross_fracture_2011} and Nemat-Nasser and Hori~\cite{nemat-nasser_micromechanics_2013}. 
It should therefore be understood as an operational definition of an effective crack resistance for a prescribed numerical experiment, not as a classical homogenization result.
The objective of this approach is the determination of an effective energetic crack resistance, denoted by $J_{x,c}^{\text{eff}}$. This quantity can subsequently be used as an input quantity in a macroscopic fracture criterion such as eq.~(\ref{eq:fracture_criterion_J}).
Such a use is justified only to the extent that the fracture process analyzed by eq.~(\ref{eq:fracture_criterion_J}) is representative of the specific material and loading situation considered in the numerical experiment.

The numerical experiment is structured as follows.
Instead of performing a physical experiment, a finite region of the material
 is selected that explicitly resolves the relevant microstructural features and whose fracture behavior is to be captured by the effective energetic crack resistance.
This region is referred to as the resolved porous region, or simply the resolved region, in the following. The term representative volume element is deliberately avoided. Hill~\cite{hill_elastic_1963} defines a representative volume element as the smallest sample of a heterogeneous material that can be treated as a statistically uniform continuum, so that its effective properties are independent of the applied boundary conditions. The resolved region considered here is a sample of the microstructure whose representativity in this sense is not established, and the separation of length scales underlying Hill's definition is lost once a crack of the size of the region develops.
This region is highlighted in dark gray in Fig.~\ref{fig:numerical_experiment}.
The corresponding length scale is referred to as the \textsl{microscale}, together with its associated \textsl{microstructure}.
The scale at which this effective energetic crack resistance is applied is denoted as the \textsl{macroscale}.

The resolved region is embedded within a surrounding homogeneous matrix with the effective elastic stiffness of a unit cell of the microstructure. For the periodic arrangements of circular or spherical pores considered in this work, this stiffness has square (two dimensions) or cubic (three dimensions) symmetry. In the simulations it is approximated by an isotropic stiffness with the effective Young's modulus $E^{\text{eff}}$ and Poisson's ratio $\nu^{\text{eff}}$. The degree of anisotropy of the unit cells and the influence of this approximation on the effective crack resistance are quantified in Sections~\ref{sec:2D_effective_medium} and~\ref{sec:3D_effective_medium}.
This surrounding material facilitates the evaluation of surface integrals at the boundary and is shown in light gray in Fig.~\ref{fig:numerical_experiment}.
The transition from the resolved region to the homogeneous embedding is a material discontinuity, and the corners of the resolved region act as geometric stress concentrators. In the two-dimensional simulations of Section~\ref{sec:2Dresults}, the corners are located at least one cell width $w_c$ away from the crack plane; in the three-dimensional specimens of Section~\ref{sec:3D_results}, thin layers of the embedding also lie on the crack plane, see Section~\ref{sec:3D_effective_medium}. The embedding remains linear elastic, so that the J-integral evaluated on the outer boundary is not affected by the local stress concentration at the corners. The influence of the corners on the crack pattern is discussed together with the results.
The union of the resolved region and the surrounding matrix is considered as a continuum~$\Omega$ with boundary $\partial\Omega$.
Macroscopic crack growth, corresponding to complete separation of the resolved region along a prescribed direction, is enforced through appropriate boundary conditions.

In this work, only mode~I loading is considered.
Following the methodology of Hossain et al.~\cite{hossain_effective_2014}, so-called surfing boundary conditions are employed.
A fictitious macroscopic crack tip, illustrated by the black dot in Fig.~\ref{fig:numerical_experiment}, is propagated along a prescribed direction with velocity $\dot{x}_{\text{bc}}$, 
corresponding to the $x$-direction in the figure.
Displacements associated with mode~I crack opening are prescribed at the boundary $\partial\Omega$ using a fictitious 
stress intensity factor $K_I^{\text{bc}}$, namely
\begin{align}
    \begin{split}
        u_x &= \dfrac{K_I^{\text{bc}}}{2\mu_{\text{bc}}}  \sqrt{\dfrac{r}{2\pi}}\left(\kappa - \cos{\varphi}\right)\cos{\dfrac{\varphi}{2}},\\
        u_y &= \dfrac{K_I^{\text{bc}}}{2\mu_{\text{bc}}}  \sqrt{\dfrac{r}{2\pi}}\left(\kappa - \cos{\varphi}\right)\sin{\dfrac{\varphi}{2}},\\
        \text{where} \quad \kappa &= 3 - 4\nu_{\text{bc}} \quad \text{and} \quad \mu_{\text{bc}} = \dfrac{E_{\text{bc}}}{2\left(1 + \nu_{\text{bc}}\right)}.
        \label{eq:surfing-bcs}
    \end{split}
\end{align}

Here, $r$ and $\varphi$ denote the polar coordinates of a system centered at the fictitious 
crack tip located at $x = \dot{x}_{\text{bc}} t$ and $y = 0$, and $\kappa$ is the Kolosov constant for plane strain.
Here, $t$ denotes the pseudo-time, or loading parameter, used to advance the fictitious crack tip.
The boundary conditions are fully defined by the 
parameters $\dot{x}_{\text{bc}}$, $K_I^{\text{bc}}$ and the auxiliary fictitious material parameters  
$E_{\text{bc}}$, and $\nu_{\text{bc}}$.

The next step consists of simulating the evolution of the microscopic crack pattern until complete separation of the resolved region is achieved.
The formulation must allow unrestricted crack evolution, since the interaction between microcracks and microstructural features governs the resulting fracture process, as indicated by the red crack paths in Fig.~\ref{fig:numerical_experiment}.

An initial microscopic crack is introduced in this work and is depicted by the thick red line in Fig.~\ref{fig:numerical_experiment}.

Throughout the simulation, the macroscopically observable energetic driving force is evaluated at the boundary, indicated by the dashed blue contour in Fig.~\ref{fig:numerical_experiment}.
This quantity is computed in terms of the J-Integral
\begin{equation}
    J_x = \int_{\partial\Omega}\left(\psi_{\text{el}}(\boldsymbol{\varepsilon})\boldsymbol{1} - \left(\nabla\boldsymbol{u}\right)^T\boldsymbol{\sigma}\right)\boldsymbol{n} \, \mathrm{d}A \cdot \boldsymbol{e}_x,
    \label{eq:J-integral}
\end{equation}
where $\boldsymbol{e}_x$ denotes the unit vector in the macroscopic crack propagation direction. Hence, the index $x$ indicates both the crack-growth direction and the direction onto which the J-Integral is projected.
The strain energy density entering the J-Integral is denoted by $\psi_{\text{el}}$ and is specified in Section~\ref{sec:phase_fiel_model_ductile}, the linearized strain tensor by $\boldsymbol{\varepsilon} = \text{sym}(\nabla{\boldsymbol{u}})$, the displacement field by $\boldsymbol{u}$, the Cauchy stress by $\boldsymbol{\sigma}$, and the outward unit normal by $\boldsymbol{n}$.
Only small strains are considered in this work.
The recorded value depends on the crack pattern inside the resolved microstructure. Following the observation reported by Hossain et al.~\cite{hossain_effective_2014}, the homogeneous surrounding material is assumed not to change the value of the boundary J-integral. Under this assumption, the J-integral evaluated at the outer boundary of the embedding equals the macroscopic driving force acting on the resolved region.

The effective crack resistance, denoted by $J_{x,c}^{\text{eff}}$, is defined by Hossain et al.~\cite{hossain_effective_2014} as the maximum value attained by the J-Integral during the complete separation process, i.e., during characteristic crack growth through the resolved region,
\begin{equation}
    J_{x,c}^{\text{eff}} = \max_t J_x(t).
\end{equation}
Hossain et al.\ take the maximum over the complete separation process. In this work, $J_{x,c}^{\text{eff}}$ is evaluated as the maximum of $J_x$ over one pore-to-pore cycle in the interior of the specimen, see Sections~\ref{sec:2D_fracture_behavior} and~\ref{sec:2D_energy_partition}, which excludes the build-up of the K-field at the start and the influence of the end of the specimen.
The index $c$ indicates that the crack-resistance value is obtained as a maximum, i.e., is \textsl{critical}, whereas the index $x$ is retained to avoid implying isotropic crack resistance.

The motivation for this definition is that loading levels below $J_{x,c}^{\text{eff}}$ lead to crack arrest within the simulation and 
therefore do not result in macroscopic crack propagation, 
understood as full separation of the resolved region.
This definition does not imply that $J_{x,c}^{\text{eff}}$ becomes independent of the size of the resolved region and of the position of the initial crack when the region is enlarged. Whether this is the case for a given class of microstructures and loading paths is a question of representativity that is not settled by the definition.

For the numerical experiment introduced by Hossain et al.~\cite{hossain_effective_2014}, two formulations of the surfing boundary conditions were compared in~\cite{hossain_effective_2014, bohnen_simulation_2024}. For the specific cases investigated there, including the formulation used in the present work, no relevant difference in the predicted value of $J_{x,c}^{\text{eff}}$ was found. This should not be understood as a general proof of independence from the boundary conditions.
Whether $J_{x,c}^{\text{eff}}$ is independent of the initial microscopic crack configuration and of the size of the resolved region cannot be taken for granted and has to be tested for every microstructure. In the present work, only one macroscopic crack-growth direction is considered. The influence of the length of the initial crack and of the number of resolved pore rows on $J_{x,c}^{\text{eff}}$ is probed in Section~\ref{sec:2D_representativity}. A representativity study in the sense of Hill~\cite{hill_elastic_1963}, i.e., a statistical assessment over microstructure realizations, sample sizes and loading conditions, is outside the scope of this work. The results should therefore be read as a controlled numerical experiment on the influence of J2-plasticity for the selected simplified microstructure, rather than as the determination of a generally valid homogenized fracture toughness.


\section{Simulation of fracture with phase field models}\label{sec:simulation_phase_field_models}



\subsection{A phase field model of ductile fracture}\label{sec:phase_fiel_model_ductile}

This section introduces the phase-field models and the numerical ingredients used in the subsequent numerical experiments.
In order to model fracture in ductile materials under quasi-static loading conditions, we employ the model proposed in~\cite{noll_3d_2020}. 
Related variational finite-strain phase-field formulations for ductile fracture and porous plasticity were developed by Miehe, Aldakheel and co-workers~\cite{miehe_phasefield_ductile_2016,miehe_phasefield_porous_2016,miehe_phasefield_implementation_2017}.
The constitutive behavior of the considered continuum $\Omega$ with boundary $\partial\Omega$ in the unfractured state is described by von Mises plasticity with isotropic linear hardening. 
The notation and numerical treatment of the underlying small-strain $J_2$ plasticity model follow standard computational plasticity concepts, see, for example, Simo and Hughes~\cite{simo_computational_1998}.
Small deformation theory is assumed to be sufficient.

In this case, the total linearized strain
\begin{equation}
    \boldsymbol{\varepsilon} = \dfrac{1}{2}\left(\nabla{\boldsymbol{u}} + \nabla{\boldsymbol{u}}^T\right)
\end{equation}
with $\boldsymbol{u}$ denoting the displacement field can be decomposed into a reversible elastic strain $\boldsymbol{\varepsilon}^e$ caused by the current stress $\boldsymbol{\sigma}$ and a plastic strain $\boldsymbol{\varepsilon}^p$ that remains after complete unloading as
\begin{equation}
  \boldsymbol{\varepsilon} =  \boldsymbol{\varepsilon}^e + \boldsymbol{\varepsilon}^p.
\end{equation}

To model linear isotropic strain hardening, a scalar internal state variable $\alpha$ is introduced. 
The yield stress is postulated as
\begin{equation}
    \sigma_Y(\alpha) = \sigma_y + H\alpha
\end{equation}
with $H$ denoting the hardening modulus and $\sigma_y$ the initial yield stress. 
Plastic deformation is assumed to be driven exclusively by the deviatoric part of the stress state, which is typical for metals. 
Accordingly, a $J_2$ yield function is introduced as
\begin{equation}
    f(\boldsymbol{\sigma},\alpha) = \|\boldsymbol{s}\| - \sqrt{\dfrac{2}{3}}\left(\sigma_y+H\alpha\right), 
    \quad \text{with} \quad 
    \boldsymbol{s} = \text{dev}\left(\boldsymbol{\sigma}\right) = \boldsymbol{\sigma}-\dfrac{1}{3}\text{tr}\left(\boldsymbol{\sigma}\right)\boldsymbol{1},
\end{equation}
where $\text{tr}(*)$ denotes the trace operator and $\boldsymbol{1}$ is the second order identity tensor.
The yield function $f$ determines the onset of plastic flow through the condition
\begin{equation}
    f(\boldsymbol{\sigma},\alpha) = 0 \quad \text{if yielding occurs}.
\end{equation}

Consistent with the principle of maximum dissipation, the evolution equations of the plastic variables are given by
\begin{equation}
    \dot{\boldsymbol{\varepsilon}}^p = \dot{\gamma}\boldsymbol{N},
    \label{eq:evolution_eps_p}
\end{equation}
and
\begin{equation}
    \dot{\alpha} = \sqrt{\dfrac{2}{3}}\,\dot{\gamma},
    \label{eq:evolution_alpha}
\end{equation}
where
\begin{equation}
    \boldsymbol{N} = \dfrac{\boldsymbol{s}}{\|\boldsymbol{s}\|}.
\end{equation}
In addition, the Karush--Kuhn--Tucker conditions
\begin{equation}
    f(\boldsymbol{\sigma},\alpha) \le 0, 
    \quad \dot{\gamma} \ge 0, 
    \quad f\,\dot{\gamma} = 0,
\end{equation}
must be satisfied to ensure a physically consistent description of plasticity.

To model fracture, an additional field variable $s(\boldsymbol{x},t) \in [0,1]$ is introduced. 
The value $s=0$ corresponds to fully fractured material, while $s=1$ represents intact material sufficiently far from cracks. 
In the transition region, the field varies smoothly between these two values, providing a diffuse representation of cracks in a continuum setting.
Irreversibility of crack growth is enforced in both the ductile and the linear elastic phase-field formulations by fixing the phase field at fully fractured points: if $s(\boldsymbol{x},t_{\boldsymbol{x}})=0$ for the first time $t_{\boldsymbol{x}}$ at a point $\boldsymbol{x}$, then a Dirichlet condition $s(\boldsymbol{x},t)=0$ is imposed for all later pseudo-times $t>t_{\boldsymbol{x}}$.

Using this phase field variable, the total energy of the continuum $\Omega$ is given by
\begin{equation}
    E_{\rm{ep}} = \int_{\Omega} \left[ 
    \psi_{\rm{el}}(\boldsymbol{\varepsilon},\boldsymbol{\varepsilon}^p,s) 
    + \psi_{\rm{fr}}(s,\nabla{s}) 
    + \psi_{\rm{pl}}(\alpha,s)
    \right] \, \mathrm{d}V,
    \label{eq:E_ep_total}
\end{equation}
where
\begin{equation}
    \psi_{\rm{el}}(\boldsymbol{\varepsilon},\boldsymbol{\varepsilon}^p,s)
    = g(s)\,w_{\rm{el}}(\boldsymbol{\varepsilon},\boldsymbol{\varepsilon}^p)
    = g(s)\dfrac{1}{2}
    \left(\boldsymbol{\varepsilon}-\boldsymbol{\varepsilon}^p\right)
    {:}
    \left(\boldsymbol{\mathbb{C}}{:}
    \left(\boldsymbol{\varepsilon}-\boldsymbol{\varepsilon}^p\right)\right).
    \label{eq:psi_el}
\end{equation}
Here, $\boldsymbol{\mathbb{C}}$ denotes the elastic stiffness tensor of the undamaged material. Throughout this work, the bulk material is assumed to be isotropic, so that $\boldsymbol{\mathbb{C}}$ is characterized by the bulk modulus $K$ and the shear modulus $\mu$, or equivalently by the Lam\'e parameters $\lambda = K - \nicefrac{2}{3}\mu$ and $\mu$. The stiffness is degraded in the presence of cracks through the degradation function
\begin{equation}
    g(s) = \beta\left(s^3 - s^2\right) + 3s^2 - 2s^3 + \eta,
    \label{eq:deg_grad}
\end{equation}
with $\beta = 0.2$ controlling the slope of $g(s)$ at $s=1$, and $\eta = 10^{-5}$ a small residual stiffness. A positive value of $\eta$ keeps the stiffness positive definite in fully fractured regions, $s = 0$, so that the equilibrium problem remains well-posed there. It is therefore a part of the model and not a device of the numerical solution procedure.
The cubic degradation function is used because it allows the prescribed yield stress and hardening parameter to be implemented in the phase-field formulation in a direct manner.
Kuhn et al.~\cite{kuhn_phase_2016} discuss that, for the standard quadratic degradation function, these parameters require a reinterpretation in order to obtain effective material properties that can be related to measured data.

To account for crack driving forces arising from the accumulation of plastic deformation, as expected for ductile fracture, the plastic contribution to the energy is also degraded according to
\begin{equation}
    \psi_{\rm{pl}}(\alpha,s) = g(s)\left(\dfrac{1}{2}H\alpha^2 + \sigma_y\alpha\right).
    \label{eq:psi_pl}
\end{equation}


Finally, the fracture energy density per unit volume is given by
\begin{equation}
    \psi_{\rm{fr}} = J_c
    \underbrace{\left(
    \dfrac{(1-s)^2}{4\epsilon} + \epsilon\|\nabla{s}\|^2
    \right)}_{\bar{A}},
    \label{eq:psi_fr}
\end{equation}
which represents the energy dissipated due to fracture. 

Here, $\bar{A}$ approximates the crack surface upon integration over a volume, $J_c$ denotes the material-specific fracture energy per unit crack surface area, and $\epsilon$ controls the width of the diffuse crack zone.


For the linear elastic case, Ambrosio--Tortorelli-type regularizations approximate the sharp-crack variational formulation in the sense of $\Gamma$-convergence as $\epsilon \to 0$. In this sense, the regularized fields $(\boldsymbol{u},s)$ approximate the displacement field and the crack set of the corresponding sharp-crack formulation. The convergence is that of the energy functionals and of their minimizers, and the region of low values of $s$ concentrates on the crack set as $\epsilon \to 0$. No pointwise convergence of $s$ is implied, see Kuhn~\cite{kuhn_numerical_2013} and Chambolle~\cite{chambolle_approximation_2004}.

A thermodynamically consistent evolution equation for the phase field variable $s$, governing crack growth, is derived using a variant of the Coleman--Noll procedure as described in~\cite{noll_3d_2020},
\begin{align}
    \begin{split}
    \dot{s} = -M \Bigg[
    g'(s)\left(w_{\rm{el}} + \left(\sigma_y+\dfrac{1}{2}H\alpha\right)\alpha\right)
    - J_c\left(2\epsilon\Delta{s} + \dfrac{1-s}{2\epsilon}\right)
    \Bigg]
    \quad \text{in} \ \Omega, \\
    \nabla{s}\cdot{\boldsymbol{n}} = 0 \quad \text{on} \ \partial\Omega.
    \label{eq:evolution_s}
    \end{split}
\end{align}

The parameter $M$ is the kinetic coefficient of the evolution equation~(\ref{eq:evolution_s}). It sets the time scale on which the phase field relaxes towards a stationary point of the energy, and for large values of $M$ the evolution of $s$ approximates a rate-independent process. The energy dissipated by this viscous evolution, $\int_\Omega \dot{s}^2/M\,\mathrm{d}V$ per unit time, is part of the energy balance of the model and is quantified for the numerical experiments in Section~\ref{sec:2D_energy_partition}.
Static equilibrium is assumed in the form
\begin{equation}
    \text{div}(\boldsymbol{\sigma}) = \boldsymbol{0} \quad \text{in} \ \Omega, 
    \quad 
    \boldsymbol{\sigma}\boldsymbol{n} = \boldsymbol{t}^* 
    \quad \text{on} \ \partial\Omega_t,
    \label{eq:equilibrium}
\end{equation}
where the stress tensor $\boldsymbol{\sigma}$ follows from eq.~(\ref{eq:psi_el}) for the isotropic stiffness as
\begin{equation}
    \boldsymbol{\sigma} = g(s)\left[
    K\,\text{tr}(\boldsymbol{\varepsilon})\,\boldsymbol{1}
    + 2\mu\left(\text{dev}(\boldsymbol{\varepsilon}) - \boldsymbol{\varepsilon}^p\right)
    \right].
    \label{eq:sigma_plast}
\end{equation}

The evolution equations for the primary fields $\boldsymbol{u}$ and $s$, together with the evolution equations for the plastic variables and the Karush--Kuhn--Tucker conditions, define the coupled phase field formulation for ductile fracture.


\subsection{A phase field model for brittle fracture}
\label{sec:phasefield_elastic}

In order to judge the influence of ductile material behavior on the effective crack resistance, 
we also consider a phase field model of brittle fracture that assumes linear elastic material behavior. 
Details of this model are described, e.g., by~\cite{kuhn_continuum_2010}. The governing equations 
are very similar to eqs.~(\ref{eq:evolution_s}) and (\ref{eq:equilibrium}) but do not contain
any contributions of the plastic variables:

\begin{equation}
    \dot{s} = -M \left[g'(s)\psi_{\rm{el}}  - J_c\left(2\epsilon\Delta{s} + \dfrac{1-s}{2\epsilon}\right)\right] \quad \text{in} \, \Omega, \quad \nabla{s}\cdot{\boldsymbol{n}} = 0\quad\text{on}\,\partial\Omega,
    \label{eq:evolution_s_elastic}
\end{equation}

\begin{equation}
    \text{div}(\boldsymbol{\sigma}) = \boldsymbol{0} \quad \text{in} \, \Omega, \quad{\boldsymbol{\sigma} \boldsymbol{n} =\boldsymbol{t}^*} \quad\text{on}\,\partial\Omega_t
    \label{eq:equilibrium_elastic}
\end{equation}

with 

\begin{equation}
    \boldsymbol{\sigma} = g(s)\left[K \text{tr}\boldsymbol{\varepsilon}\boldsymbol{1} + 2\mu\text{dev}(\boldsymbol{\varepsilon})\right].
\end{equation}


% The origins of phase field models like the one presented in the previous section can be found in the original Griffith-criterion, which states that cracks can grow if the energy released by crack growth is equal to the energy required to create the new crack surface which is a material parameter. 

\subsection{Numerical implementation}
\label{sec:numerical_implementation}

The spatial discretization uses first-order Lagrange finite elements on triangular meshes in two dimensions and tetrahedral meshes in three dimensions. Along the expected crack path, the mesh is locally refined such that the element edge length $\Delta h$ is at most half the regularization length $\epsilon$ in two dimensions, and it is usually chosen much smaller. In three dimensions, the element size in the refined band is of the order of $\epsilon$, and the mesh dependence of the result is quantified in Section~\ref{sec:3D_fracture_behavior}. The specific values of $\Delta h$ and $\epsilon$ used in the individual simulations are stated together with the corresponding numerical setups below. The evolution equation for the phase field, eq.~(\ref{eq:evolution_s}), is discretized in pseudo-time by an implicit backward Euler scheme.

The weak forms of the coupled boundary-value problems are implemented in FEniCSx using \texttt{dolfinx} version 0.7.3~\cite{fenicsx_documentation}. The displacement field, the phase field, and, where applicable, the plastic internal variables form one monolithic nonlinear system per time step. The nonlinear residual and its derivative are assembled by FEniCSx, and the resulting system is solved with the PETSc-based Newton solver provided by FEniCSx. The linear subproblems arising within Newton's method are solved with a direct sparse LU factorization (MUMPS) through PETSc. The time step is adapted automatically during the simulation. The maximum number of Newton iterations per time step is set to eight. If a step does not converge within this limit, the time step is reduced by a factor of two and the step is repeated. If convergence is reached in fewer than a prescribed number of iterations (four in the two-dimensional simulations, five for $w_s = 4L$, and five to seven in the three-dimensional simulations), the next time step is doubled, subject to the prescribed upper bound for the respective simulation.







\subsection{Parameter sets}
\label{sec:parameter_sets}
In the following sections, the influence of the material law and various parameters on the fracture behavior of the phase field model,
as well as on the effective crack resistance of a porous material, is investigated.
To isolate the impact of individual parameters, three specific parameter sets are studied.
These are defined for a given reference fracture resistance $J_c^0$, Lam\'{e} parameter $\mu_0$, and length scale $L$ as: 
\begin{itemize}
\setlength{\itemindent}{4em}


\item[\vMLabel:]
\parbox[t]{\dimexpr\linewidth-4em\relax}{
the von Mises-type plasticity model of section~\ref{sec:phase_fiel_model_ductile},
with $\sigma_y=1.0\sqrt{\nicefrac{J_c^0\mu_0}{L}}$, $H=0.22\mu_0$,
$\lambda=\mu_0$ and $J_c=1.0J_c^0$,
}

\item[\EqJcLabel:]
\parbox[t]{\dimexpr\linewidth-4em\relax}{
the linear elastic model of section~\ref{sec:phasefield_elastic}
with $J_c=1.0J_c^0$ and $\lambda=\mu_0$,
}

\item[\EqSigmaStarLabel:]
\parbox[t]{\dimexpr\linewidth-4em\relax}{
the linear elastic model of section~\ref{sec:phasefield_elastic}
with fracture resistance $J_c=0.5523J_c^0$ and $\lambda=\mu_0$.
}
\end{itemize}

The parameter set \EqJcLabel is chosen to have the same prescribed microscopic fracture resistance $J_c$ as \vMLabel, whereas \EqSigmaStarLabel is chosen such that the one-dimensional tensile strength equals that of \vMLabel.


\subsection{Relation to an Energetic Fracture Criterion}
\label{sec:steady_growth_griffith}

\begin{figure}[htb]
    \centering
    \resizebox{0.7\textwidth}{!}{\includegraphics[]{pics/griffith_deformed.png}}
    \vspace{0.2cm}
    \caption{%
        Deformed contour plot of the phase field $s$ for the parameter set \EqJcLabel in a homogeneous strip of size $16L \times 2L$ with the same surfing boundary conditions, shown to illustrate the enforced crack path. The deformation is scaled by a factor of 0.3.
    }
    \label{fig:griffith_deformed}
\begin{picture}(0,0)
    \put(133,86){$x$}
    \put(-155,136){$y$}
    \put(3,68){$s$}
\end{picture}    
\end{figure}

\begin{figure}[htb]
    \centering
    \includegraphics[width=0.7\textwidth]{pics/PAPER_rev2_Jx_vs_xct_homogeneous_references.pdf}% rev2: regenerated from 075 (S2_solid_vM, S2_strip_EqJc, S2_strip_EqSig), replaces PAPER_01_all_Jx_vs_xct_pf__vM.png (which contained a synthetic continuation of the final descent of the vM curve)
    \caption{%
        J-Integral $J_x$ vs.\ crack tip position $x_{ct}$ for the three parameter sets during crack propagation through the homogeneous specimen of size $16L \times 20L$ without pores. The shaded band marks the central part of the crack path over which the numerical fracture resistance $J_c^{\text{num}}$ is averaged.
    }
    \label{fig:PAPER_01_all_Jx_vs_xct_pf_2}
\end{figure}

\begin{figure}[htb]
    \centering
    \resizebox{0.85\textwidth}{!}{\includegraphics{pics/alpha_t_6.2495.png}}\\[4pt]
    {\small $\alpha$}~\resizebox{0.22\textwidth}{!}{\includegraphics{pics/legend_alpha.png}}
    \vspace{0.2cm}
    \caption{%
        Plastic variable $\alpha$ for the parameter set \vMLabel at $t = 6.25\nicefrac{L}{\dot{x}_{\rm bc}}$ in a homogeneous strip of size $16L \times 2L$ with the same surfing boundary conditions. Plastic deformation stays confined to the crack path.
    }
    \label{fig:plasticity_vm_jeff}
\end{figure}

In order to highlight the connection of the presented phase field model to the fracture criterion, eq.~(\ref{eq:fracture_criterion_J}), a simulation was performed in which crack growth was enforced using the surfing boundary conditions described in section~\ref{sec:numericalExperiment}. The computational domain is the simulation domain of the porous specimens of Section~\ref{sec:2Dresults} for the wall width $w_s = L$, i.e., a rectangle of size $(8w_c) \times 20L = 16L \times 20L$, but without pores and with an initial crack of length $2L$. It is meshed by the same rule as the porous specimens, with an element edge length of $\Delta h = 0.02L$ along the crack path and a coarser mesh in the far field, under plane strain assumptions with $\lambda_0 = \mu_0$, $M = 1000.0 \nicefrac{\dot{x}_{\text{bc}}}{J_c^0}$ and $\epsilon = 0.1L$. For the von~Mises set, the yield stress is assigned to the whole domain, so that the specimen is homogeneous for all three parameter sets. The boundary condition parameters were $\nu_{\text{bc}} = 0.25$, $E_{\text{bc}} = 2.5\mu_{0}$, and $K_I^{\text{bc}} = \sqrt{J_c^0 \cdot 2.5\mu_{0}}$, which enforce crack propagation in the $x$-direction, as shown in Fig.~\ref{fig:griffith_deformed}.

These simulations were conducted using the three parameter sets introduced in section~\ref{sec:parameter_sets}. 
As discussed in section~\ref{sec:numericalExperiment}, the J-Integral, eq.~(\ref{eq:J-integral}), evaluated at the boundary of the domain quantifies the energetic driving force for crack propagation. 
The J-Integral rises during a short transient at the start, remains nearly constant over the largest part of the crack path and drops when the crack approaches the end of the domain, as shown in Fig.~\ref{fig:PAPER_01_all_Jx_vs_xct_pf_2}.

As established by Tanné et al.~\cite{tanne_crack_2018}, the actual numerical fracture resistance $J_c^{\text{num}}$ obtained from finite element simulations is typically higher than the specified parameter $J_c$. For a given material law, this deviation primarily depends on the ratio $\nicefrac{\Delta{h}}{\epsilon}$, see also Tanné et al.~\cite{tanne_crack_2018}. In the current setup, the following values are obtained from the plateau of $J_x$, evaluated as the mean over the central part of the crack path between $30$ and $70\,\%$ of the final position of the K-field, shaded in Fig.~\ref{fig:PAPER_01_all_Jx_vs_xct_pf_2}:
\begin{itemize}
    \item $J_c^{\text{num}} \approx 1.425J_c^0$ for the von Mises material model \vMLabel,
    \item $J_c^{\text{num}} \approx 1.088J_c^0$ for \EqJcLabel,
    \item $J_c^{\text{num}} \approx 0.604J_c^0$ for \EqSigmaStarLabel,
\end{itemize}
The ratio of the values of \EqJcLabel and \EqSigmaStarLabel, $1.088/0.604 = 1.80$, agrees with the ratio of the prescribed fracture resistances, $1/0.5523 = 1.81$, as expected for the linear elastic model, whose energies scale with $J_c$, see Section~\ref{sec:2D_fracture_behavior}.
The corresponding three-dimensional homogeneous reference simulations, carried out for the parameter set \vMLabel only, yield $J_c^{\text{num}} = 1.42J_c^0$, and $1.445J_c^0$ for the longer specimen used for $w_s = 2L$, see Section~\ref{sec:3D_fracture_behavior}.
It is worth noting that the von Mises material model yields a significantly higher numerical fracture resistance than the linear elastic model at the same specified $J_c$. The evolution of plastic variables remains largely limited to the crack path, as can be observed in Fig.~\ref{fig:plasticity_vm_jeff}, so that the linear elastic formulation of the J-Integral, eq.~(\ref{eq:J-integral}), remains a good approximation even in the presence of plasticity. However, the ductile material law still has a significant effect on the crack driving force. In general, the larger the fracture resistance $J_c$, the higher the required crack driving force $J_x$ must be to sustain steady crack growth.

These simulations confirm that pre-existing cracks propagate according to an energetic criterion equivalent to the criterion in Eq.~(\ref{eq:fracture_criterion_J}): a crack advances if the energetic driving force equals the fracture resistance. Throughout this work, such behavior is referred to as energetically driven crack growth. In Section~\ref{sec:2Dresults}, the effective crack resistance of each parameter set is compared with the numerical fracture resistance $J_c^{\text{num}}$ of the same set.

\subsection{Stress driven crack nucleation}
\label{sec:stress_driven_nucleation}

\begin{figure}[htb]
    \centering
    \resizebox{0.7\textwidth}{!}{\includegraphics[]{pics/1D_bar.png}}
    \caption{%
       One-dimensional bar subjected to a displacement load $u_0$.
    } 
    \label{fig:setup_crack_nucleation}
\end{figure}

\begin{figure}[htb]
    \centering
    \begin{minipage}[t]{0.49\textwidth}
        \centering
        \resizebox{\linewidth}{!}{%
            \input{pics/PAPER_reaction_forces_1D_rev2.pgf}% rev2: recomputed in 077 (beta = 0.2, J_c = 0.5523 for [Eq sigma*])
        }
        \caption{Stress $\sigma$ vs.\ strain $\varepsilon$. The dots indicate the fracture stress $\sigma^*$, whereas the vertical dashed lines display the fracture strain $\varepsilon^*$.}
        \label{fig:reaction_forces}
    \end{minipage}%
    \hfill
    \begin{minipage}[t]{0.49\textwidth}
        \centering
        \resizebox{\linewidth}{!}{%
            \input{pics/PAPER_energy_vs_displacement_1D_rev2.pgf}% rev2: recomputed in 077
        }
        \caption{Energy $\Pi$ vs.\ applied displacement.}
        \label{fig:energy_vs_displacement}
    \end{minipage}
\end{figure}

In addition to satisfying the fracture criterion, eq.~(\ref{eq:fracture_criterion_J}), for pre-existing cracks, crack nucleation in previously undamaged material can also be modelled with phase-field models. Whether the nucleation loads obtained in this way are quantitatively predictive for a given material is a separate question, which depends on the regularization length and on the degradation function, see Tann\'e et al.~\cite{tanne_crack_2018}. In the present work, nucleation is used as a model mechanism whose characteristic stress and energy are known analytically.

In the one-dimensional linear elastic case with the cubic degradation function of eq.~(\ref{eq:deg_grad}) and $\beta = 0$, it can be shown analytically that crack nucleation occurs once the stress level reaches a critical value slightly below
\begin{equation}
    \sigma_c = \dfrac{81}{50}\sqrt{\dfrac{2\mu J_c}{15.0\epsilon}},
    \label{eq:sig_c}
\end{equation}
see~\cite{schlueter_phase_2018}. For the ductile phase-field fracture model, also employing the cubic degradation function of eq.~(\ref{eq:deg_grad}) with $\beta = 0$, the critical strain is given by
\begin{equation}
    \varepsilon_c = \sqrt{\dfrac{1}{3} \dfrac{E+H}{EH}
    \left[\left(\dfrac{\sigma_y^2}{H} + \dfrac{5}{18}\dfrac{J_c}{\epsilon}\right)
    + \sqrt{4\left(\dfrac{\sigma_y^2}{H}\right)^2
    + \dfrac{8}{9} \dfrac{\sigma_y^2 J_c}{H\epsilon}
    + \dfrac{25}{324} \left(\dfrac{J_c}{\epsilon}\right)^2}\,\right]},
    \label{eq:ductile_fracture_strain}
\end{equation}
for $E = 2\mu_0$, and the corresponding critical stress then follows from eq.~(\ref{eq:sigma_plast}), see Noll et al.~\cite{noll_3d_2020}. Strictly speaking, these relations hold only in one dimension. In two and three dimensions, cracks nucleate once a stress level close to $\sigma_c$ is exceeded in a region whose size is of the order of several multiples of $\epsilon$, as shown for the quadratic degradation function by Kuhn~\cite{kuhn_numerical_2013} and for the cubic degradation function by Schl\"uter~\cite{schlueter_phase_2018}. For the porous specimens studied in this work, the stress state in the ligament ahead of a trapped crack tip immediately before re-nucleation is recorded and compared with the one-dimensional values in Section~\ref{sec:2D_mechanisms}. Due to the relevance of this stress threshold, this mechanism is referred to as \textsl{stress-driven crack nucleation}, in contrast to the \textsl{energetically driven crack growth} mechanism discussed previously.

Stress-driven crack nucleation is demonstrated using the one-dimensional setup shown in Fig.~\ref{fig:setup_crack_nucleation}. The domain consists of an initially undamaged bar (i.e., $s = 1$ everywhere) of length $2L$, subjected to a linearly increasing displacement load $u_0$ applied at both ends. The field equations are reduced to one dimension by setting $\lambda = 0$, $u_y = 0$, and $u_z = 0$. The material parameters are $\mu_0$, $M = 1000.0\,\nicefrac{\dot{x}_{\text{bc}}}{J_c^0}$, $\epsilon = 0.1L$, and an element edge length of $\Delta h = 0.02L$.

Simulations were conducted for the three parameter sets introduced in Section~\ref{sec:parameter_sets}. To intentionally trigger crack nucleation at $x = 0$, the fracture resistance $J_c$ was locally reduced by 1\% within a region extending from $x=-7.5\Delta h$ to $x=7.5\Delta h$, i.e.\ $|x| \le 0.15L$.

In all cases, fracture was observed to occur at $x = 0$ once a critical load level was 
reached. For the same energetic fracture resistance $J_c$, the linear elastic model with parameters \EqJcLabel yielded a significantly higher fracture stress of $\sigma^* = 1.72\sigma_y$ than the von~Mises-type model \vMLabel, which resulted in $\sigma^* = 1.28\sigma_y$. However, the fracture strain $\varepsilon^*$, defined as the strain at which the stress drops to zero, was smaller in the linear elastic case, with $\varepsilon^* \approx 1.07\sqrt{\nicefrac{J_c^0}{\mu_0 L}}$ for \EqJcLabel compared to $\varepsilon^* \approx 1.17\sqrt{\nicefrac{J_c^0}{\mu_0 L}}$ for \vMLabel, as shown in Fig.~\ref{fig:reaction_forces}.

To match the fracture stress of the von~Mises model in a linear elastic simulation, the fracture resistance was reduced to $J_c = 0.5523J_c^0$ in \EqSigmaStarLabel, since the fracture stress of the linear elastic model is proportional to $\sqrt{J_c}$ for fixed $\epsilon$ and degradation function. This adjustment resulted in the same fracture stress of $\sigma^* = 1.28\sigma_y$ as for the von~Mises model, but a significantly lower fracture strain of only $\varepsilon^* \approx 0.80\sqrt{\nicefrac{J_c^0}{\mu_0 L}}$.

The work of the internal forces up to the strain $\varepsilon$ is defined as
\begin{equation}
   \Pi(\varepsilon) = \int_{-L}^{L}
   \left[\int_{0}^{\varepsilon} \boldsymbol{\sigma} :
   \mathrm{d}\tilde{\boldsymbol{\varepsilon}}\right] \mathrm{d}x,
   \label{eq:total_energy}
\end{equation}
and is evaluated at the fracture strain $\varepsilon = \varepsilon^*$. The evolution of this energy for all three simulations is shown in Fig.~\ref{fig:energy_vs_displacement}. The highest dissipated energy is observed in the linear elastic simulation \EqJcLabel, with $\Pi^* \approx 2.19J_c^0 L$, followed by the von~Mises simulation \vMLabel with $\Pi^* \approx 2.12J_c^0 L$. In contrast, the linear elastic simulation \EqSigmaStarLabel dissipates only $\Pi^* \approx 1.23J_c^0 L$.





\section{Effective Crack Resistance of a Simplified Two-Dimensional Porous Material}
\label{sec:2Dresults}

\subsection{Equivalent Elastic Medium and Simulation Domain}
\label{sec:2D_effective_medium}



\begin{figure}[htb]
    \centering
    \resizebox{0.9\textwidth}{!}{\includegraphics[]{pics/setup.png}}
    \vspace{10mm}
    \caption{%
        a) Simulation domain for the numerical experiment for $w_s = 1.0L$. Red indicates regions with bulk material stiffness, whereas blue shows regions with effective stiffness.
        b) Typical crack pattern.
        c) Employed triangular finite element mesh.
    }
    \label{fig:setup_crack_growth}
    \begin{picture}(0,0)
        \put(-173,207){$y$}
        \put(-63,136){$x$}
        \put(-180,60){a)}
        \put(-59,60){b)}
        \put(62,60){c)}
    \end{picture}
\end{figure}

\begin{figure}[htb]
    \centering
    \resizebox{0.45\textwidth}{!}{\input{pics/PAPER_rev2_xct_xbc_vs_t_vM_ws1.0.pgf}}% rev2: regenerated from 075 (S1_vM_ws1.0), replaces PAPER_00_xct_pf_vs_xct_KI_vonMises
    \caption{%
        Position of the surfing boundary conditions $x_{\text{bc}}$ and the crack tip position $x_{\text{ct}}$ for $w_s = L$ and the parameter set \vMLabel.
        The horizontal axis denotes the pseudo-time $t$ normalized by the time scale $L/\dot{x}_{\mathrm{bc}}$.
    }
    \label{fig:PAPER_00_xct_pf_vs_xct_KI_diamond}
\end{figure}

\begin{figure}[htb]
    \centering
    \includegraphics[width=0.72\textwidth]{pics/PAPER_rev2_Jx_vs_xct_ws1.0_three_sets.pdf}% rev2: regenerated from 075, replaces PAPER_clean_Jx_vs_xct_SIG_Jc_vonMises.png (+ holes_s_middle_row.png inset)
    \caption{%
        J-Integral vs.\ crack tip position $x_{ct}$ for $w_s = 1.0L$ and the three parameter sets. Dashed lines mark the pore edges. The picture at the top shows the phase field $s$ of the middle pore row after complete separation (\vMLabel).
    }
    \label{fig:PAPER_02_all_Jx_vs_xct_pf_diamond_holes}
\end{figure}


\begin{figure}[htb]
    \centering
    \includegraphics[width=\textwidth]{pics/PAPER_rev2_Jx_vs_xct_sections_three_sets.pdf}% rev2: one figure replaces PAPER_04_Jx_vs_xct_all_vonMises.png, PAPER_03A_Jx_vs_xct_all_linear_elastic_gc10.png and PAPER_04_Jx_vs_xct_all_SIG.png (Figs. 11-13 of v2.0)
    \caption{J-Integral $J_x$ vs.\ crack tip position $x_{ct}$ between pores two and four for varying $w_s$ and the parameter sets a) \vMLabel, b) \EqJcLabel and c) \EqSigmaStarLabel.}
    \label{fig:PAPER_rev2_Jx_sections_three_sets}
\end{figure}

\begin{figure}[htb]
    \centering
    \resizebox{0.9\textwidth}{!}{\input{pics/PAPER_rev2_Jeff_vs_wsteg_three_sets.pgf}}% rev2: replaces PAPER_06b_Jx_vs_wsteg_three_datasets.png
    \caption{Effective crack resistance $J_{x,c}^{\text{eff}}$ vs.\ wall width $w_s$: a) normalized by the reference fracture resistance $J_c^0$, b) normalized by the numerical fracture resistance $J_c^{\text{num}}$ of the respective parameter set, Section~\ref{sec:steady_growth_griffith}.}
    \label{fig:PAPER_06b_Jx_vs_wsteg_three_datasets}
\end{figure}

In order to investigate the effect of the material law on the effective crack resistance of a porous material, a numerical experiment is performed, as illustrated in Fig.~\ref{fig:setup_crack_growth}a). The domain consists of circular pores with diameter $L$, centered in square cells of edge length $w_c = w_s + L$. A $6 \times 3$ array of these cells is assembled, as depicted by the red regions in Fig.~\ref{fig:setup_crack_growth}a), where $w_s$ denotes the distance between pores. Plane strain conditions are assumed, i.e.,
\begin{equation}
    \boldsymbol{u} = u(x,y)\boldsymbol{e}_x + v(x,y)\boldsymbol{e}_y.
\end{equation}

The effective initial Young’s modulus $E_0^{\text{eff}}$ and Poisson’s ratio $\nu_0^{\text{eff}}$ of the unit cells are computed in separate preliminary simulations by prescribing linear displacement boundary conditions $\boldsymbol{u}^* = \boldsymbol{\varepsilon}_0^k \boldsymbol{x}$ at the boundary of a single unit cell.
Here, $k$ denotes the load case, i.e., the prescribed macroscopic strain for which the linear displacement boundary condition is applied. Several load cases are required to determine the full stiffness tensor. The pore surface is traction free.
The resulting stress field $\boldsymbol{\sigma}_{\text{micro}}$ and the corresponding volume-averaged stresses
\begin{equation}
    \left< \boldsymbol{\sigma} \right> = \dfrac{1}{V} \int_V \boldsymbol{\sigma}_{\text{micro}} \, \mathrm{d}V
\end{equation}
are evaluated via finite element simulations. The effective stiffness tensor $\left< \boldsymbol{\mathbb{C}} \right>$ is obtained from
\begin{equation}
    \left< \boldsymbol{\sigma} \right>^k = \left< \boldsymbol{\mathbb{C}} \right>  \left< \boldsymbol{\varepsilon} \right> = \left< \boldsymbol{\mathbb{C}} \right>  \boldsymbol{\varepsilon}_0^k.
\end{equation}
Employing Voigt notation indicated by underscores $\underline{(*)}$, the two-dimensional strain tensor
\[
\boldsymbol{\varepsilon} = 
\begin{bmatrix} \varepsilon_{xx} & \varepsilon_{xy} \\ \varepsilon_{xy} & \varepsilon_{yy} \end{bmatrix}
\]
is represented as the three-component vector
\begin{equation}
\underline{\boldsymbol{\varepsilon}} = 
\begin{bmatrix} \varepsilon_{xx} \\ \varepsilon_{yy} \\ 2 \varepsilon_{xy} \end{bmatrix},
\end{equation}
where the factor of two accounts for engineering shear strain. Similarly, the stress tensor
\[
\boldsymbol{\sigma} = 
\begin{bmatrix} \sigma_{xx} & \sigma_{xy} \\ \sigma_{xy} & \sigma_{yy} \end{bmatrix}
\]
is represented as
\begin{equation}
\underline{\boldsymbol{\sigma}} = 
\begin{bmatrix} \sigma_{xx} \\ \sigma_{yy} \\ \sigma_{xy} \end{bmatrix}.
\end{equation}
Using these definitions, we can define the unit strains in Voigt form as
\begin{equation}
\underline{\boldsymbol{\varepsilon}}_0^1 =
\begin{bmatrix} 1 \\ 0 \\ 0 \end{bmatrix}, \quad
\underline{\boldsymbol{\varepsilon}}_0^2 =
\begin{bmatrix} 0 \\ 1 \\ 0 \end{bmatrix}, \quad
\underline{\boldsymbol{\varepsilon}}_0^3 =
\begin{bmatrix} 0 \\ 0 \\ 1 \end{bmatrix}.
\end{equation}
The corresponding effective stiffness tensor is then computed from the resulting stresses as
\begin{equation}
    \left< \underline{\boldsymbol{\mathbb{C}}} \right> =
    \begin{bmatrix}
         \sigma_{xx}^1 & \sigma_{xx}^2 & \sigma_{xx}^3 \\
         \sigma_{yy}^1 & \sigma_{yy}^2 & \sigma_{yy}^3 \\
         \sigma_{xy}^1 & \sigma_{xy}^2 & \sigma_{xy}^3
     \end{bmatrix}.
\end{equation}


From this tensor, the effective Lamé parameters, Young’s modulus, and Poisson’s ratio are derived as
\begin{equation}
    \lambda^{\text{eff}} = \dfrac{\left< \underline{C} \right>_{12} + \left< \underline{C} \right>_{21}}{2}, \quad
    \mu^{\text{eff}} = \left< \underline{C} \right>_{33}, \quad
    E_0^{\text{eff}} = \dfrac{\mu^{\text{eff}} (3 \lambda^{\text{eff}} + 2 \mu^{\text{eff}})}{\lambda^{\text{eff}} + \mu^{\text{eff}}}, \quad
    \nu_0^{\text{eff}} = \dfrac{\lambda^{\text{eff}}}{2 (\lambda^{\text{eff}} + \mu^{\text{eff}})}.
    \label{eq:isotropic_approximation_2D}
\end{equation}
The unit cell has square symmetry, so that $\left< \underline{C} \right>_{11} = \left< \underline{C} \right>_{22}$ and the effective stiffness is characterized by the three independent constants $C_{11} = \left< \underline{C} \right>_{11}$, $C_{12} = \left< \underline{C} \right>_{12}$ and $C_{33} = \left< \underline{C} \right>_{33}$. The isotropic approximation of eq.~(\ref{eq:isotropic_approximation_2D}) uses $C_{12}$ and $C_{33}$ and replaces $C_{11}$ by $C_{12} + 2C_{33}$; the unit cell itself is therefore not isotropic, the isotropy is introduced as an approximation of its effective stiffness. The deviation from isotropy is measured by the Zener ratio $A = 2C_{33}/(C_{11} - C_{12})$, which equals one for an isotropic material. Table~\ref{tab:effective_stiffness_2D} lists the effective constants and the Zener ratio for five of the eleven wall widths, the other values lie monotonically in between. The anisotropy decreases with the wall width, from $A = 0.870$ for $w_s = 0.25L$ to $0.998$ for $w_s = 4L$, and the effective Poisson's ratio of the isotropic approximation increases from $\nu_0^{\text{eff}} = 0.19$ to $0.25$, Table~\ref{tab:effective_stiffness_2D}. The influence of an isotropic approximation of the embedding on the effective crack resistance is tested in Section~\ref{sec:2D_cubic_embedding}; unless stated otherwise, all simulations in this work use the isotropic approximation, and the full square-symmetric stiffness enters only these tests.
\begin{table}[htb]
\centering
\caption{Effective elastic constants of the two-dimensional unit cells (plane strain, in units of $\mu_0$, $\lambda_0 = \mu_0$), the isotropic approximation of eq.~(\ref{eq:isotropic_approximation_2D}) and the Zener ratio $A$ for five of the eleven wall widths. The values of the intermediate wall widths lie monotonically in between.}
\label{tab:effective_stiffness_2D}
\begin{tabular}{lcccccccc}
\toprule
$w_s/L$ & $C_{11}$ & $C_{12}$ & $C_{33}$ & $\lambda^{\text{eff}}$ & $\mu^{\text{eff}}$ & $E_0^{\text{eff}}$ & $\nu_0^{\text{eff}}$ & $A$ \\
\midrule
0.25  & 0.869 & 0.177 & 0.301 & 0.177 & 0.301 & 0.713 & 0.186 & 0.870 \\
0.375 & 1.087 & 0.241 & 0.376 & 0.241 & 0.376 & 0.900 & 0.195 & 0.890 \\
1.0   & 1.814 & 0.520 & 0.621 & 0.520 & 0.621 & 1.524 & 0.228 & 0.960 \\
2.0   & 2.361 & 0.755 & 0.794 & 0.755 & 0.794 & 1.975 & 0.244 & 0.989 \\
4.0   & 2.740 & 0.907 & 0.915 & 0.907 & 0.915 & 2.285 & 0.249 & 0.998 \\
\bottomrule
\end{tabular}
\end{table}

The $6 \times 3$ porous domain is embedded into a rectangular domain (blue) that is assigned the isotropic approximation of the effective stiffness, i.e., $E_0^{\text{eff}}$ and $\nu_0^{\text{eff}}$, with dimensions $(8 w_c) \times 20 L$. The two-dimensional simulations in this section are carried out under plane strain assumptions. The parameters $\sigma_y$ and $H$ entering the nonlinear part of the material law, eq.~(\ref{eq:sigma_plast}), are not homogenized.
Instead, it is assumed that nonlinear effects are locally relevant, while the effective behavior of each unit cell can be approximated as linear elastic with parameters $E_0^{\text{eff}}$ and $\nu_0^{\text{eff}}$. The embedding is therefore modelled as linear elastic, i.e., no yielding is admitted outside the resolved porous region.

\subsection{Fracture Behavior and Effective Crack Resistance}
\label{sec:2D_fracture_behavior}

An initial crack of length $w_c$ and horizontal orientation is prescribed along the center row of pores. The finite element edge length of the bilinear triangular elements along the expected crack path is $\Delta h = 0.02L$. The applied boundary conditions correspond to the surfing boundary conditions described in Section~\ref{sec:numericalExperiment}, with a stress intensity factor $K_I^{\text{bc}}$ that is chosen for every wall width such that the energy release rate of the prescribed K-field equals the microscopic fracture resistance of the respective parameter set, $J_{\text{bc}} = (1-\nu_{\text{bc}}^2)(K_I^{\text{bc}})^2/E_{\text{bc}} = J_c$, i.e., $(K_I^{\text{bc}})^2 = 2\mu_{\text{bc}}J_c/(1-\nu_{\text{bc}})$, and the effective elastic constants of the surrounding material, $E_{\text{bc}} = E_0^{\text{eff}}$ and $\nu_{\text{bc}} = \nu_0^{\text{eff}}$, as the auxiliary parameters of the prescribed K-field~(\ref{eq:surfing-bcs}). With this choice, the applied load is the same multiple of the microscopic fracture resistance for all wall widths, and the two linear elastic parameter sets, which differ only in $J_c$, are loaded self-similarly. The phase-field parameters are $\epsilon = 0.1L$, $\eta = 10^{-5}$, $M = 1000\,\nicefrac{v_{\text{bc}}}{J_c^0}$, and $\lambda_0 = \mu_0$. Simulations are conducted with wall widths ranging from $0.25L$ to $4.0L$ using the parameter sets introduced in Section~\ref{sec:parameter_sets}.

Figure~\ref{fig:setup_crack_growth}b) shows the final crack pattern for a wall width of $L$ using the von~Mises material law parameter set \vMLabel. Crack propagation occurs horizontally from pore to pore, with frequent crack trapping inside pores followed by re-nucleation. This behavior is confirmed by Fig.~\ref{fig:PAPER_00_xct_pf_vs_xct_KI_diamond}, which displays the microscopic crack tip position $x_{ct}$, identified as the node with the largest $x$-coordinate that fulfills $s=0$, and the position of the macroscopic crack tip $x_{\text{bc}}$ for a wall width of $1.0L$ and \vMLabel. During phases in which the crack is trapped inside a pore, this definition places $x_{ct}$ at the last fully damaged material point before the void. A different convention could assign the crack tip to the opposite side of the open pore, but the present definition is used consistently as a tracking quantity for the advance through intact material ligaments. The microscopic crack tip follows the macroscopic boundary condition in an unsteady manner. It is trapped at each pore until re-nucleation and then jumps across the pore, after which the ligament ahead fails unstably.

Figure~\ref{fig:PAPER_02_all_Jx_vs_xct_pf_diamond_holes} presents the J-Integral as a function of the microscopic crack tip position $x_{ct}$ for a wall width of $1.0L$ and all three parameter sets. Vertical lines indicate the pore edges. Aside from boundary effects at the beginning and end of the simulation, a characteristic periodic pattern is observed. While the crack tip remains trapped within a pore, $J_x$ increases to a peak value immediately prior to re-nucleation. Following re-nucleation, $J_x$ drops and decreases as the crack separates the ligament between pores and $x_{ct}$ catches up with $x_{\text{bc}}$. The peak value of $J_x$ therefore occurs directly before re-nucleation and determines the effective crack resistance according to the definition in Section~\ref{sec:numericalExperiment}. Throughout this work, $J_{x,c}^{\text{eff}}$ is evaluated as the maximum of $J_x$ over one pore-to-pore cycle in the sense of Section~\ref{sec:2D_energy_partition}, which in two dimensions coincides with the peak immediately before re-nucleation.

It is observed that the linear elastic simulation with equivalent $J_c$ (\EqJcLabel) exhibits significantly higher values of $J_x$ than the von~Mises model \vMLabel with the same microscopic $J_c$, while the latter still yields higher $J_x$ values than the linear elastic simulation with equivalent fracture stress \EqSigmaStarLabel. The peak values are $J_{x,c}^{\text{eff}} = 2.02$, $1.56$ and $1.12\,J_c^0$ for \EqJcLabel, \vMLabel and \EqSigmaStarLabel, respectively. The mechanisms behind this ordering, and behind its reversal for thin ligaments, are analyzed in Sections~\ref{sec:2D_energy_partition} and~\ref{sec:2D_mechanisms}.

For a wall width of $1.0L$, neither the equivalent one-dimensional fracture stress of the bulk material nor the microscopic energetic fracture resistance $J_c$ provides a reliable indicator of the effective crack resistance. Instead, the work of the internal forces until fracture, $\Pi^*$, in
the one-dimensional crack nucleation experiment has the same ordering as $J_{x,c}^{\text{eff}}$ at this wall width, $\Pi^* = 2.19$, $2.12$ and $1.23\,J_c^0L$ for \EqJcLabel, \vMLabel and \EqSigmaStarLabel, see also Fig.~\ref{fig:energy_vs_displacement}. This correspondence is limited to wide ligaments, as shown below.

The periodically repeating pattern of $J_x$, for example between pores two and four, is characteristic of the microstructure. This pattern is shown for varying $w_s$ in Fig.~\ref{fig:PAPER_rev2_Jx_sections_three_sets}a) for \vMLabel, b) for \EqJcLabel and c) for \EqSigmaStarLabel. Independent of the parameter set, $J_{x,c}^{\text{eff}}$ increases with increasing wall width up to a saturation value. This trend is summarized in Fig.~\ref{fig:PAPER_06b_Jx_vs_wsteg_three_datasets}, which shows the effective macroscopic crack resistance $J_{x,c}^{\text{eff}}$ versus $w_s$ for all parameter sets, normalized by $J_c^0$ in Fig.~\ref{fig:PAPER_06b_Jx_vs_wsteg_three_datasets}a) and by the numerical fracture resistance $J_c^{\text{num}}$ of the respective parameter set, Section~\ref{sec:steady_growth_griffith}, in Fig.~\ref{fig:PAPER_06b_Jx_vs_wsteg_three_datasets}b). The increase is steep for thin ligaments and levels off for wide ligaments. At $w_s = 4L$, the effective crack resistance is $1.90J_c^0 = 1.34J_c^{\text{num}}$ for \vMLabel, $2.21J_c^0 = 2.03J_c^{\text{num}}$ for \EqJcLabel and $1.22J_c^0 = 2.03J_c^{\text{num}}$ for \EqSigmaStarLabel. The pores thus toughen all three materials with respect to the homogeneous material, but relative to its own $J_c^{\text{num}}$ the von~Mises set is toughened less than the linear elastic sets for all wall widths, Fig.~\ref{fig:PAPER_06b_Jx_vs_wsteg_three_datasets}b).

The ordering \EqJcLabel $>$ \vMLabel $>$ \EqSigmaStarLabel observed for $w_s = L$ holds for all wall widths $w_s \ge 0.6L$. For thin walls, $w_s \le 0.5L$, the ductile material law reduces the effective crack resistance below the values of \EqSigmaStarLabel, the curves of \vMLabel and \EqSigmaStarLabel cross between $w_s = 0.5L$ and $0.6L$. Neither the one-dimensional tensile strength, which is equal for \vMLabel and \EqSigmaStarLabel, nor the microscopic fracture resistance $J_c$, which is equal for \vMLabel and \EqJcLabel, is therefore a consistent indicator of the effective crack resistance. The two regimes are analyzed in Sections~\ref{sec:2D_energy_partition} and~\ref{sec:2D_mechanisms}.

Finally, Fig.~\ref{fig:PAPER_06b_Jx_vs_wsteg_three_datasets} shows that saturation occurs at significantly larger wall widths for \vMLabel than for the linear elastic parameter sets. The linear elastic sets reach $95\,\%$ of their value at $w_s = 4L$ for $w_s \approx 1.5L$, the von~Mises set between $w_s = 1.5L$ and $2L$. Since $J_{x,c}^{\text{eff}}$ is determined by the peak value reached immediately prior to re-nucleation, saturation implies that re-nucleation is no longer influenced by neighboring pores. The delayed saturation observed for the ductile material law indicates that plasticity enlarges the interaction zone over which neighboring pores affect re-nucleation, due to the formation of extended plastic deformation zones prior to crack nucleation, see Section~\ref{sec:2D_mechanisms}.
The comparison of \vMLabel and \EqSigmaStarLabel in Fig.~\ref{fig:PAPER_06b_Jx_vs_wsteg_three_datasets} shows that, for $w_s \ge 0.6L$, the ductile material law leads to a substantially higher effective crack resistance even though the one-dimensional tensile strengths are equal.

The two linear elastic parameter sets differ only in the value of $J_c$. For the linear elastic phase-field model, $J_c$ is the only energy scale once the elastic constants and the lengths are fixed, and the amplitude of the K-field is proportional to $\sqrt{J_c}$ by the choice of $K_I^{\text{bc}}$. In the rate-independent limit $M \to \infty$, all energies, including $J_{x,c}^{\text{eff}}$, therefore scale with $J_c$, while the crack pattern is unchanged. The simulations confirm this: the ratio of the effective crack resistances of \EqJcLabel and \EqSigmaStarLabel lies between $1.809$ and $1.810$ for all wall widths, compared with $1/0.5523 = 1.811$ for the prescribed fracture resistances, so that the two sets fall onto one curve in Fig.~\ref{fig:PAPER_06b_Jx_vs_wsteg_three_datasets}b).

The three comparisons in Fig.~\ref{fig:PAPER_06b_Jx_vs_wsteg_three_datasets} answer different questions. For equal microscopic fracture resistance $J_c^0$, \vMLabel lies below \EqJcLabel for all wall widths. For equal one-dimensional strength, \vMLabel lies above \EqSigmaStarLabel from $w_s \approx 0.6L$ on, Fig.~\ref{fig:PAPER_06b_Jx_vs_wsteg_three_datasets}a). For equal crack resistance of the homogeneous material, Fig.~\ref{fig:PAPER_06b_Jx_vs_wsteg_three_datasets}b), \vMLabel lies below the linear elastic curve for all wall widths. No linear elastic phase-field model has both the strength and the $J_c^0$ of \vMLabel: in the linear elastic model the strength follows from $J_c$ and $\epsilon$, eq.~(\ref{eq:sig_c}), whereas the plastic model introduces $\sigma_y$ and $H$ as additional parameters, so that strength and microscopic fracture resistance can be chosen independently. The gain of \vMLabel over \EqSigmaStarLabel in Fig.~\ref{fig:PAPER_06b_Jx_vs_wsteg_three_datasets}a) stems from the material model, which already raises the crack resistance of the homogeneous material, $J_c^{\text{num}} = 1.425J_c^0$ against $1.088J_c^0$ for \EqJcLabel with the same $J_c^0$ ($+31\,\%$) and $0.604J_c^0$ for \EqSigmaStarLabel (factor $2.36$). Fig.~\ref{fig:PAPER_06b_Jx_vs_wsteg_three_datasets}b) isolates the toughening by the microstructure. The pores raise the crack resistance of the linear elastic materials by a factor of up to $2.03$ relative to the homogeneous value, that of the plastic material only by $1.34$, and below $w_s \approx 0.9L$ they even lower it, $J_{x,c}^{\text{eff}} < J_c^{\text{num}}$. A linear elastic set with $J_c^{\text{num}} = 1.425J_c^0$ would reach $2.03 \times 1.425J_c^0 = 2.89J_c^0$ at $w_s = 4L$ instead of the $1.90J_c^0$ of \vMLabel. The microstructure thus preserves only part of the gain that plasticity gives to the homogeneous material. The reason is twofold, see Sections~\ref{sec:2D_energy_partition} and~\ref{sec:2D_mechanisms}. Trapping and forced re-nucleation at the blunt pore wall raise the peak driving force of a linear elastic material far above its $J_c^{\text{num}}$, at the price of $47$ to $83\,\%$ of the dissipated energy going into the unstable snap-through of the ligament. The homogeneous plastic material, in contrast, already blunts its crack by the plastic zone, and this contribution is part of $J_c^{\text{num}}$ itself, so re-nucleation at the pore wall raises the peak driving force above $J_c^{\text{num}}$ by a smaller factor. For thin walls, yielding of the ligament keeps the peak driving force low, and it even falls below $J_c^{\text{num}}$, Section~\ref{sec:2D_mechanisms}.

\subsection{Energy Partition of the Numerical Experiment}
\label{sec:2D_energy_partition}

\begin{figure}[htb]
    \centering
    \resizebox{0.9\textwidth}{!}{\input{pics/PAPER_rev2_energy_partition_cycle.pgf}}% rev2: new figure from 075 (energy_cycles.csv)
    \caption{Energy partition of one pore-to-pore cycle vs.\ wall width $w_s$. a) Fractions of the energy $\Delta D_{\text{tot}}$ dissipated per cycle, eq.~(\ref{eq:D_tot_cycle}), that create crack surface (fracture), remain stored as degraded plastic energy (plastic) and are dissipated by the viscous evolution of $s$ (viscous). The two linear elastic sets give the same fractions, only \EqSigmaStarLabel is shown. b) Energy dissipated per unit crack advance, $G_{\text{diss}} = \Delta D_{\text{tot}}/w_c$, relative to the effective crack resistance.}
    \label{fig:energy_partition_cycle}
\end{figure}

This section quantifies where the energy supplied in the numerical experiment goes, in order to identify the dissipation channels that distinguish the ductile from the linear elastic sets. The model of Section~\ref{sec:phase_fiel_model_ductile} obeys an exact energy balance, which follows from the time derivative of the total energy, eq.~(\ref{eq:E_ep_total}), with the equilibrium condition, eq.~(\ref{eq:equilibrium}), and the evolution equation of the phase field, eq.~(\ref{eq:evolution_s}). The work $W_{\text{ext}}$ of the tractions on the outer boundary, where the surfing displacements are prescribed, is
\begin{equation}
    W_{\text{ext}} = E_{\text{el}} + E_{\text{pl}} + E_{\text{fr}} + D_{\text{visc}},
    \quad
    D_{\text{visc}} = \int_0^t \int_\Omega \dfrac{\dot{s}^2}{M}\,\mathrm{d}V\,\mathrm{d}\tilde{t},
    \label{eq:energy_balance}
\end{equation}
where $E_{\text{el}}$, $E_{\text{pl}}$ and $E_{\text{fr}}$ denote the integrals of $\psi_{\text{el}}$, $\psi_{\text{pl}}$ and $\psi_{\text{fr}}$ over $\Omega$, i.e., the three contributions to $E_{\text{ep}}$ in eq.~(\ref{eq:E_ep_total}), and $D_{\text{visc}}$ is the energy dissipated by the viscous evolution of the phase field. The plastic work is stored in the degraded form $\psi_{\text{pl}} = g(s)(\nicefrac{1}{2}H\alpha^2 + \sigma_y\alpha)$, eq.~(\ref{eq:psi_pl}), and acts as a crack driving force through $g'(s)$ in eq.~(\ref{eq:evolution_s}). Plastic work done in material that fails later is released by the degradation and reappears in $E_{\text{fr}}$ and $D_{\text{visc}}$, whereas $E_{\text{pl}}$ records the plastic energy stored in intact material. The viscous term collects the energy released in the unstable failure of a ligament, which would become kinetic energy in a dynamic formulation. For the linear elastic model, $E_{\text{pl}} = 0$. All terms are recorded in every time step. The partition below is built from the recorded sinks $E_{\text{fr}}$, $E_{\text{pl}}$, $D_{\text{visc}}$ and the stored elastic energy; the recorded external work of the two-dimensional runs is not used. The balance of eq.~(\ref{eq:energy_balance}) against the external work is verified for the three-dimensional simulations, Section~\ref{sec:3D_fracture_behavior}.

The partition is evaluated per pore-to-pore cycle. A cycle starts when the crack enters a pore and ends when it enters the next pore, i.e., it comprises the loading of one ligament while the crack tip is trapped, the re-nucleation and the failure of this ligament. The energy dissipated during a cycle is
\begin{equation}
    \Delta D_{\text{tot}} = \Delta E_{\text{fr}} + \Delta E_{\text{pl}} + \Delta D_{\text{visc}},
    \label{eq:D_tot_cycle}
\end{equation}
and $G_{\text{diss}} = \Delta D_{\text{tot}}/w_c$ is the energy dissipated per unit crack advance, averaged over the cycle. Figure~\ref{fig:energy_partition_cycle} shows the three fractions of $\Delta D_{\text{tot}}$ and the ratio $G_{\text{diss}}/J_{x,c}^{\text{eff}}$ for the cycle that contains the ligament in which $J_{x,c}^{\text{eff}}$ is attained. The energy dissipated per cycle is nearly constant from cycle to cycle, see Section~\ref{sec:2D_representativity}. The stored elastic energy decreases over a cycle, by $0.05$ to $0.6\,\Delta D_{\text{tot}}$, because the K-field applied at $t = 0$ stores elastic energy that is drawn down as the crack advances, so that the specimen is not in a strict steady state.

For the two linear elastic sets, which give the same fractions to two digits, the viscous share is large. Between $83\,\%$ of the dissipated energy for $w_s = 0.25L$ and $47\,\%$ for $w_s = 4L$ are lost in the unstable snap-through of the ligament, and only $17$ to $53\,\%$ create crack surface, since the elastic energy stored up to re-nucleation exceeds the fracture energy of the ligament by far. For the von~Mises set, the viscous share is $12$ to $37\,\%$, the fracture share $49$ to $74\,\%$, and $12$ to $16\,\%$ remain stored as plastic energy in the intact material along the crack path. The difference is largest for thin ligaments. $G_{\text{diss}}$ shows the same crossover as $J_{x,c}^{\text{eff}}$: $0.48$ vs.\ $0.71\,J_c^0$ for \vMLabel and \EqSigmaStarLabel at $w_s = 0.25L$, equal values at $w_s = 0.6L$, and $1.39$ vs.\ $0.89\,J_c^0$ at $w_s = 4L$. The ratio $G_{\text{diss}}/J_{x,c}^{\text{eff}}$, Fig.~\ref{fig:energy_partition_cycle}b), decreases from $0.93$ to $1.10$ at $w_s = 0.25L$ to $0.73$ at $w_s = 4L$, since $J_{x,c}^{\text{eff}}$ is the peak of $J_x$ in the cycle and $G_{\text{diss}}$ its average. For thin walls, the ductile material law thus reduces not only the peak driving force but also the energy dissipated per unit crack advance, which is the quantity closest to an effective fracture energy in the sense of the energy dissipated along the crack path, see Michel and Suquet~\cite{michel_merits_2022}. Part of this reduction must be attributed to the quasi-static model rather than to the physics. The energy released in the unstable snap-through of a ligament is absorbed by the viscous term of eq.~(\ref{eq:evolution_s}), i.e., by the viscous shares quoted above ($12$ to $37\,\%$ for \vMLabel, $47$ to $83\,\%$ for the linear elastic sets). In a dynamic computation, this energy would instead become kinetic energy and could drive further crack advance or additional plastic dissipation, so the reduction of $G_{\text{diss}}$ by plasticity for thin walls is in part a property of the viscous regularization. The underlying mechanism, namely that the lower peak stress of the yielded ligament suppresses re-nucleation ahead of the pore and with it the snap-through, is physical, see Section~\ref{sec:2D_mechanisms}. Resolving how the released energy would partition in reality requires dynamic simulations, which are beyond the scope of this work. The mechanism behind the crossover is identified in the next section.

\subsection{Stress State and Plastic Zone before Re-Nucleation}
\label{sec:2D_mechanisms}

\begin{figure}[htb]
    \centering
    \resizebox{0.9\textwidth}{!}{\input{pics/PAPER_rev2_ligament_state_before_renucleation.pgf}}% rev2: new figure from 075 (ligament stress records at the J_eff peak)
    \caption{State of the next intact ligament ahead of the trapped crack tip at the peak of $J_x$, i.e., immediately before re-nucleation, vs.\ wall width $w_s$. a) Maximum normal stress $\sigma_{yy,\max}$ in the ligament, normalized by the one-dimensional fracture stress $\sigma^*_{\text{1D}}$ of the respective parameter set, Section~\ref{sec:stress_driven_nucleation}. b) Ratio $\sigma_{h,\max}/\sigma_{\text{vM},\max}$ of the recorded maxima of the hydrostatic and the von~Mises stress in the ligament and, dotted, the maximum of the plastic variable $\alpha$ in the ligament for \vMLabel.}
    \label{fig:ligament_state}
\end{figure}

To relate the re-nucleation in the porous specimen to the one-dimensional results of Section~\ref{sec:stress_driven_nucleation}, the maxima of the normal stress $\sigma_{yy}$, the von~Mises stress $\sigma_{\text{vM}}$, the hydrostatic stress $\sigma_h = \nicefrac{1}{3}\,\text{tr}(\boldsymbol{\sigma})$ and the plastic variable $\alpha$ in the next intact ligament ahead of the trapped crack tip are recorded in every time step. Figure~\ref{fig:ligament_state} shows these values at the peak of $J_x$, i.e., immediately before re-nucleation, with $\sigma_{yy,\max}$ normalized by the one-dimensional fracture stress $\sigma^*_{\text{1D}}$ of the respective parameter set, $1.28\sigma_y$ for \vMLabel and \EqSigmaStarLabel and $1.72\sigma_y$ for \EqJcLabel, $\sigma_y = \sqrt{\nicefrac{J_c^0\mu_0}{L}}$, and the triaxiality measure $\sigma_{h,\max}/\sigma_{\text{vM},\max}$ formed from the recorded maxima, which are not necessarily attained at the same point.

For the two linear elastic sets, re-nucleation occurs when the maximum normal stress in the ligament reaches the one-dimensional fracture stress, $\sigma_{yy,\max} = 0.95$ to $1.02\,\sigma^*_{\text{1D}}$ for all wall widths and both values of $J_c$, at $\sigma_{h,\max}/\sigma_{\text{vM},\max} = 0.49$ to $0.53$. The nucleation in the ligament is thus stress controlled, and the one-dimensional value is a good estimate of the nucleation stress in the two-dimensional porous specimen. This confirms, for the present geometry and the cubic degradation function, the observations of Kuhn~\cite{kuhn_numerical_2013} and Schl\"uter~\cite{schlueter_phase_2018} referred to in Section~\ref{sec:stress_driven_nucleation}.

For the von~Mises set, two regimes can be distinguished in Fig.~\ref{fig:ligament_state}. For thin ligaments, $w_s \le 0.6L$, the recorded maxima indicate that the ligament yields over its width before re-nucleation: the maximum of $\alpha$ in the ligament reaches $0.66$ to $0.73$, the ratio $\sigma_{h,\max}/\sigma_{\text{vM},\max}$ is $0.67$ to $0.74$, and re-nucleation occurs at $\sigma_{yy,\max} = 0.95$ to $1.11\,\sigma^*_{\text{vM}}$, i.e., at about the one-dimensional fracture stress after a large plastic strain, as in the bar of Section~\ref{sec:stress_driven_nucleation}. Once the ligament has yielded, its tangent stiffness drops to the hardening modulus $H$, so that its opening grows with the advancing K-field at nearly constant stress, and $J_x$ at re-nucleation stays low compared with the linear elastic sets: their ligaments are loaded elastically up to the fracture stress, store correspondingly more energy, and release $47$ to $83\,\%$ of the dissipated energy per cycle in the unstable snap-through, against $12$ to $37\,\%$ for \vMLabel, Section~\ref{sec:2D_energy_partition}. Yielding lowers the peak stress ahead of the pore; it does not remove the plastic work from the energy balance, eq.~(\ref{eq:energy_balance}). In this nucleation-controlled regime, which in terms of the effective crack resistance extends up to $w_s = 0.5L$ (the curves of \vMLabel and \EqSigmaStarLabel cross between $0.5L$ and $0.6L$, Section~\ref{sec:2D_fracture_behavior}), the ductile material law reduces the effective crack resistance below that of \EqSigmaStarLabel, whose linear elastic ligament has to be loaded to a higher driving force before it reaches $\sigma^*$, Fig.~\ref{fig:PAPER_06b_Jx_vs_wsteg_three_datasets}. The energy dissipated per unit crack advance is reduced as well, Section~\ref{sec:2D_energy_partition}, since the linear elastic ligament stores more elastic energy up to failure.

For wide ligaments, $w_s \ge L$, the maximum of $\alpha$ in the ligament drops to $0.44$ to $0.60$, which indicates that the plastic zone no longer spans the ligament, and the trapped crack tip at the pore surface acts as a notch constrained by the surrounding elastic material. The ratio $\sigma_{h,\max}/\sigma_{\text{vM},\max}$ rises to $1.0$ to $1.1$, the recorded maximum of the von~Mises stress remains at the yield-stress level, $0.98$ to $1.06\,\sigma_y$, and the normal stress reaches $\sigma_{yy,\max} = 1.67$ to $1.71\sqrt{\nicefrac{J_c^0\mu_0}{L}} = 1.31$ to $1.34\,\sigma^*_{\text{vM}}$, the level of the one-dimensional fracture stress of \EqJcLabel. Since the hydrostatic stress does not enter the yield condition, the constrained ligament carries a normal stress well above the one-dimensional fracture stress before the driving force suffices for re-nucleation, and the plastic work ahead of the trapped tip has to be supplied before the crack advances. In this propagation regime, the ductile material law raises the effective crack resistance above that of \EqSigmaStarLabel, and the interaction length between neighboring pores is set by the size of the plastic zone rather than by the pore diameter or the regularization length $\epsilon = 0.1L$, consistent with the later saturation of the von~Mises set in Fig.~\ref{fig:PAPER_06b_Jx_vs_wsteg_three_datasets}.

\begin{table}[htb]
\centering
\caption{Sensitivity of the effective crack resistance of the parameter set \vMLabel with respect to the yield stress $\sigma_y$ and the hardening modulus $H$, in the nucleation regime ($w_s = 0.375L$) and in the propagation regime ($w_s = L$). Reference values $\sigma_y = 1.0\sqrt{\nicefrac{J_c^0\mu_0}{L}}$ and $H = 0.22\mu_0$, all other parameters unchanged.}
\label{tab:sensitivity_2D}
\begin{tabular}{lcccc}
\toprule
 & \multicolumn{2}{c}{$w_s = 0.375L$} & \multicolumn{2}{c}{$w_s = L$} \\
\cmidrule(lr){2-3}\cmidrule(lr){4-5}
 & $J_{x,c}^{\text{eff}}/J_c^0$ & change & $J_{x,c}^{\text{eff}}/J_c^0$ & change \\
\midrule
reference                                        & 0.718 & --        & 1.563 & --       \\
$\sigma_y = 0.75\sqrt{\nicefrac{J_c^0\mu_0}{L}}$ & 0.700 & $-2.4\,\%$  & 1.474 & $-5.7\,\%$ \\
$\sigma_y = 1.5\sqrt{\nicefrac{J_c^0\mu_0}{L}}$  & 0.981 & $+36.7\,\%$ & 1.608 & $+2.9\,\%$ \\
$H = 0.11\mu_0$                                  & 0.633 & $-11.9\,\%$ & 1.493 & $-4.4\,\%$ \\
$H = 0.44\mu_0$                                  & 0.853 & $+18.9\,\%$ & 1.658 & $+6.1\,\%$ \\
\bottomrule
\end{tabular}
\end{table}

The sensitivity of $J_{x,c}^{\text{eff}}$ with respect to the plastic parameters confirms the two regimes, Table~\ref{tab:sensitivity_2D}. Halving or doubling the hardening modulus changes the effective crack resistance by $-12\,\%$ and $+19\,\%$ at $w_s = 0.375L$, but only by $-4\,\%$ and $+6\,\%$ at $w_s = L$. Increasing the yield stress by $50\,\%$ raises it by $37\,\%$ at $w_s = 0.375L$ and by $3\,\%$ at $w_s = L$, a reduction by $25\,\%$ changes it by $-2\,\%$ and $-6\,\%$. In the nucleation regime, the effective crack resistance is governed by the stress-strain response of the yielded ligament and is therefore sensitive to $\sigma_y$ and $H$. In the propagation regime, it is governed by the re-nucleation at the constrained notch and is less sensitive to both. For the linear elastic sets, the corresponding sensitivity is the scaling with $J_c$ of Section~\ref{sec:2D_fracture_behavior}. The changes in Table~\ref{tab:sensitivity_2D} are reported in units of $J_c^0$; the fracture resistance $J_c^{\text{num}}$ of the homogeneous material was not recomputed for the varied parameters, so the table quantifies the sensitivity of the recorded value, not of the ratio $J_{x,c}^{\text{eff}}/J_c^{\text{num}}$.

\subsection{Robustness of the Effective Crack Resistance}
\label{sec:2D_cubic_embedding}\label{sec:2D_representativity}

\begin{table}[htb]
\centering
\caption{Change of the effective crack resistance of the parameter set \vMLabel (last row: \EqSigmaStarLabel) by the full square-symmetric stiffness of the embedding (Zener ratio $A$ of Table~\ref{tab:effective_stiffness_2D}), by a longer initial crack, whose tip is thereby moved relative to the pores, and by five instead of three resolved pore rows.}
\label{tab:robustness_2D}
\begin{tabular}{llc}
\toprule
variation & case & change of $J_{x,c}^{\text{eff}}$ \\
\midrule
square-symmetric embedding & $w_s = 0.375L$, $A = 0.890$ & $+0.18\,\%$ \\
                           & $w_s = L$, $A = 0.960$      & $+0.08\,\%$ \\
                           & $w_s = 4L$, $A = 0.998$     & $+0.01\,\%$ \\
initial crack $1.5w_c$     & $w_s = L$                   & $+0.02\,\%$ \\
                           & $w_s = 0.375L$              & $+0.04\,\%$ \\
initial crack $2w_c$       & $w_s = L$                   & $+0.08\,\%$ \\
                           & $w_s = 0.375L$              & $-0.06\,\%$ \\
five pore rows             & $w_s = L$                   & $-0.02\,\%$ \\
                           & $w_s = 0.375L$              & $+1.1\,\%$ \\
                           & $w_s = L$, \EqSigmaStarLabel & $-0.09\,\%$ \\
\bottomrule
\end{tabular}
\end{table}

This section tests how robust the effective crack resistance is with respect to modelling choices that are not part of the microstructure: the stiffness assigned to the embedding, the length of the initial crack and the size of the resolved region. The embedding of the resolved porous region is assigned the isotropic approximation, eq.~(\ref{eq:isotropic_approximation_2D}), of the square-symmetric stiffness of Table~\ref{tab:effective_stiffness_2D}. To test this approximation, the simulations for $w_s = 0.375L$, $L$ and $4L$ with the parameter set \vMLabel were repeated with the full stiffness $C_{11}$, $C_{12}$, $C_{33}$ assigned to the embedding. The effective crack resistance changes by at most $+0.18\,\%$, Table~\ref{tab:robustness_2D}, even for the thinnest wall of the test series, for which the anisotropy is largest, since the embedding only transmits the K-field and the J-integral is evaluated on its outer boundary, whereas the re-nucleation that determines $J_{x,c}^{\text{eff}}$ takes place in the resolved ligaments. The three-dimensional unit cells of Section~\ref{sec:3D_effective_medium} are closer to isotropic, $A \ge 0.978$ for $w_s \ge 0.375L$, so no separate three-dimensional test was performed. The corners of the resolved region act as stress concentrators, see Section~\ref{sec:numericalExperiment}. In none of the simulations did a crack nucleate at a corner, the crack pattern remained confined to the central row of pores, Fig.~\ref{fig:setup_crack_growth}b), and since the embedding is linear elastic and the J-integral is evaluated on its outer boundary, the stress concentration at the corners does not enter the recorded driving force.

The definition of $J_{x,c}^{\text{eff}}$ in Section~\ref{sec:numericalExperiment} refers to one initial crack and one resolved region. Two of the dependencies that can be tested within the present setup were probed, Table~\ref{tab:robustness_2D}. Increasing the length of the initial crack from $w_c$ to $1.5w_c$ and $2w_c$, which moves its tip from the edge of the first cell to the center of the first pore and to the middle of the first ligament, changes $J_{x,c}^{\text{eff}}$ by at most $0.1\,\%$, both at $w_s = L$ and at $w_s = 0.375L$, i.e., in the propagation and in the nucleation regime, and enlarging the resolved region from three to five rows of pores changes it by at most $1.1\,\%$. In addition, the response along the crack path is periodic. The energy dissipated per cycle, Section~\ref{sec:2D_energy_partition}, is nearly constant from the second cycle on; the first cycle, during which the K-field builds up, dissipates less, the last cycle is affected by the end of the resolved region, and the inner cycles agree within about $7\,\%$ ($10\,\%$ for \EqJcLabel at $w_s = 0.25L$). The peaks of $J_x$ in successive ligaments, Fig.~\ref{fig:PAPER_02_all_Jx_vs_xct_pf_diamond_holes}, decrease slightly from cycle to cycle, e.g.\ $1.59$, $1.58$, $1.56$ and $1.53J_c^0$ for \vMLabel at $w_s = L$. The pore-free reference specimens show the same drift of $J_x$ along the crack path, by $1$ to $3\,\%$ over the same stretch, in two and in three dimensions, so this decrease is a property of the finite-domain surfing setup rather than of the pores; the averaging window of Section~\ref{sec:steady_growth_griffith} covers it.

These tests do not establish representativity in the sense of Hill, see Section~\ref{sec:numericalExperiment}. The length of the resolved region in the crack-growth direction was varied only in three dimensions, Section~\ref{sec:3D_fracture_behavior}, where four instead of six pores in this direction change the normalized effective crack resistance by at most $1.5\,\%$. Not tested were the influence of other crack-growth directions relative to the pore lattice and of the microstructure itself, since one periodic arrangement of equal pores is considered and no statistics over microstructure realizations is attempted. The amplitude of the applied K-field is fixed by the rule of Section~\ref{sec:2D_fracture_behavior} and was not varied. The conclusions of this work are therefore restricted to the stated microstructure, crack-growth direction and loading.

 


% Consequently, dissipative materials enhance fracture toughness better than materials with high fracture stress do.

% Fig.~\ref{fig:PAPER_06a_KIc_vs_wsteg_hole&diamond}  shows the Irwin effective fracture toughness $K_{Ic}^{\text{eff}}$ obtained from (\ref{eq:kic_eff}), which is often more relevant for practical applications. The graph of $K_{Ic}^{\text{eff}}$ shows the same general trends as that of $J_{x,c}^{\text{eff}}$.

\section{Effective Crack Resistance of a Simplified Three-Dimensional Porous Material}
\label{sec:3D_results}

\subsection{Three-Dimensional Model and Equivalent Elastic Medium}
\label{sec:3D_effective_medium}

\begin{figure}[htb]
        \centering
        \includegraphics[width=0.9\textwidth, trim=10 46 2 36, clip]{pics/PAPER_rev2_3D_mesh_6x3x2_ws1.0.pdf}% rev2 (2026-10-27): mesh of the 6x3x2 run S1G_3D_vM_ws1.0 (076, all 36 pores resolved), replaces the section of the 4x3x2 mesh (PAPER_rev2_3D_mesh_section_ws0.375.pdf) and pics/mesh_with_cs.png of v2.0; trim = white margins of the PDF (content bbox 15 52 450 241 of 454 x 281 bp)
        \caption{Tetrahedral mesh of the $6 \times 3 \times 2$ specimen for $w_s = L$ ($353\,921$ tetrahedra). a) View of the domain with the surface mesh of the back face $z = z_{\max}$, which shows the refined band around the crack plane, and the surfaces of all 36 pores. b) Section through the plane of the pore centers $z = -w_c/2$ with the crack plane $y = 0$ (dashed). The polygons are the intersections of the tetrahedra with the section plane.}
        \label{fig:mesh3D}
\end{figure}

% rev2 (2026-10-23): figure of the fields s and alpha for w_s = 0.25L removed. It came from a simulation of the previous
% version (coarser w_s-scaled mesh, eta = 1e-3, beta = 0.1, w_s = 0.25L), which is not part of the fixed-mesh series; the
% statement it illustrated (continuous crack advance around the pores) is now carried by Fig. position_3D and the text.
% Re-enable with renderings of the fixed-mesh run S1F if wanted (caption below).
% \begin{figure}[t]
% \centering
% \begin{tabular}{c c}
% \llap{{a)}\hspace{0mm}}%
% \includegraphics[width=0.51\textwidth]{pics/s_t_2.9998_wsteg_0.25.png} &
% \llap{{b)}\hspace{0mm}}%
% \includegraphics[width=0.51\textwidth]{pics/alpha_t_2.9998_wsteg_0.25.png} \\
% \multicolumn{2}{c}{\small $t = 2.998 \nicefrac{L}{\dot{x}_{\rm{bc}}}$}
% \end{tabular}
% \caption{a) Evolution of the phase field $s$ and b) evolution of the internal variable $\alpha$ for $w_s=0.25L$.}
% \label{fig:one_row_two_cols}
% \begin{picture}(0,0)
%     \put(-33,186){\small$s$}
%     \put(168,186){\small$\alpha$}
% \end{picture}
% \end{figure}

The two-dimensional plane strain model discussed in Section~\ref{sec:2Dresults} provides a useful indication of the general influence of voids on the effective crack resistance of a material. However, in real porous metals, pore geometries are more accurately represented by spherical voids rather than by cylindrical pores extending in the out-of-plane direction, as assumed in the two-dimensional model.

To investigate the three-dimensional character that more closely resembles real porous metals, spherical pores of diameter $L$ are considered. These pores are centered in cubic unit cells with edge length $w_c = w_s + L$. The effective stiffness of these unit cells is computed following the procedure described in the previous section, with the exception that six independent load cases $\underline{\boldsymbol{\varepsilon}}_0^k$ are applied and the resulting effective stiffness tensor in Voigt notation 
$\left< \underline{\boldsymbol{C}} \right>$ is of size $6\times6$. The unit cell has cubic symmetry with the three independent constants $C_{11}$, $C_{12}$ and $C_{44}$ and the Zener ratio $A = 2C_{44}/(C_{11}-C_{12})$, which are listed in Table~\ref{tab:effective_stiffness_3D}. The effective Lam\'e parameters of the isotropic approximation are determined as
\begin{equation}
    \lambda^{\text{eff}} = \dfrac{\left< \underline{C} \right>_{12} + \left< \underline{C} \right>_{13} + \left< \underline{C} \right>_{21} + \left< \underline{C} \right>_{31} + \left< \underline{C} \right>_{23} + \left< \underline{C} \right>_{32}}{6}, \quad
    \mu^{\text{eff}} = \dfrac{\left< \underline{C} \right>_{44} + \left< \underline{C} \right>_{55} + \left< \underline{C} \right>_{66}}{3}.
\end{equation}
Table~\ref{tab:effective_stiffness_3D} lists the effective constants for the five wall widths considered in this section. The three-dimensional cells are closer to isotropic than the two-dimensional ones, $A \ge 0.978$ compared with $A = 0.890$ for $w_s = 0.375L$ in Table~\ref{tab:effective_stiffness_2D}, and the effective Poisson's ratio of the isotropic approximation increases from $\nu_0^{\text{eff}} = 0.232$ to $0.249$. The influence of the isotropic approximation of the embedding on the effective crack resistance is discussed in Section~\ref{sec:2D_cubic_embedding}.
\begin{table}[htb]
\centering
\caption{Effective elastic constants of the three-dimensional unit cells (in units of $\mu_0$, $\lambda_0 = \mu_0$), the isotropic approximation and the Zener ratio $A$ for the five wall widths of the three-dimensional series.}
\label{tab:effective_stiffness_3D}
\begin{tabular}{lccccccc}
\toprule
$w_s/L$ & $C_{11}$ & $C_{12}$ & $C_{44}$ & $\lambda^{\text{eff}}$ & $\mu^{\text{eff}}$ & $E_0^{\text{eff}}$ & $A$ \\
\midrule
0.375 & 2.046 & 0.608 & 0.703 & 0.608 & 0.703 & 1.732 & 0.978 \\
0.5   & 2.230 & 0.682 & 0.762 & 0.682 & 0.762 & 1.884 & 0.984 \\
0.75  & 2.483 & 0.787 & 0.841 & 0.787 & 0.841 & 2.090 & 0.992 \\
1.0   & 2.640 & 0.853 & 0.890 & 0.853 & 0.890 & 2.216 & 0.996 \\
2.0   & 2.889 & 0.955 & 0.966 & 0.955 & 0.966 & 2.413 & 1.000 \\
\bottomrule
\end{tabular}
\end{table}

A $6\times3\times2$ array of these unit cells is assembled to represent a simplified porous material. This array is centered about the plane $y=0$ and embedded into a surrounding cuboidal domain of effective stiffness with dimensions $(8 w_c)\times 20L \times (2.5 w_c)$, i.e., one cell width of embedding before and after the porous block in the $x$-direction and $w_c/4$ on each side in the $z$-direction. Thus, the unit cell used to determine the effective stiffness is cubic, whereas the surrounding homogeneous domain is a larger cuboidal embedding region assigned this effective stiffness. The two layers of thickness $w_c/4$ between the porous block and the planes $z = z_{\rm min}$ and $z = z_{\rm max}$ belong to the linear elastic embedding and lie on the crack plane $y = 0$ along the entire crack path, i.e., about $20\,\%$ of the specimen thickness fractures within embedding material with the same $J_c$. The mesh of tetrahedral finite elements is refined in a band around the expected fracture plane $y=0$, Fig.~\ref{fig:mesh3D}, with a target element size of $h = 0.1L$ for $|y| < 0.35L$, $2h$ for $|y| < 0.65L$ and at most $L$ in the remaining domain, independent of the wall width, and the surfaces of all 36 pores are resolved with $h = 0.1L$, from which the element size grows to the far-field value over a distance of $L$. The median edge length of the elements in the crack band is $1.49\epsilon$, the mean edge length on the pore surfaces is $0.10L$ in the middle row and $0.11L$ in the outer rows, and the meshes contain $184\,000$ ($w_s = 0.375L$) to $741\,000$ ($w_s = 2L$) elements. The checks reported in Section~\ref{sec:3D_fracture_behavior} use a shorter specimen with four instead of six pores in the $x$-direction, embedded in $(6 w_c)\times 20L \times (2.5 w_c)$ and meshed by the same rule in the crack band, with the outer pore rows resolved more coarsely, on which the simulation for $w_s = L$ was repeated with $h = 0.07L$ (median edge length $1.03\epsilon$, $505\,000$ elements).

An initial crack is prescribed in the plane $y=0$ from $x=x_{\rm min}$ to $x=x_{\rm min}+w_c$, where $x_{\rm min}$ denotes the minimum $x$-coordinate of the mesh. 

The phase-field parameters are chosen as $\epsilon = 0.1L$, $\eta = 10^{-5}$, $M = 1000 \nicefrac{v_{\text{bc}}}{J_c^0}$, and $\lambda_0=\mu_0$. The applied boundary conditions are given by the surfing boundary conditions described in Section~\ref{sec:numericalExperiment}, with $K_I^{\text{bc}}=\sqrt{2.5J_c^0\mu_0}$ and, as in the two-dimensional case, the effective elastic constants of the surrounding material as the auxiliary parameters $E_{\text{bc}}$ and $\nu_{\text{bc}}$. The surfing displacements with $u_z = 0$ are prescribed on the entire outer boundary, including the planes $z=z_{\rm min}$ and $z=z_{\rm max}$, for the porous specimens and the homogeneous reference alike, which mimics plane strain while still accounting for the three-dimensional character of the geometry of the microstructure. The amplitude of the K-field is not adapted to the wall width here. Its energy release rate, $J_{\text{bc}} = (1-\nu_{\text{bc}}^2)(K_I^{\text{bc}})^2/E_{\text{bc}}$, is $1.37$, $1.25$, $1.13$, $1.06$ and $0.97\,J_c^0$ for $w_s = 0.375L$, $0.5L$, $0.75L$, $L$ and $2L$, which is close to the effective crack resistances found below (ratio $0.7$ to $1.1$).

Simulations are conducted with wall widths $w_s = 0.375L$, $0.5L$, $0.75L$, $L$ and $2L$ using the parameter set \vMLabel. The time step is controlled as in Section~\ref{sec:numerical_implementation}, except that it is increased when a step converges in fewer than seven Newton iterations, since the plastic ligaments require five to six iterations per step. On the four-pore specimen with $w_s = L$, the thresholds seven and five give values of $J_{x,c}^{\text{eff}}$ that differ by $0.5\,\%$ and records of $J_x$ that agree within $0.1\,\%$ in the median, with $37\,\%$ fewer time steps for the higher threshold.

\subsection{Fracture Behavior and Comparison with the Two-Dimensional Model}
\label{sec:3D_fracture_behavior}

\begin{figure}[htb]
    \centering
    \resizebox{0.45\textwidth}{!}{\input{pics/PAPER_rev2_3D_xct_xbc_vs_t_ws1.0.pgf}}% rev2: regenerated from 076 (S1F_3D_vM_ws1.0), replaces PAPER_VM3D_xct_vs_time.pgf
    \caption{Position of the surfing boundary conditions $x_{\text{bc}}$ and crack tip position $x_{\text{ct}}$ for $w_s=1.0L$ and the parameter set \vMLabel, three-dimensional model.}
    \label{fig:position_3D}
\end{figure}

\begin{figure}[htb]
    \centering
    \resizebox{0.45\textwidth}{!}{\includegraphics{pics/s_t_2.9998_wsteg_0.25.png}}
    \caption{Crack propagation through the central row of pores in the three-dimensional model (phase field $s$, illustration; red: intact material, blue: crack). The half of the specimen with $y > 0$ is indicated by its transparent surface mesh.}
    \label{fig:crack3D_illustration}
\end{figure}

\begin{figure}[htb]
    \centering
    \includegraphics[width=0.72\textwidth]{pics/PAPER_rev2_3D_Jx_vs_xct_2D_vs_3D_ws1.0.pdf}% rev2: regenerated from 075/076, replaces PAPER_VM_Jx_vs_xct_2D_vs_3D.png (+ holes_s_middle_row.png overlay, now part of the figure)
    \caption{J-Integral $J_x$ vs.\ crack tip position $x_{ct}$ for $w_s = L$ and \vMLabel in the two- and the three-dimensional model, each normalized by the numerical fracture resistance $J_c^{\text{num}}$ of the homogeneous material of the respective dimension. Dashed lines mark the pore edges, which are the same in both models (six pores). The picture at the top shows the phase field $s$ of the middle pore row of the two-dimensional specimen after complete separation, Fig.~\ref{fig:PAPER_02_all_Jx_vs_xct_pf_diamond_holes}, for orientation.}
    \label{fig:PAPER_VM_Jx_all_3D}
\end{figure}

\begin{figure}[htb]
    \centering
    \includegraphics[width=0.72\textwidth]{pics/PAPER_rev2_3D_Jx_vs_xct_sections.pdf}% rev2: regenerated from 076, replaces PAPER_VM3D_Jx_vs_xct_all_wsteg.png
    \caption{J-Integral $J_x$ vs.\ crack tip position $x_{ct}$ between pores two and four of the three-dimensional specimen, the same section as in Fig.~\ref{fig:PAPER_rev2_Jx_sections_three_sets}, for the parameter set \vMLabel and varying $w_s$, normalized by $J_c^{\text{num}}$ of the three-dimensional homogeneous material.}
    \label{fig:test12q2_3D}
\end{figure}

\begin{figure}[htb]
    \centering
    \resizebox{0.85\textwidth}{!}{\input{pics/PAPER_rev2_3D_Jeff_and_partition_2D_vs_3D.pgf}}% rev2: regenerated from 075/076, replaces PAPER_VM_Jeff_vs_wsteg_2D_vs_3D.png
    \caption{a) Effective crack resistance $J_{x,c}^{\text{eff}}$ vs.\ wall width $w_s$ for the parameter set \vMLabel in the two- and the three-dimensional model, normalized by the numerical fracture resistance $J_c^{\text{num}}$ of the respective dimension. Filled triangles: $6 \times 3 \times 2$ specimen, open triangles: shorter specimen with four pores in the $x$-direction, open square: the latter on the finer mesh, $h = 0.07L$. b) Fractions of the energy dissipated per pore-to-pore cycle, eq.~(\ref{eq:D_tot_cycle}), that create crack surface (fracture) and that are dissipated by the viscous evolution of $s$ (viscous) in the two- and the three-dimensional ($6 \times 3 \times 2$) model.}
    \label{fig:PAPER_06b_Jx_vs_wsteg_2D_3D}
\end{figure}

This section describes the crack patterns of the three-dimensional series, compares the effective crack resistance with the two-dimensional results and closes with the energy partition. In all simulations, the crack propagates through the central row of pores, Fig.~\ref{fig:crack3D_illustration}. Figure~\ref{fig:position_3D} shows the crack tip position $x_{ct}$, identified as the node with the largest $x$-coordinate that fulfills $s=0$, together with the position $x_{\text{bc}}$ of the K-field for $w_s = L$. In contrast to the two-dimensional case, Fig.~\ref{fig:PAPER_00_xct_pf_vs_xct_KI_diamond}, the crack tip follows the K-field nearly continuously, with small steps at the pores instead of the jumps of about $1.2L$ across the pores. The crack front advances by at most $0.11$ to $0.15L$ per time step, and $x_{ct}$ lies within the $x$-range of a pore in $21$ to $75\,\%$ of the time steps ($w_s = 2L$ to $0.375L$, counted up to the end of the evaluation cycle), which is only possible if the front has advanced beside the sphere. The front thus advances beside the spherical pores instead of being trapped as a whole, and the ligament between two pores is separated progressively rather than at once.

The J-integral is computed according to eq.~(\ref{eq:J-integral}) and normalized by the domain width in the $z$-direction for comparison with the two-dimensional plane strain results. This normalization does not change the definition of the J-integral. It converts the three-dimensional boundary integral to a value per unit out-of-plane thickness.
\begin{equation}
    J_x = \dfrac{1}{2.5 w_c}\int_{\partial\Omega}\left(\psi_{\text{el}}(\boldsymbol{\varepsilon})\boldsymbol{1} - (\nabla\boldsymbol{u})^T\boldsymbol{\sigma}\right)\boldsymbol{n}\,\mathrm{d}A\cdot\boldsymbol{e}_x.
    \label{eq:J_3D}
\end{equation}
Since only the parameter set \vMLabel is considered in this section, and the intention of the following discussion is to highlight the effect of the geometric representation of the microstructure and not of the material law, the records of $J_x$ are normalized by the numerical fracture resistance $J_c^{\text{num}}$ of the homogeneous material in the respective dimension. For the three-dimensional model, $J_c^{\text{num}}$ is determined from a specimen without pores, with the yield stress assigned to the whole domain, constructed and meshed by the same rules as the porous specimens, as the mean of $J_x$ over the central part of the crack path. It is $1.421J_c^0$ for the domain of $w_s = L$ and $1.445J_c^0$ for the larger domain of $w_s = 2L$, which agrees with the two-dimensional value $1.425J_c^0$ of Section~\ref{sec:steady_growth_griffith} within $2\,\%$. The simulations for $w_s \le L$ are normalized by the former value and the simulation for $w_s = 2L$ by the latter. The four-pore specimen is normalized by its own references, $1.415J_c^0$ ($w_s \le L$) and $1.444J_c^0$ ($w_s = 2L$), which differ from the above by less than $0.5\,\%$, and by $1.351J_c^0$ for the finer mesh, $h = 0.07L$.
\prelim{The values normalized by $J_c^{\text{num}}$ in this section are provisional: the homogeneous reference is plastic through the whole thickness, whereas in the porous specimens about $20\,\%$ of the thickness fractures within the elastic side layers of Section~\ref{sec:3D_effective_medium}. A reference with identical side layers is being computed; the normalized values are expected to shift upward by a few percent.}

Figure~\ref{fig:PAPER_VM_Jx_all_3D} compares $J_x$ as a function of $x_{ct}$ for the two- and the three-dimensional model at $w_s = L$. The two-dimensional record shows the sharp re-nucleation peaks of Section~\ref{sec:2D_fracture_behavior}: $J_x$ rises above $J_c^{\text{num}}$ while the crack tip is trapped at the pore wall and drops during the unstable failure of the ligament. The three-dimensional record is a smooth undulation with one maximum per cell. $J_x$ is lowest when the crack tip enters a pore and highest shortly after it has left the pore and entered the ligament, with maxima of \prelim{$0.98$, $0.96$, $0.95$, $0.94$ and $0.90J_c^{\text{num}}$} over the five pore-to-pore cycles, i.e., \prelim{below the crack resistance of the homogeneous material}. Part of the decrease along the specimen is shared by the homogeneous reference, whose $J_x$ drops by about $5\,\%$ towards the end of the domain, relative to which the inner cycles agree within about $1.5\,\%$. Since the front is never trapped as a whole, the forced re-nucleation at a blunt notch, which produces the peaks of the two-dimensional model, does not occur, and since the ligament does not fail at once, there is no phase of unstable crack advance.

The pattern between pores two and four, which contains the ligament in which $J_{x,c}^{\text{eff}}$ is evaluated, is shown for the wall widths in Fig.~\ref{fig:test12q2_3D}. $J_x$ varies smoothly for all wall widths, and the undulation becomes shallower as the pores move apart. For thin walls, the maximum of $J_x$ is attained while $x_{ct}$ is still within the $x$-range of the pore, since the front passes beside the sphere. The effective crack resistance, i.e., the maximum of $J_x$ over the pore-to-pore cycle of the middle ligament between the third and the fourth pore, as in two dimensions, Section~\ref{sec:2D_fracture_behavior}, increases only weakly with the wall width, Fig.~\ref{fig:PAPER_06b_Jx_vs_wsteg_2D_3D}a): $J_{x,c}^{\text{eff}} = 1.22$, $1.25$, $1.31$, $1.34$ and $1.40J_c^0$, i.e., \prelim{$0.86$, $0.88$, $0.92$, $0.95$ and $0.97J_c^{\text{num}}$} for $w_s = 0.375L$, $0.5L$, $0.75L$, $L$ and $2L$. On the four-pore specimen of Section~\ref{sec:3D_effective_medium}, refining the mesh to $h = 0.07L$ decreases the absolute values for $w_s = L$ by $3.6\,\%$ ($J_{x,c}^{\text{eff}}$) and $4.5\,\%$ ($J_c^{\text{num}}$), whereas the normalized value changes by only $+0.9\,\%$, open square in Fig.~\ref{fig:PAPER_06b_Jx_vs_wsteg_2D_3D}a). The normalized effective crack resistance is therefore far less sensitive to the mesh than the absolute values and is used in the comparison with the two-dimensional results. The four-pore specimen itself, with the maximum of $J_x$ taken over the cycle of its middle ligament between the second and the third pore, gives \prelim{$0.87$, $0.89$, $0.93$, $0.95$ and $0.96J_c^{\text{num}}$}, open triangles in Fig.~\ref{fig:PAPER_06b_Jx_vs_wsteg_2D_3D}a), i.e., $0.4$ to $1.5\,\%$ above the six-pore specimen for $w_s \le L$ and $0.7\,\%$ below it for $w_s = 2L$, so that the number of pores in the crack-growth direction has little influence on the normalized effective crack resistance.

The comparison with the two-dimensional model in Fig.~\ref{fig:PAPER_06b_Jx_vs_wsteg_2D_3D}a) shows two differences. In two dimensions, the pores raise the crack resistance of the von~Mises material above the homogeneous value for $w_s \gtrsim 0.9L$, up to $1.28J_c^{\text{num}}$ at $w_s = 2L$ and $1.34J_c^{\text{num}}$ at $w_s = 4L$. In three dimensions, \prelim{the effective crack resistance does not exceed the homogeneous value for any of the computed wall widths, $0.86$ to $0.97J_c^{\text{num}}$ for $w_s = 0.375L$ to $2L$}, and the increase with the wall width is weak. For thin walls, on the other hand, the three-dimensional specimen loses much less of its crack resistance than the two-dimensional one, \prelim{$0.86$} against $0.50J_c^{\text{num}}$ at $w_s = 0.375L$. The two curves cross \prelim{between $w_s = 0.75L$ ($0.92$ against $0.91J_c^{\text{num}}$) and $w_s = L$ ($0.95$ against $1.10J_c^{\text{num}}$)}, the three-dimensional values lie above the two-dimensional ones for thinner and below them for wider walls. Both differences follow from the crack-front evolution described above. The toughening of the two-dimensional specimen at large wall widths stems from the forced re-nucleation at the blunt pore wall, Section~\ref{sec:2D_mechanisms}, which requires the whole front to be trapped. In three dimensions, the front advances beside the spherical pores, so that this mechanism is absent. Equally, a thin ligament does not fail as a whole, which limits the loss of crack resistance in the nucleation regime of Section~\ref{sec:2D_mechanisms}.

The energy partition per pore-to-pore cycle, evaluated as in Section~\ref{sec:2D_energy_partition} for the cycle that contains the evaluation ligament, confirms the different mechanism, Fig.~\ref{fig:PAPER_06b_Jx_vs_wsteg_2D_3D}b). In three dimensions, $91$ to $94\,\%$ of the dissipated energy create crack surface, $6$ to $9\,\%$ remain stored as plastic energy and at most $1\,\%$ is dissipated by the viscous evolution of the phase field, for $w_s = 0.375L$ to $2L$, compared with a viscous share of $12$ to $37\,\%$ in two dimensions, which collects the energy released in the unstable snap-through of the ligaments. In the three-dimensional simulations, the external work and the sum of the stored and dissipated energies, eq.~(\ref{eq:energy_balance}), agree to within $0.4\,\%$ ($0.1\,\%$ for $w_s = L$). The energy dissipated per unit crack advance, evaluated for the same cycle as $J_{x,c}^{\text{eff}}$, is $G_{\text{diss}} = 1.17$ to $1.44J_c^0$, i.e., $G_{\text{diss}}/J_{x,c}^{\text{eff}} = 0.95$ to $1.03$ for $w_s = 0.375L$ to $2L$, whereas the ratio is $0.73$ to $0.93$ in two dimensions, where $J_{x,c}^{\text{eff}}$ is a sharp peak and $G_{\text{diss}}$ the average over the cycle. Across the other pore-to-pore cycles, the ratio varies between $0.82$ and $1.05$, since the stored elastic energy decreases during each cycle and $G_{\text{diss}}$ therefore also contains released elastic energy. In three dimensions, the effective crack resistance is thus close to the energy dissipated per unit crack advance for the evaluation cycle, not pointwise along the crack path.


% The 2D plane strain model from Section~\ref{sec:2Dresults} gives a good indication of the general 
% effect of introducing voids into material on its effective fracture toughness. 
% However, in real porous metals, the pore shape more closely resembles spheres than 
% it does cylindrical which extend in the out of plane direction as it is 
% modelled by the 2D model in Section~\ref{sec:2Dresults}. 

% In order to investigate the threedimensional character that is a closer match
% to real porous metals, we consider spherical pores of diameter $L$, that are centered 
% in cubic cells of edge length $w_c=w_s+L$. Again, we compute the effective stiffness
% of these cubic cells as explained in the previous section, with the 
% exceptions that we now have six load cases $\underline{\boldsymbol{\varepsilon}}_0^k$ 
% and that $\left< \underline{\boldsymbol{C}} \right>$ is now a 6$\times$6 matrix.
% The effective Lam\'{e} parameters are determined by 
% \begin{equation}
%     \lambda^{\text{eff}} = \dfrac{\left< \underline{C} \right>_{12} + 
%     \left< \underline{C} \right>_{13} 
%     + \left< \underline{C} \right>_{21} 
%     + \left< \underline{C} \right>_{31} 
%     + \left< \underline{C} \right>_{23} 
%     + \left< \underline{C} \right>_{32} }{6}, \quad
%     \mu^{\text{eff}} = \dfrac{\left< \underline{C} 
%     \right>_{44} + \left< \underline{C} \right>_{55} + 
%     \left< \underline{C} \right>_{66}}{3}.
% \end{equation}




% A $6\times{3}\times{2}$ array of these unit 
% cells is assembled in order to represent
% a simplified porous material. This array is centered about the $y=0$ plane 
% and again embedded in a 
% a cuboidal domain of effective stiffness with dimensions
% %\footnote{Recall that the width of the domain has a negligible effect on the nucleation length $\Delta a$, but the height does. We want to avoid an influence of the height in this study.}
% $(8 w_c) \times 20 L \times (4 w_c) $. The mesh of bilinear tetrahedral finite elements is refined around the expected 
% fracture plane $y=0$, see Fig.~\ref{fig:mesh3D}.

% An initial crack in the plane $y=0$ for $x-x_{\rm{min}}\le{w}_c$ where $x_{\rm{min}}$ is the 
% minimal $x$-coordinate of any node in the mesh is defined. The finite element edge length
% along the expected crack path is set to $\Delta{h}=0.02L$. 


% The parameters of the phase field model are 
% $\epsilon = 0.1 L$, $\eta = 10^{-5}$, $M = 1000 \nicefrac{v_{\text{bc}}}{J_c^0}$, and $\lambda_0 = \mu_0$. 
% The applied boundary conditions are given by 
% the surfing boundary conditions from section~\ref{sec:numericalExperiment} 
% with $K_I^{\text{bc}} = \sqrt{2.5 J_c^0 \mu_0}$, $E_{\text{bc}} = 2.5 \mu_0$, 
% and $\nu_{\text{bc}} = 0.25$. In addition it is $u_z=0$ at the planes
% $z=z_{\rm{min}}$ and $z=z_{\rm{max}}$ and where $z_{\rm{min}}$ and 
% $z_{\rm{max}}$ are the minimum and maximum $z$-coordinates found in the 
% mesh respectively.

% We conduct simulations for different 
% wall widths $w_s \in \{ 0.25 L, 0.375 L, 0.5 L, 0.75 L, 1.0 L, 2.0 L, 3.0 L, 4.0 L \}$ 
% using only the \vMLabel parameter set from section~\ref{sec:parameter_sets}.

% Figure~\ref{fig:one_row_two_cols}~a) shows the evolution for the phase field $s$ for $w_s=0.25L$ at 
% $t\approx{2.998}\nicefrac{L}{\dot{x}_{bc}}$ with the subdomain $y\ge{0}$ only represented
% as a transparent mesh. The crack again grows through the center 'row' of pores, but 
% in contrast to the 2D model can now continuously grow through the space between pores
% without crack fronts being fully trapped inside pores. Thus, the current crack tip position
% $x_{ct}$ which is again defined as the node with the largest $x$-coordinate and $s=0$, is 
% expected to be a steady function in contrast to the 2D case. 

% Figure~\ref{fig:one_row_two_cols}~b) shows the plastic variable $\alpha$ also for $w_s=0.25L$ at
%  $t\approx{2.998}\nicefrac{L}{\dot{x}_{bc}}$. Plastic deformation mainly occurs 
%  along the crack but also just in front of the advancing crack front. 

%  Figure~\ref{fig:position_3D} illustrates the evolution of the microscale crack tip position
%  $x_{ct}$ as well as the position of the fictional macroscopic crack tip $x_{\rm{bc}}$ for
%  $w_s=1.0L$ and \vMLabel. As expected, the crack evolution is much smoother than what we
%  observe for the 2D case in Fig.~\ref{fig:PAPER_00_xct_pf_vs_xct_KI_diamond} as the crack
%  front can advance continuously in between pores and generally follows the 
%  applied boundary conditions. However, the pores still force 
%  periods of unsteady, jump-like crack propagation. 

% The J-integral is computed again from (\ref{eq:J-integral}). However, in order to be able 
% to compare the obtained values to the 2D results from the previous section, which are values per unit thickness, these values are 
% normalized by the width of the domain in $z$-direction, i.e.
% \begin{equation}
%     J_x = \dfrac{1}{4 w_c}\int_{\partial\Omega}\left(\psi_{\text{el}}(\boldsymbol{\varepsilon})\boldsymbol{1} - \left(\nabla\boldsymbol{u}\right)^T\boldsymbol{\sigma}\right)\boldsymbol{n} \, \mathrm{d}A \cdot \boldsymbol{e}_x. 
% \end{equation}

% In Fig.~\ref{fig:fig:PAPER_VM_Jx_all_3D} the record of $J_x$ vs. $x_ct$ 
% is displayed for $w_s=1.0L$ and \vMLabel for the 2D simulation from the previous section
% and for the 3D model presented here. Also for the 3D simulation, a characteristic periodic pattern emerges with peak values 
% in phases where the crack has to pass the pores and unsteady crack propagation is 
% forced, see also  Fig.~\ref{fig:position_3D}. These peaks define the maxima of 
% $J_x$ and thus the effective fracture toughness. As expected the 3D record is steady in contrast 
% to the 2D record, since the crack front can propagate in between pores uninterrupted. 
% Noteworthy, the 3D simulation also quantitatively aligns well with the 2D simulation as in both 
% case the effective fracture tougnness $J_{x,c}^{\rm{eff}}$ reaches similar values. 

% Again, we choose the pattern of $J_x$ as recorded between pores two and four as 
% representative of the fracture behavior of the given microstructure. This pattern
% is plotted for the 3D model for varying $w_s$ in Fig.~\ref{fig:test12q2_3D}. As in 
% the 2D case $J_{x,c}^{\rm{eff}}$ increases as the wall width $w_s$ gets larger.

% This can also be verified in Fig.~\ref{fig:PAPER_06b_Jx_vs_wsteg_2D_3D} which shows 
% only the effective crack resistance $J_{x,c}^{\rm{eff}}$ versus
% the wall width for the 2D and 3D models and \vMLabel. 

% For the 3D model, very similar observations can be made. The effective fracture
% toughness increases for larger $w_s$ up to a certain saturation value 
% with enhanced fracture toughness compared to a homogeneous material. For the 3D case and 
% \vMLabel this value is about $\approx XXX J_c^{\text{num}}$. This value is close to 
% but slightly smaller than the value of the 2D model for the same material law. 
% In general, the dependence of $J_{x,c}^{\rm{eff}}$ on $w_s$ is less pronounced than 
% in the 2D case, as it is expected since the 3D model mitigates the effect of the pores as 
% it does not force a complete renucleation of trapped cracks. 


% \begin{figure}[t]
% \centering

% \begin{tabular}{c c}
% % Left image with (a) label
% \llap{{a)}\hspace{0mm}}%
% \includegraphics[width=0.49\textwidth]{pics/s_t_2.9998_wsteg_0.25.png} &
% % Right image with (b) label
% \llap{{b)}\hspace{0mm}}%
% \includegraphics[width=0.49\textwidth]{pics/alpha_t_2.9998_wsteg_0.25.png} \\
% \multicolumn{2}{c}{\small $t = 2.998 \nicefrac{L}{\dot{x}_{\rm{bc}}}$}
% \end{tabular}

% \caption{a) evolution of the phase field $s$ and b) evolution of the internal variable $\alpha$ for $w_s=0.25L$}
% \label{fig:one_row_two_cols}
%         \begin{picture}(0,0)
%                     \put(-38,182){\small$s$} % Adjust width as needed
%                     \put(152,182){\small$\alpha$}
%         \end{picture}

% \end{figure}


% \begin{figure}[htb]
%     \centering
%         \resizebox{0.4\textwidth}{!}{\input{pics/PAPER_VM3D_xct_vs_time.pgf}}
%         \caption{%
%             Position of the surfing boundary conditions $x_{\text{bc}}$ and crack tip position $x_{\text{ct}}$ for $w_s=1.0L$ and the
%             parameter set \vMLabel. 
%         }
%         \label{fig:position_3D}
% \end{figure}


% \begin{figure}[htb]
%     \centering
%     \includegraphics[width=0.8\textwidth]{pics/PAPER_VM_Jx_vs_xct_2D_vs_3D.png}
%     \caption{%
%         J-Integral vs. crack tip position $x_{ct}$ for $w_s=1.0L$.
%     }
%     \begin{picture}(0,0)
%         \put(-85,60){%
%             \includegraphics[width=0.57\textwidth]{pics/holes_s_middle_row.png}%
%         }
%     \end{picture}
%     \label{fig:PAPER_VM_Jx_all_3D}
% \end{figure}


% \begin{figure}[htb]
%     \centering
%     \begin{minipage}{0.95\textwidth}
%         \centering
%                 % \resizebox{\textwidth}{!}{\input{pics/PAPER_04_Jx_vs_xct_all_vonMises.pgf}}
%                 %     % \resizebox{0.3\textwidth}{!}{\includegraphics[]{pics/half_notch_specimen.png}}
%          \resizebox{\textwidth}{!}{\includegraphics[width=\textwidth]{pics/PAPER_VM3D_Jx_vs_xct_all_wsteg.png}}
%         \caption{Section of the J-Integral $J_x$ vs. crack tip position $x_{ct}$ between pore two and four for the threedimensional case and for the parameter set \vMLabel and varying $w_s$.}
%     \label{fig:test12q2_3D}
%     \end{minipage}
% %         \begin{minipage}{0.49\textwidth}
% %         \centering
% %         % \resizebox{\textwidth}{!}{\input{pics/PAPER_03_Jx_vs_xct_all_linear_elastic_gc05679.pgf}}
% %         {\includegraphics[width=\textwidth]{pics/PAPER_03B_Jx_vs_xct_all_linear_elastic_Pi.png}}
% %         \caption{Section of the J-Integral $J_x$ vs. crack tip position $x_{ct}$ between pore two and four for parameter set \textbf{Eq}$\boldsymbol{\Pi^*}$ and varying $w_s$. }
% % \label{fig:PAPER_03B_Jx_vs_xct_all_linear_elastic_Pi}
% %     \end{minipage}
%     \label{fig:1212312}
% \end{figure}


% \begin{figure}[htb]
%         \centering
%         \resizebox{0.6\textwidth}{!}{\includegraphics[]{pics/PAPER_VM_Jeff_vs_wsteg_2D_vs_3D.png}}
%         % \resizebox{\textwidth}{!}{\input{pics/PAPER_04_Jx_vs_xct_all_ramberg_osgoods.pgf}}
%         \caption{Effective crack resistance $J_{x,c}^{\text{eff}}$ vs. wall width $w_s$.}
%         \label{fig:PAPER_06b_Jx_vs_wsteg_2D_3D}
%     \label{fig:side_by_side_plots}
% \end{figure}


\section{Conclusion and Outlook}

In this work, we investigated the effective crack resistance of porous materials under the
influence of ductile material behavior using a phase-field modeling framework with a von~Mises-type constitutive law and
compared the results with multiple linear elastic reference models. The methodology followed
the numerical experiment proposed by Hossain et al.~\cite{hossain_effective_2014}, applying
surfing boundary conditions to enforce macroscopic crack propagation at a prescribed speed in a resolved porous region.

In the two-dimensional model with cylindrical pores, the effective crack resistance increases with the wall width between the pores and saturates for all three parameter sets. The von~Mises set lies below the linear elastic set with the same microscopic fracture resistance $J_c^0$ for all wall widths and above the linear elastic set with the same one-dimensional strength only from a wall width of about $0.6L$ on. Normalized by the crack resistance of the homogeneous material of the same set, the pores raise the crack resistance of the linear elastic materials by up to a factor of $2.03$, that of the plastic material only by $1.34$, and for wall widths below about $0.9L$ they lower it. The gain of the plastic material over the linear elastic material of equal strength thus stems from the material model, which raises the crack resistance of the homogeneous material by a factor of $2.36$ at equal strength ($31\,\%$ at equal $J_c^0$), and the microstructure preserves only part of it. The two linear elastic sets are exactly self-similar, their effective crack resistances scale with $J_c$ to within $0.1\,\%$.

The energy balance of the model and the stress state in the ligament ahead of the trapped crack tip identify the mechanisms. In the linear elastic materials, $47$ to $83\,\%$ of the energy dissipated per pore-to-pore cycle are lost in the unstable snap-through of the ligament, in the plastic material only $12$ to $37\,\%$, whereas $49$ to $74\,\%$ create crack surface and $12$ to $16\,\%$ remain stored as plastic energy. The linear elastic ligaments re-nucleate at the one-dimensional fracture stress for all wall widths. Thin ligaments of the plastic material ($w_s \le 0.6L$ in the recorded runs) yield across their width and reach this stress only after a large plastic strain, which keeps the driving force at re-nucleation low and makes the result sensitive to the yield stress and the hardening modulus; in terms of the effective crack resistance, this nucleation-controlled regime extends to $w_s = 0.5L$, where the plastic set lies below the linear elastic set of equal strength, and the two curves cross between $0.5L$ and $0.6L$. In wide ligaments, $w_s \ge L$, the constraint ahead of the trapped crack raises the triaxiality to $1.0$ to $1.1$ and the normal stress at re-nucleation to $1.31$ to $1.34$ times the one-dimensional fracture stress, and the sensitivity to the plastic parameters is weak. The isotropic approximation of the embedding changes the effective crack resistance by at most $0.2\,\%$, the length of the initial crack and the number of resolved pore rows by at most $1.1\,\%$.

The three-dimensional model with spherical pores behaves differently. The crack front advances beside the pores in small increments, no unstable snap-through of a ligament occurs, and the viscous share of the dissipated energy is at most $1\,\%$, so that neither the toughening by forced re-nucleation nor the strong weakening of thin ligaments of the two-dimensional model occurs. The effective crack resistance, normalized by the crack resistance of the three-dimensional homogeneous material, \prelim{increases only from $0.86$ to $0.97$ for $w_s = 0.375L$ to $2L$ and does not exceed the homogeneous value}; the normalized value changes by $+0.9\,\%$ under a mesh refinement from $h = 0.1L$ to $0.07L$ and by at most $1.5\,\%$ with four instead of six pores in the crack-growth direction. It lies above the two-dimensional values for thin walls, $0.86$ against $0.50$ at $w_s = 0.375L$, and below them for $w_s \ge L$. The two-dimensional idealization of cylindrical pores thus does not carry over quantitatively to spherical pores.

The effective crack resistance determined here is an operational quantity of the numerical experiment for the stated microstructure, crack-growth direction and loading, not a homogenized material property. Representativity in the sense of Hill is not established, since statistics over microstructure realizations and other crack-growth directions were not tested, and the size of the specimen in the crack-growth direction only in three dimensions. Such a study is the next step, followed by microstructures derived from micro-CT scans of porous or additively manufactured components and by a comparison with experimentally determined fracture resistances.




% \begin{acknowledgement}
%   Funded by the Deutsche Forschungsgemeinschaft (DFG, German Research Foundation) – Project-ID 511263698 – TRR 375.\\
%   Funded by the Deutsche Forschungsgemeinschaft (DFG, German Research Foundation) – Project-ID 523648868 - FOR 5701.\\
%   Calculations for this research were conducted on the Lichtenberg high performance computer of the TU Darmstadt. \\
% This research was enabled by the FEniCSx library, and we thank its developers and the community for their valuable contributions.\\
% \end{acknowledgement}





% Sample body text. Sample body text. Sample body text. Sample body text. Sample body text. Sample body text. Sample body text. Sample body text.

% \section{This is an example for first level head---section head}\label{sec3}

% \subsection{This is an example for second level head---subsection head}\label{subsec2}

% \subsubsection{This is an example for third level head---subsubsection head}\label{subsubsec2}

% Sample body text. Sample body text. Sample body text. Sample body text. Sample body text. Sample body text. Sample body text. Sample body text. 

% \section{Equations}\label{sec4}

% Equations in \LaTeX\ can either be inline or on-a-line by itself (``display equations''). For
% inline equations use the \verb+$...$+ commands. E.g.: The equation
% $H\psi = E \psi$ is written via the command \verb+$H \psi = E \psi$+.

% For display equations (with auto generated equation numbers)
% one can use the equation or align environments:
% \begin{equation}
% \|\tilde{X}(k)\|^2 \leq\frac{\sum\limits_{i=1}^{p}\left\|\tilde{Y}_i(k)\right\|^2+\sum\limits_{j=1}^{q}\left\|\tilde{Z}_j(k)\right\|^2 }{p+q}.\label{eq1}
% \end{equation}
% where,
% \begin{align}
% D_\mu &=  \partial_\mu - ig \frac{\lambda^a}{2} A^a_\mu \nonumber \\
% F^a_{\mu\nu} &= \partial_\mu A^a_\nu - \partial_\nu A^a_\mu + g f^{abc} A^b_\mu A^a_\nu \label{eq2}
% \end{align}
% Notice the use of \verb+\nonumber+ in the align environment at the end
% of each line, except the last, so as not to produce equation numbers on
% lines where no equation numbers are required. The \verb+\label{}+ command
% should only be used at the last line of an align environment where
% \verb+\nonumber+ is not used.
% \begin{equation}
% Y_\infty = \left( \frac{m}{\textrm{GeV}} \right)^{-3}
%     \left[ 1 + \frac{3 \ln(m/\textrm{GeV})}{15}
%     + \frac{\ln(c_2/5)}{15} \right]
% \end{equation}
% The class file also supports the use of \verb+\mathbb{}+, \verb+\mathscr{}+ and
% \verb+\mathcal{}+ commands. As such \verb+\mathbb{R}+, \verb+\mathscr{R}+
% and \verb+\mathcal{R}+ produces $\mathbb{R}$, $\mathscr{R}$ and $\mathcal{R}$
% respectively (refer Subsubsection~\ref{subsubsec2}).

% \section{Tables}\label{sec5}

% Tables can be inserted via the normal table and tabular environment. To put
% footnotes inside tables you should use \verb+\footnotetext[]{...}+ tag.
% The footnote appears just below the table itself (refer Tables~\ref{tab1} and \ref{tab2}). 
% For the corresponding footnotemark use \verb+\footnotemark[...]+

% \begin{table}[h]
% \caption{Caption text}\label{tab1}%
% \begin{tabular}{@{}llll@{}}
% \toprule
% Column 1 & Column 2  & Column 3 & Column 4\\
% \midrule
% row 1    & data 1   & data 2  & data 3  \\
% row 2    & data 4   & data 5\footnotemark[1]  & data 6  \\
% row 3    & data 7   & data 8  & data 9\footnotemark[2]  \\
% \botrule
% \end{tabular}
% \footnotetext{Source: This is an example of table footnote. This is an example of table footnote.}
% \footnotetext[1]{Example for a first table footnote. This is an example of table footnote.}
% \footnotetext[2]{Example for a second table footnote. This is an example of table footnote.}
% \end{table}

% \noindent
% The input format for the above table is as follows:

% %%=============================================%%
% %% For presentation purpose, we have included  %%
% %% \bigskip command. Please ignore this.       %%
% %%=============================================%%
% \bigskip
% \begin{verbatim}
% \begin{table}[<placement-specifier>]
% \caption{<table-caption>}\label{<table-label>}%
% \begin{tabular}{@{}llll@{}}
% \toprule
% Column 1 & Column 2 & Column 3 & Column 4\\
% \midrule
% row 1 & data 1 & data 2	 & data 3 \\
% row 2 & data 4 & data 5\footnotemark[1] & data 6 \\
% row 3 & data 7 & data 8	 & data 9\footnotemark[2]\\
% \botrule
% \end{tabular}
% \footnotetext{Source: This is an example of table footnote. 
% This is an example of table footnote.}
% \footnotetext[1]{Example for a first table footnote.
% This is an example of table footnote.}
% \footnotetext[2]{Example for a second table footnote. 
% This is an example of table footnote.}
% \end{table}
% \end{verbatim}
% \bigskip
% %%=============================================%%
% %% For presentation purpose, we have included  %%
% %% \bigskip command. Please ignore this.       %%
% %%=============================================%%

% \begin{table}[h]
% \caption{Example of a lengthy table which is set to full textwidth}\label{tab2}
% \begin{tabular*}{\textwidth}{@{\extracolsep\fill}lcccccc}
% \toprule%
% & \multicolumn{3}{@{}c@{}}{Element 1\footnotemark[1]} & \multicolumn{3}{@{}c@{}}{Element 2\footnotemark[2]} \\\cmidrule{2-4}\cmidrule{5-7}%
% Project & Energy & $\sigma_{calc}$ & $\sigma_{expt}$ & Energy & $\sigma_{calc}$ & $\sigma_{expt}$ \\
% \midrule
% Element 3  & 990 A & 1168 & $1547\pm12$ & 780 A & 1166 & $1239\pm100$\\
% Element 4  & 500 A & 961  & $922\pm10$  & 900 A & 1268 & $1092\pm40$\\
% \botrule
% \end{tabular*}
% \footnotetext{Note: This is an example of table footnote. This is an example of table footnote this is an example of table footnote this is an example of~table footnote this is an example of table footnote.}
% \footnotetext[1]{Example for a first table footnote.}
% \footnotetext[2]{Example for a second table footnote.}
% \end{table}

% In case of double column layout, tables which do not fit in single column width should be set to full text width. For this, you need to use \verb+\begin{table*}+ \verb+...+ \verb+\end{table*}+ instead of \verb+\begin{table}+ \verb+...+ \verb+\end{table}+ environment. Lengthy tables which do not fit in textwidth should be set as rotated table. For this, you need to use \verb+\begin{sidewaystable}+ \verb+...+ \verb+\end{sidewaystable}+ instead of \verb+\begin{table*}+ \verb+...+ \verb+\end{table*}+ environment. This environment puts tables rotated to single column width. For tables rotated to double column width, use \verb+\begin{sidewaystable*}+ \verb+...+ \verb+\end{sidewaystable*}+.

% \begin{sidewaystable}
% \caption{Tables which are too long to fit, should be written using the ``sidewaystable'' environment as shown here}\label{tab3}
% \begin{tabular*}{\textheight}{@{\extracolsep\fill}lcccccc}
% \toprule%
% & \multicolumn{3}{@{}c@{}}{Element 1\footnotemark[1]}& \multicolumn{3}{@{}c@{}}{Element\footnotemark[2]} \\\cmidrule{2-4}\cmidrule{5-7}%
% Projectile & Energy	& $\sigma_{calc}$ & $\sigma_{expt}$ & Energy & $\sigma_{calc}$ & $\sigma_{expt}$ \\
% \midrule
% Element 3 & 990 A & 1168 & $1547\pm12$ & 780 A & 1166 & $1239\pm100$ \\
% Element 4 & 500 A & 961  & $922\pm10$  & 900 A & 1268 & $1092\pm40$ \\
% Element 5 & 990 A & 1168 & $1547\pm12$ & 780 A & 1166 & $1239\pm100$ \\
% Element 6 & 500 A & 961  & $922\pm10$  & 900 A & 1268 & $1092\pm40$ \\
% \botrule
% \end{tabular*}
% \footnotetext{Note: This is an example of table footnote this is an example of table footnote this is an example of table footnote this is an example of~table footnote this is an example of table footnote.}
% \footnotetext[1]{This is an example of table footnote.}
% \end{sidewaystable}

% \section{Figures}\label{sec6}

% As per the \LaTeX\ standards you need to use eps images for \LaTeX\ compilation and \verb+pdf/jpg/png+ images for \verb+PDFLaTeX+ compilation. This is one of the major difference between \LaTeX\ and \verb+PDFLaTeX+. Each image should be from a single input .eps/vector image file. Avoid using subfigures. The command for inserting images for \LaTeX\ and \verb+PDFLaTeX+ can be generalized. The package used to insert images in \verb+LaTeX/PDFLaTeX+ is the graphicx package. Figures can be inserted via the normal figure environment as shown in the below example:

% %%=============================================%%
% %% For presentation purpose, we have included  %%
% %% \bigskip command. Please ignore this.       %%
% %%=============================================%%
% \bigskip
% \begin{verbatim}
% \begin{figure}[<placement-specifier>]
% \centering
% \includegraphics{<eps-file>}
% \caption{<figure-caption>}\label{<figure-label>}
% \end{figure}
% \end{verbatim}
% \bigskip
% %%=============================================%%
% %% For presentation purpose, we have included  %%
% %% \bigskip command. Please ignore this.       %%
% %%=============================================%%

% \begin{figure}[h]
% \centering
% \includegraphics[width=0.9\textwidth]{fig.eps}
% \caption{This is a widefig. This is an example of long caption this is an example of long caption  this is an example of long caption this is an example of long caption}\label{fig1}
% \end{figure}

% In case of double column layout, the above format puts figure captions/images to single column width. To get spanned images, we need to provide \verb+\begin{figure*}+ \verb+...+ \verb+\end{figure*}+.

% For sample purpose, we have included the width of images in the optional argument of \verb+\includegraphics+ tag. Please ignore this. 

% \section{Algorithms, Program codes and Listings}\label{sec7}

% Packages \verb+algorithm+, \verb+algorithmicx+ and \verb+algpseudocode+ are used for setting algorithms in \LaTeX\ using the format:

% %%=============================================%%
% %% For presentation purpose, we have included  %%
% %% \bigskip command. Please ignore this.       %%
% %%=============================================%%
% \bigskip
% \begin{verbatim}
% \begin{algorithm}
% \caption{<alg-caption>}\label{<alg-label>}
% \begin{algorithmic}[1]
% . . .
% \end{algorithmic}
% \end{algorithm}
% \end{verbatim}
% \bigskip
% %%=============================================%%
% %% For presentation purpose, we have included  %%
% %% \bigskip command. Please ignore this.       %%
% %%=============================================%%

% You may refer above listed package documentations for more details before setting \verb+algorithm+ environment. For program codes, the ``verbatim'' package is required and the command to be used is \verb+\begin{verbatim}+ \verb+...+ \verb+\end{verbatim}+. 

% Similarly, for \verb+listings+, use the \verb+listings+ package. \verb+\begin{lstlisting}+ \verb+...+ \verb+\end{lstlisting}+ is used to set environments similar to \verb+verbatim+ environment. Refer to the \verb+lstlisting+ package documentation for more details.

% A fast exponentiation procedure:

% \lstset{texcl=true,basicstyle=\small\sf,commentstyle=\small\rm,mathescape=true,escapeinside={(*}{*)}}
% \begin{lstlisting}
% begin
%   for $i:=1$ to $10$ step $1$ do
%       expt($2,i$);  
%       newline() od                (*\textrm{Comments will be set flush to the right margin}*)
% where
% proc expt($x,n$) $\equiv$
%   $z:=1$;
%   do if $n=0$ then exit fi;
%      do if odd($n$) then exit fi;                 
%         comment: (*\textrm{This is a comment statement;}*)
%         $n:=n/2$; $x:=x*x$ od;
%      { $n>0$ };
%      $n:=n-1$; $z:=z*x$ od;
%   print($z$). 
% end
% \end{lstlisting}

% \begin{algorithm}
% \caption{Calculate $y = x^n$}\label{algo1}
% \begin{algorithmic}[1]
% \Require $n \geq 0 \vee x \neq 0$
% \Ensure $y = x^n$ 
% \State $y \Leftarrow 1$
% \If{$n < 0$}\label{algln2}
%         \State $X \Leftarrow 1 / x$
%         \State $N \Leftarrow -n$
% \Else
%         \State $X \Leftarrow x$
%         \State $N \Leftarrow n$
% \EndIf
% \While{$N \neq 0$}
%         \If{$N$ is even}
%             \State $X \Leftarrow X \times X$
%             \State $N \Leftarrow N / 2$
%         \Else[$N$ is odd]
%             \State $y \Leftarrow y \times X$
%             \State $N \Leftarrow N - 1$
%         \EndIf
% \EndWhile
% \end{algorithmic}
% \end{algorithm}

% %%=============================================%%
% %% For presentation purpose, we have included  %%
% %% \bigskip command. Please ignore this.       %%
% %%=============================================%%
% \bigskip
% \begin{minipage}{\hsize}%
% \lstset{frame=single,framexleftmargin=-1pt,framexrightmargin=-17pt,framesep=12pt,linewidth=0.98\textwidth,language=pascal}% Set your language (you can change the language for each code-block optionally)
% %%% Start your code-block
% \begin{lstlisting}
% for i:=maxint to 0 do
% begin
% { do nothing }
% end;
% Write('Case insensitive ');
% Write('Pascal keywords.');
% \end{lstlisting}
% \end{minipage}

% \section{Cross referencing}\label{sec8}

% Environments such as figure, table, equation and align can have a label
% declared via the \verb+\label{#label}+ command. For figures and table
% environments use the \verb+\label{}+ command inside or just
% below the \verb+\caption{}+ command. You can then use the
% \verb+\ref{#label}+ command to cross-reference them. As an example, consider
% the label declared for Figure~\ref{fig1} which is
% \verb+\label{fig1}+. To cross-reference it, use the command 
% \verb+Figure \ref{fig1}+, for which it comes up as
% ``Figure~\ref{fig1}''. 

% To reference line numbers in an algorithm, consider the label declared for the line number 2 of Algorithm~\ref{algo1} is \verb+\label{algln2}+. To cross-reference it, use the command \verb+\ref{algln2}+ for which it comes up as line~\ref{algln2} of Algorithm~\ref{algo1}.

% \subsection{Details on reference citations}\label{subsec7}

% Standard \LaTeX\ permits only numerical citations. To support both numerical and author-year citations this template uses \verb+natbib+ \LaTeX\ package. For style guidance please refer to the template user manual.

% Here is an example for \verb+\cite{...}+: \cite{bib1}. Another example for \verb+\citep{...}+: \citep{bib2}. For author-year citation mode, \verb+\cite{...}+ prints Jones et al. (1990) and \verb+\citep{...}+ prints (Jones et al., 1990).

% All cited bib entries are printed at the end of this article: \cite{bib3}, \cite{bib4}, \cite{bib5}, \cite{bib6}, \cite{bib7}, \cite{bib8}, \cite{bib9}, \cite{bib10}, \cite{bib11}, \cite{bib12} and \cite{bib13}.

% \section{Examples for theorem like environments}\label{sec10}

% For theorem like environments, we require \verb+amsthm+ package. There are three types of predefined theorem styles exists---\verb+thmstyleone+, \verb+thmstyletwo+ and \verb+thmstylethree+ 

% %%=============================================%%
% %% For presentation purpose, we have included  %%
% %% \bigskip command. Please ignore this.       %%
% %%=============================================%%
% \bigskip
% \begin{tabular}{|l|p{19pc}|}
% \hline
% \verb+thmstyleone+ & Numbered, theorem head in bold font and theorem text in italic style \\\hline
% \verb+thmstyletwo+ & Numbered, theorem head in roman font and theorem text in italic style \\\hline
% \verb+thmstylethree+ & Numbered, theorem head in bold font and theorem text in roman style \\\hline
% \end{tabular}
% \bigskip
% %%=============================================%%
% %% For presentation purpose, we have included  %%
% %% \bigskip command. Please ignore this.       %%
% %%=============================================%%

% For mathematics journals, theorem styles can be included as shown in the following examples:

% \begin{theorem}[Theorem subhead]\label{thm1}
% Example theorem text. Example theorem text. Example theorem text. Example theorem text. Example theorem text. 
% Example theorem text. Example theorem text. Example theorem text. Example theorem text. Example theorem text. 
% Example theorem text. 
% \end{theorem}

% Sample body text. Sample body text. Sample body text. Sample body text. Sample body text. Sample body text. Sample body text. Sample body text.

% \begin{proposition}
% Example proposition text. Example proposition text. Example proposition text. Example proposition text. Example proposition text. 
% Example proposition text. Example proposition text. Example proposition text. Example proposition text. Example proposition text. 
% \end{proposition}

% Sample body text. Sample body text. Sample body text. Sample body text. Sample body text. Sample body text. Sample body text. Sample body text.

% \begin{example}
% Phasellus adipiscing semper elit. Proin fermentum massa
% ac quam. Sed diam turpis, molestie vitae, placerat a, molestie nec, leo. Maecenas lacinia. Nam ipsum ligula, eleifend
% at, accumsan nec, suscipit a, ipsum. Morbi blandit ligula feugiat magna. Nunc eleifend consequat lorem. 
% \end{example}

% Sample body text. Sample body text. Sample body text. Sample body text. Sample body text. Sample body text. Sample body text. Sample body text.

% \begin{remark}
% Phasellus adipiscing semper elit. Proin fermentum massa
% ac quam. Sed diam turpis, molestie vitae, placerat a, molestie nec, leo. Maecenas lacinia. Nam ipsum ligula, eleifend
% at, accumsan nec, suscipit a, ipsum. Morbi blandit ligula feugiat magna. Nunc eleifend consequat lorem. 
% \end{remark}

% Sample body text. Sample body text. Sample body text. Sample body text. Sample body text. Sample body text. Sample body text. Sample body text.

% \begin{definition}[Definition sub head]
% Example definition text. Example definition text. Example definition text. Example definition text. Example definition text. Example definition text. Example definition text. Example definition text. 
% \end{definition}

% Additionally a predefined ``proof'' environment is available: \verb+\begin{proof}+ \verb+...+ \verb+\end{proof}+. This prints a ``Proof'' head in italic font style and the ``body text'' in roman font style with an open square at the end of each proof environment. 

% \begin{proof}
% Example for proof text. Example for proof text. Example for proof text. Example for proof text. Example for proof text. Example for proof text. Example for proof text. Example for proof text. Example for proof text. Example for proof text. 
% \end{proof}

% Sample body text. Sample body text. Sample body text. Sample body text. Sample body text. Sample body text. Sample body text. Sample body text.

% \begin{proof}[Proof of Theorem~{\upshape\ref{thm1}}]
% Example for proof text. Example for proof text. Example for proof text. Example for proof text. Example for proof text. Example for proof text. Example for proof text. Example for proof text. Example for proof text. Example for proof text. 
% \end{proof}

% \noindent
% For a quote environment, use \verb+\begin{quote}...\end{quote}+
% \begin{quote}
% Quoted text example. Aliquam porttitor quam a lacus. Praesent vel arcu ut tortor cursus volutpat. In vitae pede quis diam bibendum placerat. Fusce elementum
% convallis neque. Sed dolor orci, scelerisque ac, dapibus nec, ultricies ut, mi. Duis nec dui quis leo sagittis commodo.
% \end{quote}

% Sample body text. Sample body text. Sample body text. Sample body text. Sample body text (refer Figure~\ref{fig1}). Sample body text. Sample body text. Sample body text (refer Table~\ref{tab3}). 

% \section{Methods}\label{sec11}

% Topical subheadings are allowed. Authors must ensure that their Methods section includes adequate experimental and characterization data necessary for others in the field to reproduce their work. Authors are encouraged to include RIIDs where appropriate. 

% \textbf{Ethical approval declarations} (only required where applicable) Any article reporting experiment/s carried out on (i)~live vertebrate (or higher invertebrates), (ii)~humans or (iii)~human samples must include an unambiguous statement within the methods section that meets the following requirements: 

% \begin{enumerate}[1.]
% \item Approval: a statement which confirms that all experimental protocols were approved by a named institutional and/or licensing committee. Please identify the approving body in the methods section

% \item Accordance: a statement explicitly saying that the methods were carried out in accordance with the relevant guidelines and regulations

% \item Informed consent (for experiments involving humans or human tissue samples): include a statement confirming that informed consent was obtained from all participants and/or their legal guardian/s
% \end{enumerate}

% If your manuscript includes potentially identifying patient/participant information, or if it describes human transplantation research, or if it reports results of a clinical trial then  additional information will be required. Please visit (\url{https://www.nature.com/nature-research/editorial-policies}) for Nature Portfolio journals, (\url{https://www.springer.com/gp/authors-editors/journal-author/journal-author-helpdesk/publishing-ethics/14214}) for Springer Nature journals, or (\url{https://www.biomedcentral.com/getpublished/editorial-policies\#ethics+and+consent}) for BMC.

% \section{Discussion}\label{sec12}

% Discussions should be brief and focused. In some disciplines use of Discussion or `Conclusion' is interchangeable. It is not mandatory to use both. Some journals prefer a section `Results and Discussion' followed by a section `Conclusion'. Please refer to Journal-level guidance for any specific requirements. 

% \section{Conclusion}\label{sec13}

% Conclusions may be used to restate your hypothesis or research question, restate your major findings, explain the relevance and the added value of your work, highlight any limitations of your study, describe future directions for research and recommendations. 

% In some disciplines use of Discussion or 'Conclusion' is interchangeable. It is not mandatory to use both. Please refer to Journal-level guidance for any specific requirements. 

% \backmatter

% \bmhead{Supplementary information}

% If your article has accompanying supplementary file/s please state so here. 

% Authors reporting data from electrophoretic gels and blots should supply the full unprocessed scans for key as part of their Supplementary information. This may be requested by the editorial team/s if it is missing.

% Please refer to Journal-level guidance for any specific requirements.

\bmhead{Acknowledgements}

  Funded by the Deutsche Forschungsgemeinschaft (DFG, German Research Foundation) – Project-ID 511263698 – TRR 375.\\
  Funded by the Deutsche Forschungsgemeinschaft (DFG, German Research Foundation) – Project-ID 523648868 - FOR 5701.\\
  Calculations for this research were conducted on the Lichtenberg high performance computer of the TU Darmstadt. \\
This research was enabled by the FEniCSx library~\cite{fenicsx_documentation}, and we thank its developers and the community for their valuable contributions.\\

% Acknowledgements are not compulsory. Where included they should be brief. Grant or contribution numbers may be acknowledged.

% Please refer to Journal-level guidance for any specific requirements.

\section*{Declarations}

\bmhead{Ethical Approval} not applicable.

\bmhead{Funding} Funded by the Deutsche Forschungsgemeinschaft (DFG, German Research Foundation) – Project-ID 511263698 – TRR 375.\\
  Funded by the Deutsche Forschungsgemeinschaft (DFG, German Research Foundation) – Project-ID 523648868 - FOR 5701.\\

  \bmhead{Availability of data and materials} Research data will be published after review. The data will be published on Zenodo and referenced here in the final version of the manuscript.
% Some journals require declarations to be submitted in a standardised format. Please check the Instructions for Authors of the journal to which you are submitting to see if you need to complete this section. If yes, your manuscript must contain the following sections under the heading `Declarations':

% \begin{itemize}
% \item Funding
% \item Conflict of interest/Competing interests (check journal-specific guidelines for which heading to use)
% \item Ethics approval and consent to participate
% \item Consent for publication
% \item Data availability 
% \item Materials availability
% \item Code availability 
% \item Author contribution
% \end{itemize}

% \noindent
% If any of the sections are not relevant to your manuscript, please include the heading and write `Not applicable' for that section. 

% %%===================================================%%
% %% For presentation purpose, we have included        %%
% %% \bigskip command. Please ignore this.             %%
% %%===================================================%%
% \bigskip
% \begin{flushleft}%
% Editorial Policies for:

% \bigskip\noindent
% Springer journals and proceedings: \url{https://www.springer.com/gp/editorial-policies}

% \bigskip\noindent
% Nature Portfolio journals: \url{https://www.nature.com/nature-research/editorial-policies}

% \bigskip\noindent
% \textit{Scientific Reports}: \url{https://www.nature.com/srep/journal-policies/editorial-policies}

% \bigskip\noindent
% BMC journals: \url{https://www.biomedcentral.com/getpublished/editorial-policies}
% \end{flushleft}

% \begin{appendices}

% \section{Section title of first appendix}\label{secA1}

% An appendix contains supplementary information that is not an essential part of the text itself but which may be helpful in providing a more comprehensive understanding of the research problem or it is information that is too cumbersome to be included in the body of the paper.

% %%=============================================%%
% %% For submissions to Nature Portfolio Journals %%
% %% please use the heading ``Extended Data''.   %%
% %%=============================================%%

% %%=============================================================%%
% %% Sample for another appendix section			       %%
% %%=============================================================%%

% %% \section{Example of another appendix section}\label{secA2}%
% %% Appendices may be used for helpful, supporting or essential material that would otherwise 
% %% clutter, break up or be distracting to the text. Appendices can consist of sections, figures, 
% %% tables and equations etc.

% \end{appendices}

%%===========================================================================================%%
%% If you are submitting to one of the Nature Portfolio journals, using the eJP submission   %%
%% system, please include the references within the manuscript file itself. You may do this  %%
%% by copying the reference list from your .bbl file, paste it into the main manuscript .tex %%
%% file, and delete the associated \verb+\bibliography+ commands.                            %%
%%===========================================================================================%%

 % \bibliographystyle{sn-basic}
% \bibliography{bib/literatur.bib}% common bib file
%% if required, the content of .bbl file can be included here once bbl is generated
%%\input sn-article.bbl
% \input{./schueter_medda_mueller_klaim_2025.bbl}

\begin{thebibliography}{44}
% BibTex style file: bmc-mathphys.bst (version 2.1), 2014-07-24
\ifx \bisbn   \undefined \def \bisbn  #1{ISBN #1}\fi
\ifx \binits  \undefined \def \binits#1{#1}\fi
\ifx \bauthor  \undefined \def \bauthor#1{#1}\fi
\ifx \batitle  \undefined \def \batitle#1{#1}\fi
\ifx \bjtitle  \undefined \def \bjtitle#1{#1}\fi
\ifx \bvolume  \undefined \def \bvolume#1{\textbf{#1}}\fi
\ifx \byear  \undefined \def \byear#1{#1}\fi
\ifx \bissue  \undefined \def \bissue#1{#1}\fi
\ifx \bfpage  \undefined \def \bfpage#1{#1}\fi
\ifx \blpage  \undefined \def \blpage #1{#1}\fi
\ifx \burl  \undefined \def \burl#1{\textsf{#1}}\fi
\ifx \doiurl  \undefined \def \doiurl#1{\url{https://doi.org/#1}}\fi
\ifx \betal  \undefined \def \betal{\textit{et al.}}\fi
\ifx \binstitute  \undefined \def \binstitute#1{#1}\fi
\ifx \binstitutionaled  \undefined \def \binstitutionaled#1{#1}\fi
\ifx \bctitle  \undefined \def \bctitle#1{#1}\fi
\ifx \beditor  \undefined \def \beditor#1{#1}\fi
\ifx \bpublisher  \undefined \def \bpublisher#1{#1}\fi
\ifx \bbtitle  \undefined \def \bbtitle#1{#1}\fi
\ifx \bedition  \undefined \def \bedition#1{#1}\fi
\ifx \bseriesno  \undefined \def \bseriesno#1{#1}\fi
\ifx \blocation  \undefined \def \blocation#1{#1}\fi
\ifx \bsertitle  \undefined \def \bsertitle#1{#1}\fi
\ifx \bsnm \undefined \def \bsnm#1{#1}\fi
\ifx \bsuffix \undefined \def \bsuffix#1{#1}\fi
\ifx \bparticle \undefined \def \bparticle#1{#1}\fi
\ifx \barticle \undefined \def \barticle#1{#1}\fi
\bibcommenthead
\ifx \bconfdate \undefined \def \bconfdate #1{#1}\fi
\ifx \botherref \undefined \def \botherref #1{#1}\fi
\ifx \url \undefined \def \url#1{\textsf{#1}}\fi
\ifx \bchapter \undefined \def \bchapter#1{#1}\fi
\ifx \bbook \undefined \def \bbook#1{#1}\fi
\ifx \bcomment \undefined \def \bcomment#1{#1}\fi
\ifx \oauthor \undefined \def \oauthor#1{#1}\fi
\ifx \citeauthoryear \undefined \def \citeauthoryear#1{#1}\fi
\ifx \endbibitem  \undefined \def \endbibitem {}\fi
\ifx \bconflocation  \undefined \def \bconflocation#1{#1}\fi
\ifx \arxivurl  \undefined \def \arxivurl#1{\textsf{#1}}\fi
\csname PreBibitemsHook\endcsname

%%% 1
%% 1
%% 1
%% 1
%% 1
%% 1
%% 1
%% 1
%% 1
%% 1
%% 1
%% 1
\bibitem[\protect\citeauthoryear{Hossain et~al.}{2014}]{hossain_effective_2014}
\begin{barticle}
\bauthor{\bsnm{Hossain}, \binits{M.Z.}},
\bauthor{\bsnm{Hsueh}, \binits{C.-J.}},
\bauthor{\bsnm{Bourdin}, \binits{B.}},
\bauthor{\bsnm{Bhattacharya}, \binits{K.}}:
\batitle{Effective toughness of heterogeneous media}.
\bjtitle{Journal of the Mechanics and Physics of Solids}
\bvolume{71},
\bfpage{15}--\blpage{32}
(\byear{2014})
\doiurl{10.1016/j.jmps.2014.06.002}
\end{barticle}
\endbibitem

%% 2
\bibitem[\protect\citeauthoryear{Hauser et~al.}{2021}]{hauser_porosity_2021}
\begin{barticle}
\bauthor{\bsnm{Hauser}, \binits{T.}},
\bauthor{\bsnm{Reisch}, \binits{R.T.}},
\bauthor{\bsnm{Breese}, \binits{P.P.}},
\bauthor{\bsnm{Lutz}, \binits{B.S.}},
\bauthor{\bsnm{Pantano}, \binits{M.}},
\bauthor{\bsnm{Nalam}, \binits{Y.}},
\bauthor{\bsnm{Bela}, \binits{K.}},
\bauthor{\bsnm{Kamps}, \binits{T.}},
\bauthor{\bsnm{Volpp}, \binits{J.}},
\bauthor{\bsnm{Kaplan}, \binits{A.F.H.}}:
\batitle{Porosity in wire arc additive manufacturing of aluminium alloys}.
\bjtitle{Additive Manufacturing}
\bvolume{41},
\bfpage{101993}
(\byear{2021})
\doiurl{10.1016/j.addma.2021.101993}
\end{barticle}
\endbibitem

%% 3
\bibitem[\protect\citeauthoryear{Gross and Seelig}{2011}]{gross_fracture_2011}
\begin{bbook}
\bauthor{\bsnm{Gross}, \binits{D.}},
\bauthor{\bsnm{Seelig}, \binits{T.}}:
\bbtitle{Fracture {{Mechanics}}: {{With}} an {{Introduction}} to {{Micromechanics}}}.
\bsertitle{Mechanical {{Engineering Series}}}.
\bpublisher{Springer},
\blocation{Berlin, Heidelberg}
(\byear{2011}).
\doiurl{10.1007/978-3-642-19240-1}
\end{bbook}
\endbibitem

%% 4
\bibitem[\protect\citeauthoryear{Hill}{1963}]{hill_elastic_1963}
\begin{barticle}
\bauthor{\bsnm{Hill}, \binits{R.}}:
\batitle{Elastic properties of reinforced solids: Some theoretical principles}.
\bjtitle{Journal of the Mechanics and Physics of Solids}
\bvolume{11}(\bissue{5}),
\bfpage{357}--\blpage{372}
(\byear{1963})
\doiurl{10.1016/0022-5096(63)90036-X}
\end{barticle}
\endbibitem

%% 5
\bibitem[\protect\citeauthoryear{{Nemat-Nasser} and Hori}{2013}]{nemat-nasser_micromechanics_2013}
\begin{bbook}
\bauthor{\bsnm{{Nemat-Nasser}}, \binits{S.}},
\bauthor{\bsnm{Hori}, \binits{M.}}:
\bbtitle{Micromechanics: Overall Properties of Heterogeneous Materials}.
\bpublisher{Elsevier},
\blocation{Amsterdam}
(\byear{2013})
\end{bbook}
\endbibitem

%% 6
\bibitem[\protect\citeauthoryear{Braides et~al.}{1996}]{braides_homogenization_1996}
\begin{barticle}
\bauthor{\bsnm{Braides}, \binits{A.}},
\bauthor{\bsnm{Defranceschi}, \binits{A.}},
\bauthor{\bsnm{Vitali}, \binits{E.}}:
\batitle{Homogenization of free discontinuity problems}.
\bjtitle{Archive for Rational Mechanics and Analysis}
\bvolume{135}(\bissue{4}),
\bfpage{297}--\blpage{356}
(\byear{1996})
\doiurl{10.1007/BF02198476}
\end{barticle}
\endbibitem

%% 7
\bibitem[\protect\citeauthoryear{Francfort and Marigo}{1998}]{francfort_revisiting_1998}
\begin{barticle}
\bauthor{\bsnm{Francfort}, \binits{G.A.}},
\bauthor{\bsnm{Marigo}, \binits{J.-J.}}:
\batitle{Revisiting brittle fracture as an energy minimization problem}.
\bjtitle{Journal of the Mechanics and Physics of Solids}
\bvolume{46}(\bissue{8}),
\bfpage{1319}--\blpage{1342}
(\byear{1998})
\doiurl{10.1016/S0022-5096(98)00034-9}
\end{barticle}
\endbibitem

%% 8
\bibitem[\protect\citeauthoryear{Bourdin et~al.}{2008}]{bourdin_variational_2008}
\begin{barticle}
\bauthor{\bsnm{Bourdin}, \binits{B.}},
\bauthor{\bsnm{Francfort}, \binits{G.A.}},
\bauthor{\bsnm{Marigo}, \binits{J.-J.}}:
\batitle{The {{Variational Approach}} to {{Fracture}}}.
\bjtitle{Journal of Elasticity}
\bvolume{91}(\bissue{1}),
\bfpage{5}--\blpage{148}
(\byear{2008})
\doiurl{10.1007/s10659-007-9107-3}
\end{barticle}
\endbibitem

%% 9
\bibitem[\protect\citeauthoryear{Cagnetti et~al.}{2019}]{cagnetti_stochastic_2019}
\begin{barticle}
\bauthor{\bsnm{Cagnetti}, \binits{F.}},
\bauthor{\bsnm{Dal Maso}, \binits{G.}},
\bauthor{\bsnm{Scardia}, \binits{L.}},
\bauthor{\bsnm{Zeppieri}, \binits{C.I.}}:
\batitle{Stochastic homogenisation of free-discontinuity problems}.
\bjtitle{Archive for Rational Mechanics and Analysis}
\bvolume{233}(\bissue{2}),
\bfpage{935}--\blpage{974}
(\byear{2019})
\doiurl{10.1007/s00205-019-01372-x}
\end{barticle}
\endbibitem

%% 10
\bibitem[\protect\citeauthoryear{Schneider}{2020}]{schneider_fftbased_2020}
\begin{barticle}
\bauthor{\bsnm{Schneider}, \binits{M.}}:
\batitle{An {{FFT-based}} method for computing weighted minimal surfaces in microstructures with applications to the computational homogenization of brittle fracture}.
\bjtitle{International Journal for Numerical Methods in Engineering}
\bvolume{121}(\bissue{7}),
\bfpage{1367}--\blpage{1387}
(\byear{2020})
\doiurl{10.1002/nme.6270}
\end{barticle}
\endbibitem

%% 11
\bibitem[\protect\citeauthoryear{Ernesti and Schneider}{2021}]{ernesti_fast_2021}
\begin{barticle}
\bauthor{\bsnm{Ernesti}, \binits{F.}},
\bauthor{\bsnm{Schneider}, \binits{M.}}:
\batitle{A fast {{Fourier}} transform based method for computing the effective crack energy of a heterogeneous material on a combinatorially consistent grid}.
\bjtitle{International Journal for Numerical Methods in Engineering}
\bvolume{122}(\bissue{21}),
\bfpage{6283}--\blpage{6307}
(\byear{2021})
\doiurl{10.1002/nme.6792}
\end{barticle}
\endbibitem

%% 12
\bibitem[\protect\citeauthoryear{Ernesti and Schneider}{2022}]{ernesti_computing_2022}
\begin{barticle}
\bauthor{\bsnm{Ernesti}, \binits{F.}},
\bauthor{\bsnm{Schneider}, \binits{M.}}:
\batitle{Computing the effective crack energy of heterogeneous and anisotropic microstructures via anisotropic minimal surfaces}.
\bjtitle{Computational Mechanics}
\bvolume{69}(\bissue{1}),
\bfpage{45}--\blpage{57}
(\byear{2022})
\doiurl{10.1007/s00466-021-02082-6}
\end{barticle}
\endbibitem

%% 13
\bibitem[\protect\citeauthoryear{Michel and Suquet}{2022}]{michel_merits_2022}
\begin{barticle}
\bauthor{\bsnm{Michel}, \binits{J.-C.}},
\bauthor{\bsnm{Suquet}, \binits{P.}}:
\batitle{Merits and limits of a variational definition of the effective toughness of heterogeneous materials}.
\bjtitle{Journal of the Mechanics and Physics of Solids}
\bvolume{164},
\bfpage{104889}
(\byear{2022})
\doiurl{10.1016/j.jmps.2022.104889}
\end{barticle}
\endbibitem

%% 14
\bibitem[\protect\citeauthoryear{Lebihain}{2021}]{lebihain_brittle_2021}
\begin{barticle}
\bauthor{\bsnm{Lebihain}, \binits{M.}}:
\batitle{Towards brittle materials with tailored fracture properties: The decisive influence of the material disorder and its microstructure}.
\bjtitle{International Journal of Fracture}
\bvolume{230}(\bissue{1-2}),
\bfpage{99}
(\byear{2021})
\doiurl{10.1007/s10704-021-00538-7}
\end{barticle}
\endbibitem

%% 15
\bibitem[\protect\citeauthoryear{Gao and Rice}{1989}]{gao_firstorder_1989}
\begin{barticle}
\bauthor{\bsnm{Gao}, \binits{H.}},
\bauthor{\bsnm{Rice}, \binits{J.R.}}:
\batitle{A first-order perturbation analysis of crack trapping by arrays of obstacles}.
\bjtitle{Journal of Applied Mechanics}
\bvolume{56}(\bissue{4}),
\bfpage{828}--\blpage{836}
(\byear{1989})
\doiurl{10.1115/1.3176178}
\end{barticle}
\endbibitem

%% 16
\bibitem[\protect\citeauthoryear{Bower and Ortiz}{1991}]{bower_threedimensional_1991}
\begin{barticle}
\bauthor{\bsnm{Bower}, \binits{A.F.}},
\bauthor{\bsnm{Ortiz}, \binits{M.}}:
\batitle{A three-dimensional analysis of crack trapping and bridging by tough particles}.
\bjtitle{Journal of the Mechanics and Physics of Solids}
\bvolume{39}(\bissue{6}),
\bfpage{815}--\blpage{858}
(\byear{1991})
\doiurl{10.1016/0022-5096(91)90026-K}
\end{barticle}
\endbibitem

%% 17
\bibitem[\protect\citeauthoryear{Lebihain et~al.}{2020}]{lebihain_effective_2020}
\begin{barticle}
\bauthor{\bsnm{Lebihain}, \binits{M.}},
\bauthor{\bsnm{Leblond}, \binits{J.-B.}},
\bauthor{\bsnm{Ponson}, \binits{L.}}:
\batitle{Effective toughness of periodic heterogeneous materials: the effect of out-of-plane excursions of cracks}.
\bjtitle{Journal of the Mechanics and Physics of Solids}
\bvolume{137},
\bfpage{103876}
(\byear{2020})
\doiurl{10.1016/j.jmps.2020.103876}
\end{barticle}
\endbibitem

%% 18
\bibitem[\protect\citeauthoryear{Hsueh et~al.}{2018}]{hsueh_stress_2018}
\begin{barticle}
\bauthor{\bsnm{Hsueh}, \binits{C.-J.}},
\bauthor{\bsnm{Avellar}, \binits{L.}},
\bauthor{\bsnm{Bourdin}, \binits{B.}},
\bauthor{\bsnm{Ravichandran}, \binits{G.}},
\bauthor{\bsnm{Bhattacharya}, \binits{K.}}:
\batitle{Stress fluctuation, crack renucleation and toughening in layered materials}.
\bjtitle{Journal of the Mechanics and Physics of Solids}
\bvolume{120},
\bfpage{68}--\blpage{78}
(\byear{2018})
\doiurl{10.1016/j.jmps.2018.04.011}
\end{barticle}
\endbibitem

%% 19
\bibitem[\protect\citeauthoryear{Brach et~al.}{2019}]{brach_anisotropy_2019}
\begin{barticle}
\bauthor{\bsnm{Brach}, \binits{S.}},
\bauthor{\bsnm{Hossain}, \binits{M.Z.}},
\bauthor{\bsnm{Bourdin}, \binits{B.}},
\bauthor{\bsnm{Bhattacharya}, \binits{K.}}:
\batitle{Anisotropy of the effective toughness of layered media}.
\bjtitle{Journal of the Mechanics and Physics of Solids}
\bvolume{131},
\bfpage{96}--\blpage{111}
(\byear{2019})
\doiurl{10.1016/j.jmps.2019.06.021}
\end{barticle}
\endbibitem

%% 20
\bibitem[\protect\citeauthoryear{Liu and Shen}{2025}]{liu_effective_2025}
\begin{barticle}
\bauthor{\bsnm{Liu}, \binits{S.}},
\bauthor{\bsnm{Shen}, \binits{Y.}}:
\batitle{Effective fracture toughness for heterogeneous materials accounting for the worst-case scenario}.
\bjtitle{Journal of the Mechanics and Physics of Solids}
\bvolume{205},
\bfpage{106333}
(\byear{2025})
\doiurl{10.1016/j.jmps.2025.106333}
\end{barticle}
\endbibitem

%% 21
\bibitem[\protect\citeauthoryear{Aranda and Segurado}{2025}]{aranda_cracklength_2025}
\begin{barticle}
\bauthor{\bsnm{Aranda}, \binits{P.}},
\bauthor{\bsnm{Segurado}, \binits{J.}}:
\batitle{A crack-length control technique for phase-field fracture in FFT homogenization}.
\bjtitle{International Journal for Numerical Methods in Engineering}
\bvolume{126},
\bfpage{e7664}
(\byear{2025})
\doiurl{10.1002/nme.7664}
\end{barticle}
\endbibitem

%% 22
\bibitem[\protect\citeauthoryear{Man et~al.}{2025}]{man_collaborative_2025}
\begin{barticle}
\bauthor{\bsnm{Man}, \binits{S.}},
\bauthor{\bsnm{Yu}, \binits{H.}},
\bauthor{\bsnm{Wang}, \binits{J.}}:
\batitle{Collaborative control of crack guiding and trapping in bioinspired interfaces on effective toughness}.
\bjtitle{Composite Structures}
\bvolume{362},
\bfpage{119094}
(\byear{2025})
\doiurl{10.1016/j.compstruct.2025.119094}
\end{barticle}
\endbibitem

%% 23
\bibitem[\protect\citeauthoryear{Henry}{2025}]{henry_phasefield_2025}
\begin{barticle}
\bauthor{\bsnm{Henry}, \binits{H.}}:
\batitle{Phase-field study of the effective fracture energy increase during dynamic crack propagation in disordered heterogeneous materials}.
\bjtitle{Physical Review E}
\bvolume{111},
\bfpage{055502}
(\byear{2025})
\doiurl{10.1103/PhysRevE.111.055502}
\end{barticle}
\endbibitem

%% 24
\bibitem[\protect\citeauthoryear{Hao et~al.}{2026}]{hao_variable_2026}
\begin{barticle}
\bauthor{\bsnm{Hao}, \binits{T.}},
\bauthor{\bsnm{Piel}, \binits{A.}},
\bauthor{\bsnm{Hossain}, \binits{Z.M.}}:
\batitle{Variable stiffness boundary condition to determine effective toughness of heterogeneous materials}.
\bjtitle{Computer Methods in Applied Mechanics and Engineering}
\bvolume{448},
\bfpage{118414}
(\byear{2026})
\doiurl{10.1016/j.cma.2025.118414}
\end{barticle}
\endbibitem

%% 25
\bibitem[\protect\citeauthoryear{Ambati et~al.}{2015}]{ambati_review_2015}
\begin{barticle}
\bauthor{\bsnm{Ambati}, \binits{M.}},
\bauthor{\bsnm{Gerasimov}, \binits{T.}},
\bauthor{\bsnm{De~Lorenzis}, \binits{L.}}:
\batitle{A review on phase-field models of brittle fracture and a new fast hybrid formulation}.
\bjtitle{Computational Mechanics}
\bvolume{55}(\bissue{2}),
\bfpage{383}--\blpage{405}
(\byear{2015})
\doiurl{10.1007/s00466-014-1109-y}
\end{barticle}
\endbibitem

%% 26
\bibitem[\protect\citeauthoryear{Alessi et~al.}{2014}]{alessi_gradient_2014}
\begin{barticle}
\bauthor{\bsnm{Alessi}, \binits{R.}},
\bauthor{\bsnm{Marigo}, \binits{J.-J.}},
\bauthor{\bsnm{Vidoli}, \binits{S.}}:
\batitle{Gradient damage models coupled with plasticity and nucleation of cohesive cracks}.
\bjtitle{Archive for Rational Mechanics and Analysis}
\bvolume{214}(\bissue{2}),
\bfpage{575}--\blpage{615}
(\byear{2014})
\doiurl{10.1007/s00205-014-0763-8}
\end{barticle}
\endbibitem

%% 27
\bibitem[\protect\citeauthoryear{Miehe et~al.}{2016a}]{miehe_phasefield_ductile_2016}
\begin{barticle}
\bauthor{\bsnm{Miehe}, \binits{C.}},
\bauthor{\bsnm{Aldakheel}, \binits{F.}},
\bauthor{\bsnm{Raina}, \binits{A.}}:
\batitle{Phase field modeling of ductile fracture at finite strains: A variational gradient-extended plasticity-damage theory}.
\bjtitle{International Journal of Plasticity}
\bvolume{84},
\bfpage{1}--\blpage{32}
(\byear{2016})
\doiurl{10.1016/j.ijplas.2016.04.011}
\end{barticle}
\endbibitem

%% 28
\bibitem[\protect\citeauthoryear{Miehe et~al.}{2016b}]{miehe_phasefield_porous_2016}
\begin{barticle}
\bauthor{\bsnm{Miehe}, \binits{C.}},
\bauthor{\bsnm{Kienle}, \binits{D.}},
\bauthor{\bsnm{Aldakheel}, \binits{F.}},
\bauthor{\bsnm{Teichtmeister}, \binits{S.}}:
\batitle{Phase field modeling of fracture in porous plasticity: A variational gradient-extended Eulerian framework for the macroscopic analysis of ductile failure}.
\bjtitle{Computer Methods in Applied Mechanics and Engineering}
\bvolume{312},
\bfpage{3}--\blpage{50}
(\byear{2016})
\doiurl{10.1016/j.cma.2016.09.028}
\end{barticle}
\endbibitem

%% 29
\bibitem[\protect\citeauthoryear{Borden et~al.}{2016}]{borden_phasefield_2016}
\begin{barticle}
\bauthor{\bsnm{Borden}, \binits{M.J.}},
\bauthor{\bsnm{Hughes}, \binits{T.J.R.}},
\bauthor{\bsnm{Landis}, \binits{C.M.}},
\bauthor{\bsnm{Anvari}, \binits{A.}},
\bauthor{\bsnm{Lee}, \binits{I.J.}}:
\batitle{A phase-field formulation for fracture in ductile materials: Finite deformation balance law derivation, plastic degradation, and stress triaxiality effects}.
\bjtitle{Computer Methods in Applied Mechanics and Engineering}
\bvolume{312},
\bfpage{130}--\blpage{166}
(\byear{2016})
\doiurl{10.1016/j.cma.2016.09.005}
\end{barticle}
\endbibitem

%% 30
\bibitem[\protect\citeauthoryear{Ambati et~al.}{2015}]{ambati_ductile_2015}
\begin{barticle}
\bauthor{\bsnm{Ambati}, \binits{M.}},
\bauthor{\bsnm{Gerasimov}, \binits{T.}},
\bauthor{\bsnm{De Lorenzis}, \binits{L.}}:
\batitle{Phase-field modeling of ductile fracture}.
\bjtitle{Computational Mechanics}
\bvolume{55}(\bissue{5}),
\bfpage{1017}--\blpage{1040}
(\byear{2015})
\doiurl{10.1007/s00466-015-1151-4}
\end{barticle}
\endbibitem

%% 31
\bibitem[\protect\citeauthoryear{Kuhn et~al.}{2016}]{kuhn_phase_2016}
\begin{barticle}
\bauthor{\bsnm{Kuhn}, \binits{C.}},
\bauthor{\bsnm{Noll}, \binits{T.}},
\bauthor{\bsnm{M{\"u}ller}, \binits{R.}}:
\batitle{On phase field modeling of ductile fracture}.
\bjtitle{GAMM-Mitteilungen}
\bvolume{39}(\bissue{1}),
\bfpage{35}--\blpage{54}
(\byear{2016})
\doiurl{10.1002/gamm.201610003}
\end{barticle}
\endbibitem

%% 32
\bibitem[\protect\citeauthoryear{Noll et~al.}{2020}]{noll_3d_2020}
\begin{barticle}
\bauthor{\bsnm{Noll}, \binits{T.}},
\bauthor{\bsnm{Kuhn}, \binits{C.}},
\bauthor{\bsnm{Olesch}, \binits{D.}},
\bauthor{\bsnm{M{\"u}ller}, \binits{R.}}:
\batitle{{{3D}} phase field simulations of ductile fracture}.
\bjtitle{GAMM-Mitteilungen}
\bvolume{43}(\bissue{2}),
\bfpage{202000008}
(\byear{2020})
\doiurl{10.1002/gamm.202000008}
\end{barticle}
\endbibitem

%% 33
\bibitem[\protect\citeauthoryear{Schl{\"u}ter and M{\"u}ller}{2025a}]{schlueter_ermittlung_2025}
\begin{bchapter}
\bauthor{\bsnm{Schl{\"u}ter}, \binits{A.}},
\bauthor{\bsnm{M{\"u}ller}, \binits{R.}}:
\bctitle{Ermittlung des effektiven Risswiderstands in por{\"o}sen Materialien mit einem Bruch-Phasenfeldmodell}.
In: \bbtitle{57. Tagung des DVM-Arbeitskreises Bruchmechanik und Bauteilsicherheit}.
\bpublisher{Deutscher Verband f{\"u}r Materialforschung und -pr{\"u}fung e.V.},
\blocation{Berlin}
(\byear{2025})
\end{bchapter}
\endbibitem

%% 34
\bibitem[\protect\citeauthoryear{Schl{\"u}ter and M{\"u}ller}{2025b}]{schluter_determination_2025}
\begin{barticle}
\bauthor{\bsnm{Schl{\"u}ter}, \binits{A.}},
\bauthor{\bsnm{M{\"u}ller}, \binits{R.}}:
\batitle{Determination of the effective crack resistance in porous materials using a fracture phase-field model}.
\bjtitle{Engineering Fracture Mechanics}
\bvolume{326},
\bfpage{111348}
(\byear{2025})
\doiurl{10.1016/j.engfracmech.2025.111348}
\end{barticle}
\endbibitem

%% 35
\bibitem[\protect\citeauthoryear{Schl{\"u}ter et~al.}{2025}]{schluter_influence_2025}
\begin{barticle}
\bauthor{\bsnm{Schl{\"u}ter}, \binits{A.}},
\bauthor{\bsnm{Medda}, \binits{R.}},
\bauthor{\bsnm{M{\"u}ller}, \binits{R.}}:
\batitle{Influence of {{Material Nonlinearity}} on the {{Effective Fracture Resistance}} of {{Porous Materials}}}.
\bjtitle{PAMM}
\bvolume{25}(\bissue{4}),
\bfpage{70044}
(\byear{2025})
\doiurl{10.1002/pamm.70044}
\end{barticle}
\endbibitem

%% 36
\bibitem[\protect\citeauthoryear{Bohnen et~al.}{2024}]{bohnen_simulation_2024}
\begin{bchapter}
\bauthor{\bsnm{Bohnen}, \binits{M.}},
\bauthor{\bsnm{M{\"u}ller}, \binits{R.}},
\bauthor{\bsnm{Gross}, \binits{D.}}:
\bctitle{Simulation of~{{Crack Propagation}} in~{{Heterogeneous Materials}} by~a~{{Fracture Phase Field}}}.
In: \beditor{\bsnm{M{\"u}ller}, \binits{W.H.}},
\beditor{\bsnm{Noe}, \binits{A.}},
\beditor{\bsnm{Ferber}, \binits{F.}} (eds.)
\bbtitle{New {{Achievements}} in {{Mechanics}}: {{A Tribute}} to {{Klaus Peter Herrmann}}},
pp. \bfpage{191}--\blpage{215}.
\bpublisher{Springer},
\blocation{Cham}
(\byear{2024}).
\doiurl{10.1007/978-3-031-56132-0_10}
\end{bchapter}
\endbibitem

%% 37
\bibitem[\protect\citeauthoryear{Miehe et~al.}{2017}]{miehe_phasefield_implementation_2017}
\begin{barticle}
\bauthor{\bsnm{Miehe}, \binits{C.}},
\bauthor{\bsnm{Aldakheel}, \binits{F.}},
\bauthor{\bsnm{Teichtmeister}, \binits{S.}}:
\batitle{Phase-field modeling of ductile fracture at finite strains: A robust variational-based numerical implementation of a gradient-extended theory by micromorphic regularization}.
\bjtitle{International Journal for Numerical Methods in Engineering}
\bvolume{111}(\bissue{9}),
\bfpage{816}--\blpage{863}
(\byear{2017})
\doiurl{10.1002/nme.5484}
\end{barticle}
\endbibitem

%% 38
\bibitem[\protect\citeauthoryear{Simo and Hughes}{1998}]{simo_computational_1998}
\begin{bbook}
\bauthor{\bsnm{Simo}, \binits{J.C.}},
\bauthor{\bsnm{Hughes}, \binits{T.J.R.}}:
\bbtitle{Computational Inelasticity}.
\bsertitle{Interdisciplinary Applied Mathematics}.
\bvolume{7}.
\bpublisher{Springer},
\blocation{New York}
(\byear{1998}).
\doiurl{10.1007/b98904}
\end{bbook}
\endbibitem

%% 39
\bibitem[\protect\citeauthoryear{Kuhn}{2013}]{kuhn_numerical_2013}
\begin{botherref}
\oauthor{\bsnm{Kuhn}, \binits{C.}}:
Numerical and {{Analytical Investigation}} of a {{Phase Field Model}} for {{Fracture}}.
PhD thesis,
Technische Universit\"at Kaiserslautern
(March 2013)
\end{botherref}
\endbibitem

%% 40
\bibitem[\protect\citeauthoryear{Chambolle}{2004}]{chambolle_approximation_2004}
\begin{barticle}
\bauthor{\bsnm{Chambolle}, \binits{A.}}:
\batitle{An approximation result for special functions with bounded deformation}.
\bjtitle{Journal de Math\'ematiques Pures et Appliqu\'ees}
\bvolume{83}(\bissue{7}),
\bfpage{929}--\blpage{954}
(\byear{2004})
\doiurl{10.1016/j.matpur.2004.02.004}
\end{barticle}
\endbibitem

%% 41
\bibitem[\protect\citeauthoryear{Kuhn and M{\"u}ller}{2010}]{kuhn_continuum_2010}
\begin{barticle}
\bauthor{\bsnm{Kuhn}, \binits{C.}},
\bauthor{\bsnm{M{\"u}ller}, \binits{R.}}:
\batitle{A continuum phase field model for fracture}.
\bjtitle{Engineering Fracture Mechanics}
\bvolume{77}(\bissue{18}),
\bfpage{3625}--\blpage{3634}
(\byear{2010})
\doiurl{10.1016/j.engfracmech.2010.08.009}
\end{barticle}
\endbibitem

%% 42
\bibitem[\protect\citeauthoryear{{FEniCS Project}}{{FEniCS Project}}{2025}]{fenicsx_documentation}
\begin{botherref}
\oauthor{\bsnm{{FEniCS Project}}}:
FEniCSx Documentation.
\url{https://fenicsproject.org/}.
Accessed 22 July 2025
\end{botherref}
\endbibitem

%% 43
\bibitem[\protect\citeauthoryear{Tann{\'e} et~al.}{2018}]{tanne_crack_2018}
\begin{barticle}
\bauthor{\bsnm{Tann{\'e}}, \binits{E.}},
\bauthor{\bsnm{Li}, \binits{T.}},
\bauthor{\bsnm{Bourdin}, \binits{B.}},
\bauthor{\bsnm{Marigo}, \binits{J.-J.}},
\bauthor{\bsnm{Maurini}, \binits{C.}}:
\batitle{Crack nucleation in variational phase-field models of brittle fracture}.
\bjtitle{Journal of the Mechanics and Physics of Solids}
\bvolume{110},
\bfpage{80}--\blpage{99}
(\byear{2018})
\doiurl{10.1016/j.jmps.2017.09.006}
\end{barticle}
\endbibitem

%% 44
\bibitem[\protect\citeauthoryear{Schlueter}{2018}]{schlueter_phase_2018}
\begin{botherref}
\oauthor{\bsnm{Schlueter}, \binits{A.}}:
Phase {{Field Modeling}} of {{Dynamic Brittle Fracture}}.
PhD thesis,
Technische Universit\"at Kaiserslautern
(June 2018)
\end{botherref}
\endbibitem

\end{thebibliography}

\end{document}
